% !TEX program = xelatex
% 中文译本：Katate Masatsuka (Hiroaki Nishikawa), I do like CFD, VOL.1, Second Edition.
\documentclass[11pt,a4paper,openany,fontset=none]{ctexbook}
\usepackage[margin=2.55cm,headheight=15pt]{geometry}
\usepackage{microtype}
\usepackage{fancyhdr}
\usepackage{hyperref}
\usepackage{xurl}
\usepackage{titlesec}
\usepackage{amsmath,amssymb,amsfonts,bm,mathtools}
\usepackage[export]{adjustbox}
\usepackage{enumitem}
\usepackage{graphicx}
\usepackage{longtable,booktabs,array}
\usepackage{needspace}
\setmainfont[AutoFakeBold=2,AutoFakeSlant=.2]{Arial Unicode MS}
\setCJKmainfont{Songti SC}
\setCJKsansfont{Songti SC}
\setCJKmonofont{Songti SC}
\setlength{\parindent}{2em}
\setlength{\parskip}{0.62em}
\setlength{\emergencystretch}{3em}
\linespread{1.32}\selectfont
\xeCJKsetup{CJKglue=\hskip 0.08em plus 0.04em minus 0.03em}
\setlist[itemize]{leftmargin=2.8em,itemsep=0.32em,topsep=0.48em}
\titleformat{\chapter}[display]{\bfseries\Huge\centering}{第\thechapter 章}{0.75em}{}
\titleformat{\section}{\large\bfseries}{}{0pt}{}
\titleformat{\subsection}{\normalsize\bfseries}{}{0pt}{}
\titlespacing*{\section}{0pt}{1.45em}{0.72em}
\titlespacing*{\subsection}{0pt}{1.20em}{0.58em}
\pagestyle{fancy}
\fancyhf{}
\fancyhead[C]{\small\nouppercase{\leftmark}}
\fancyfoot[C]{\thepage}
\renewcommand{\headrulewidth}{0.3pt}
\hypersetup{hypertexnames=false,colorlinks=true,linkcolor=black,urlcolor=blue,
  pdftitle={我喜欢计算流体力学：第一卷——控制方程与精确解},
  pdfauthor={Katate Masatsuka（Hiroaki Nishikawa）}}
\graphicspath{{./}}
\title{我喜欢计算流体力学：第一卷\\[0.45em]
\Large 控制方程与精确解（第二版）\\[0.7em]\large 中文译文}
\author{Katate Masatsuka（Hiroaki Nishikawa）著\\\small AI 辅助翻译与排版}
\date{原著 PDF 2.6 版（2018）}
\clubpenalty=10000
\widowpenalty=10000
\displaywidowpenalty=10000
\raggedbottom
\allowdisplaybreaks[1]
\setlength{\abovedisplayskip}{0.82\baselineskip plus 0.2\baselineskip minus 0.1\baselineskip}
\setlength{\belowdisplayskip}{0.82\baselineskip plus 0.2\baselineskip minus 0.1\baselineskip}
\setlength{\abovedisplayshortskip}{0.58\baselineskip}
\setlength{\belowdisplayshortskip}{0.58\baselineskip}
\newenvironment{sourceblock}
  {\par\addvspace{0.65\baselineskip}\noindent\begin{minipage}{\linewidth}\centering}
  {\end{minipage}\par\addvspace{0.48\baselineskip}}
% 参数依次为：PDF 页号、左/下/右/上裁边（bp）、最大高度。
% 原 PDF 带有嵌入权限限制；本次成品编译时使用了临时 Ghostscript 重写副本，
% 并在完成后删除。若重新编译，可先生成同名 Nishikawa_embed.pdf；否则尝试原文件。
\IfFileExists{Nishikawa_embed.pdf}{\def\NishikawaSource{Nishikawa_embed.pdf}}
  {\def\NishikawaSource{Nishikawa.pdf}}
\newcommand{\sourcecrop}[6]{%
  \begin{sourceblock}%
  \includegraphics[page=#1,pagebox=mediabox,viewport={#2bp #3bp #4bp #5bp},clip,
    width=0.96\textwidth,max height=#6\textheight]{\NishikawaSource}%
  \end{sourceblock}%
}
\newcommand{\eqnote}[1]{\par\smallskip\noindent\textit{式中：}#1\par}

\begin{document}
\frontmatter
\maketitle
\chapter*{译本说明}
\markboth{译本说明}{译本说明}
本书中文译文按正常教材阅读顺序重新排版，不与原 PDF 逐页对齐，也不保留原书页码目录。
章节另起页，页眉显示当前章名；正文采用较宽松的行距、字距与段间距。

说明性文字由 AI 按计算流体力学、气体动力学和数值分析的通行术语翻译并统一；短公式重写为
LaTeX，复杂矩阵、长推导与图形从原 PDF 以矢量裁片保留，以免文字层抽取损坏上下标、分式和
矩阵结构。裁片与题注作为不可拆分的版块处理，避免跨页、侵入页眉或只剩题注的空白页。
译文供学习与检索使用；涉及公式、模型常数及边界条件时，建议同时核对原著。

\begingroup\small\tableofcontents\endgroup

\chapter*{常用符号与中文释义}
\addcontentsline{toc}{chapter}{常用符号与中文释义}
\markboth{常用符号与中文释义}{常用符号与中文释义}
\begingroup\small
\renewcommand{\arraystretch}{1.20}
\begin{longtable}{>{\centering\arraybackslash}p{0.22\textwidth}p{0.68\textwidth}}
\toprule 符号 & 中文释义 \\ \midrule
\endfirsthead
\toprule 符号 & 中文释义 \\ \midrule
\endhead
$x,y,z$；$t$ & 笛卡尔空间坐标；时间 \\
$u,v,w$；$\bm{v}$ & 三个方向的速度分量；速度矢量 \\
$\rho,p,T$ & 密度、压力、温度 \\
$e,h,E,H$ & 内能、焓、总能、总焓；具体为单位质量量或体积量时依原式定义 \\
$\mu,\nu,k,\kappa$ & 动力黏度、运动黏度、热导率及文中相应扩散/模型系数 \\
$\gamma,R,c_p,c_v$ & 比热比、气体常数、定压比热和定容比热 \\
$a,c$ & 声速（部分章节以 $c$ 表示）；亦可能作为一般常数，依上下文判断 \\
$\bm{U}$ & 守恒变量矢量或一般未知量矢量 \\
$\bm{F},\bm{G},\bm{H}$ & $x,y,z$ 方向的通量矢量 \\
$\tau_{ij}$；$\bm{q}$ & 黏性应力张量分量；热流密度矢量 \\
$\nabla,\operatorname{grad},\operatorname{div},\operatorname{curl}$ & 纳布拉算子、梯度、散度和旋度 \\
$D/Dt$ & 随体导数（物质导数、拉格朗日导数） \\
$\delta_{ij},\epsilon_{ijk}$ & 克罗内克 delta 与 Levi--Civita（Eddington）符号 \\
$\lambda$；$\bm{R},\bm{L}$ & 特征值；右、左特征向量矩阵 \\
$M,Re,Pr$ & 马赫数、雷诺数和普朗特数 \\
$\overline{(\cdot)},(\cdot)'$ & 平均量与脉动量 \\
$k_{\mathrm{t}},\varepsilon,\omega$ & 湍动能、湍动能耗散率和比耗散率（按章节约定） \\
\bottomrule
\end{longtable}
\endgroup

\clearpage
\makeatletter\@mainmattertrue\makeatother
\pagenumbering{arabic}

% ---- 合并分片 1 ----
% Nishikawa.pdf，PDF 第 11--70 页中文译文
% 覆盖：前言、第 1 章，以及第 2 章 2.1--2.15。

\chapter*{前言}
\addcontentsline{toc}{chapter}{前言}
\markboth{前言}{前言}

\section*{第二版（2013）}

第一版出版以后，我仍然喜欢 CFD，而且又发现了更多令我着迷的主题。我还收到了许多同样喜欢 CFD 的读者发来的勘误报告（非常感谢！）。因此，我决定修订本书。书中的选材依然带有很强的个人偏好，不过对于学习 CFD 基础的学生，以及开发基础 CFD 算法的研究人员而言，它仍可作为一本有用的参考书。我希望第二版能让他们更加喜欢 CFD。

\begin{flushright}
Katate Masatsuka\par
2013 年 10 月，约克敦
\end{flushright}

\section*{第一版（2009）}

CFD，即计算流体力学（Computational Fluid Dynamics）。

我喜欢 CFD。正因为如此喜欢，我决定写一本书，说说自己究竟喜欢它什么。本卷集中讨论 CFD 中使用的基本记号与公式、控制方程以及精确解。我希望你也会喜欢这些内容。

\begin{itemize}
  \item \textbf{基础：}一些基本内容，包括记号、公式和定理。
  \item \textbf{模型方程：}算法开发中常用的模型方程。
  \item \textbf{Euler 方程：}无黏流动的 Euler 方程及其相关方程。
  \item \textbf{Navier--Stokes 方程：}黏性流动的 Navier--Stokes 方程及其相关方程。
  \item \textbf{湍流方程：}平均 Navier--Stokes 方程，以及少量湍流模型知识。
  \item \textbf{精确解 I：}一般解以及推导精确解的技巧。
  \item \textbf{精确解 II：}若干可用于精度研究的方程的精确解。
\end{itemize}

自然，本书内容带有很强的个人偏好：这里写的，就是我所喜欢的 CFD；我不喜欢的主题没有收入。因此，这当然不能算一本标准教材，但仍可以是一本不错的参考书。尤其是对正在学习 CFD 基础的学生，以及开发基础 CFD 算法的研究人员，本卷或许会有所帮助（我也希望他们喜欢 CFD）。

\begin{flushright}
Katate Masatsuka\par
2009 年 2 月，约克敦
\end{flushright}

\mainmatter
\chapter{基础}
\markboth{第 1 章\quad 基础}{第 1 章\quad 基础}

\section{微分记号}

下列导数记号都很常用：
\begin{align}
  \frac{\partial u}{\partial x}&=\partial_xu=u_x, &&\text{一阶导数}, \tag{1.1.1}\\
  \frac{\partial^2u}{\partial x^2}&=\partial_{xx}u=u_{xx}, &&\text{二阶导数}, \tag{1.1.2}\\
  \frac{\partial^nu}{\partial x^n}&=\underbrace{\partial_{x\cdots x}}_{n\text{ 次}}u
  =u_{\underbrace{x\cdots x}_{n\text{ 次}}}, &&n\text{ 阶导数}. \tag{1.1.3}
\end{align}

同一件事有多种表达方式很方便。我尤其喜欢 $u_x$ 这样的下标记号，因为它易写，而且比其他记号少占许多竖直空间。当然，导数阶数很高时，下标会变得很长；不过 CFD 书籍里通常不会频繁出现那样的高阶导数。

物质导数、随体导数或 Lagrange 导数写作
\begin{equation}
  \frac{D\alpha}{Dt}=\frac{\partial\alpha}{\partial t}
  +u\frac{\partial\alpha}{\partial x}
  +v\frac{\partial\alpha}{\partial y}
  +w\frac{\partial\alpha}{\partial z}. \tag{1.1.4}
\end{equation}
我很喜欢这个导数。它表示量 $\alpha$ 随速度 $(u,v,w)$ 被流动携带时的时间变化率。因此，若 $D\alpha/Dt=0$，空间分布 $\alpha(x,y,z)$（例如其等值线图）便以速度 $(u,v,w)$ 平移，同时保持初始形状。反之，若普通时间偏导 $\partial_t\alpha=0$，则该分布不随时间移动，而始终固定在空间中。

\section{矢量与运算}

\subsection{矢量与张量}

矢量和张量通常用黑体字母表示，而不是在字母上方加箭头。这种写法简单又有效。本书中矢量几乎总是列矢量，行矢量则写成列矢量的转置，上标 $t$ 表示转置：
\begin{equation}
 \bm a=\begin{bmatrix}a_1\\a_2\\a_3\end{bmatrix},\quad
 \bm b=\begin{bmatrix}b_1\\b_2\\b_3\end{bmatrix},\quad
 \bm a^t=[a_1,a_2,a_3],\quad \bm b^t=[b_1,b_2,b_3]. \tag{1.2.1}
\end{equation}
二阶张量及其转置例如
\begin{equation}
 \bm A=\begin{bmatrix}A_{11}&A_{12}&A_{13}\\A_{21}&A_{22}&A_{23}\\A_{31}&A_{32}&A_{33}\end{bmatrix},\quad
 \bm A^t=\begin{bmatrix}A_{11}&A_{21}&A_{31}\\A_{12}&A_{22}&A_{32}\\A_{13}&A_{23}&A_{33}\end{bmatrix}. \tag{1.2.2}
\end{equation}
转置就是交换矩阵的行与列。

\subsection{点积、内积与张量双点积}

一对矢量的点积（内积）产生一个标量：
\begin{align}
 \bm a\cdot\bm b&=\bm a^t\bm b=a_1b_1+a_2b_2+a_3b_3, \tag{1.2.3}\\
 a^2&=\bm a\cdot\bm a=\bm a^t\bm a. \tag{1.2.4}
\end{align}
一对张量的双点积为
\begin{equation}
 \bm A:\bm B=\sum_{i,j}A_{ij}B_{ij}. \tag{1.2.5}
\end{equation}
它也常被称为张量内积。两个点对应于对两个指标求和。无论矢量或张量有多少分量，这种运算最后都给出一个实数，而且
\begin{equation}
 \bm a\cdot\bm b=\bm b\cdot\bm a,\qquad \bm A:\bm B=\bm B:\bm A. \tag{1.2.6}
\end{equation}

\subsection{叉积或外积}

矢量叉积产生一个新矢量：
\begin{equation}
 \bm a\times\bm b=
 \begin{vmatrix}\bm e_1&\bm e_2&\bm e_3\\a_1&a_2&a_3\\b_1&b_2&b_3\end{vmatrix}
 =\begin{bmatrix}a_2b_3-a_3b_2\\a_3b_1-a_1b_3\\a_1b_2-a_2b_1\end{bmatrix}. \tag{1.2.7}
\end{equation}
这里 $\bm e_1,\bm e_2,\bm e_3$ 是单位基矢。$\bm a\times\bm b$ 与 $\bm b\times\bm a$ 大小相同、方向相反；这一几何意义十分直观。

\subsection{并矢张量}

并矢积以 $\otimes$ 表示，由两个矢量构成二阶张量：
\begin{equation}
 \bm a\otimes\bm b=\bm a\bm b^t=
 \begin{bmatrix}a_1b_1&a_1b_2&a_1b_3\\a_2b_1&a_2b_2&a_2b_3\\a_3b_1&a_3b_2&a_3b_3\end{bmatrix},
 \qquad (\bm a\otimes\bm b)_{ij}=a_ib_j. \tag{1.2.8--9}
\end{equation}
若 $\bm c=[c_1,c_2,c_3]^t$，也可用
\begin{equation}
 (\bm a\otimes\bm b)\bm c=\bm a(\bm b\cdot\bm c) \tag{1.2.10}
\end{equation}
来定义并矢积。把右端逐分量展开并整理成左端的矩阵形式，便能重新得到上式。一般而言，$\bm a\otimes\bm b\ne\bm b\otimes\bm a$。它不像叉积那样具有立即可见的几何意义，但形式足够简洁。

\subsection{标量三重积}

三个矢量的标量三重积定义为
\begin{equation} \bm a\cdot(\bm b\times\bm c). \tag{1.2.12}\end{equation}
它等于由 $\bm a,\bm b,\bm c$ 张成的平行六面体的有向体积。因此，下列循环置换关系只是“同一体积”的不同写法：
\begin{align}
 \bm a\cdot(\bm b\times\bm c)&=\bm b\cdot(\bm c\times\bm a)=\bm c\cdot(\bm a\times\bm b), \tag{1.2.13}\\
 (\bm b\times\bm c)\cdot\bm a&=(\bm c\times\bm a)\cdot\bm b=(\bm a\times\bm b)\cdot\bm c. \tag{1.2.14}
\end{align}
交换叉积中的两个矢量会改变有向体积的符号，例如
\begin{equation}
 \bm a\cdot(\bm b\times\bm c)=-\bm a\cdot(\bm c\times\bm b). \tag{1.2.15}
\end{equation}
标量三重积也可写成行列式
\begin{equation}
 \bm a\cdot(\bm b\times\bm c)=
 \begin{vmatrix}a_1&a_2&a_3\\b_1&b_2&b_3\\c_1&c_2&c_3\end{vmatrix}. \tag{1.2.18}
\end{equation}
任意两个矢量相同，平行六面体就退化为平面，体积为零，例如 $\bm a\cdot(\bm a\times\bm c)=0$。由三个共顶点棱矢量构成的四面体体积则是
\begin{equation}
 V=\frac16\left|\bm a\cdot(\bm b\times\bm c)\right|. \tag{1.2.20}
\end{equation}
四面体是 CFD 中常用、也很讨人喜欢的计算单元。绝对值保证体积为正。坐标系旋转不改变几何体体积，所以若 $\bm a'=\bm R\bm a$、$\bm b'=\bm R\bm b$、$\bm c'=\bm R\bm c$，则
\begin{equation}
 \bm a\cdot(\bm b\times\bm c)=\bm a'\cdot(\bm b'\times\bm c'). \tag{1.2.21}
\end{equation}

\subsection{矢量三重积}

矢量三重积 $\bm a\times(\bm b\times\bm c)$ 满足 Lagrange 展开式（也称 BAC--CAB 恒等式）：
\begin{equation}
 \bm a\times(\bm b\times\bm c)=\bm b(\bm a\cdot\bm c)-\bm c(\bm a\cdot\bm b). \tag{1.2.24}
\end{equation}
因为 $[\bm a\times(\bm b\times\bm c)]\cdot\bm a=0$，该三重积必与 $\bm a$ 垂直。这在 CFD 中很有用：我们常由计算单元面的法向量构造切向量，或者更一般地，需要寻找与给定矢量正交的方向。

旋转混合 Riemann 求解器就是一个例子。有限体积法通常沿面单位法向 $\hat{\bm n}$ 计算一个数值通量；旋转混合法把 $\hat{\bm n}$ 分解到两个正交方向，在物理上有意义的方向 $\hat{\bm n}_1$（例如激波法向）使用稳健但耗散较大的通量，在其正交方向 $\hat{\bm n}_2$ 使用稳健性略低但耗散小得多的通量。这样既能提高强激波计算的稳健性，又不会在光滑区加入过多耗散。给定 $\hat{\bm n}_1$ 后，可由
\begin{equation}
 \hat{\bm n}_2=\frac{\hat{\bm n}_1\times(\hat{\bm n}\times\hat{\bm n}_1)}
 {|\hat{\bm n}_1\times(\hat{\bm n}\times\hat{\bm n}_1)|} \tag{1.2.26}
\end{equation}
构造单位切向量；显然 $\hat{\bm n}_2\cdot\hat{\bm n}_1=0$。

\subsection{单位矩阵}

单位矩阵对角元全为 1，其余元素为零。例如
\begin{equation}
 \bm I=\begin{bmatrix}1&0&0\\0&1&0\\0&0&1\end{bmatrix}=\operatorname{diag}(1,1,1),\qquad
 \bm I\bm A=\bm A\bm I=\bm A. \tag{1.2.27--29}
\end{equation}
$n\times n$ 单位矩阵含有 $n^2-n$ 个零；用 $\operatorname{diag}$ 表示高维单位矩阵尤其紧凑。

\section{正交矢量}

考虑三个两两正交的单位矢量 $\bm n,\bm\ell,\bm m$：
\begin{gather}
 \bm n\cdot\bm n=\bm\ell\cdot\bm\ell=\bm m\cdot\bm m=1,\qquad
 \bm n\cdot\bm\ell=\bm\ell\cdot\bm m=\bm m\cdot\bm n=0, \tag{1.3.1--2}\\
 \bm n=\bm\ell\times\bm m,\quad \bm\ell=\bm m\times\bm n,\quad \bm m=\bm n\times\bm\ell, \tag{1.3.3--5}\\
 \bm n\otimes\bm n+\bm\ell\otimes\bm\ell+\bm m\otimes\bm m=\bm I. \tag{1.3.6}
\end{gather}
若 $\bm n=[n_x,n_y,n_z]^t$、$\bm\ell=[\ell_x,\ell_y,\ell_z]^t$、$\bm m=[m_x,m_y,m_z]^t$，则上述关系逐分量给出：每个矢量的三个分量平方和为 1；任意两矢量的对应分量乘积之和为 0；叉积给出九个分量关系；而完备性关系给出
\begin{gather}
 n_x^2+\ell_x^2+m_x^2=n_y^2+\ell_y^2+m_y^2=n_z^2+\ell_z^2+m_z^2=1,\\
 n_xn_y+\ell_x\ell_y+m_xm_y=0,\quad
 n_xn_z+\ell_x\ell_z+m_xm_z=0,\quad
 n_yn_z+\ell_y\ell_z+m_ym_z=0. \tag{1.3.14--19}
\end{gather}

对任意矢量 $\bm v=[u,v,w]^t$，定义法向和两个切向分量
\begin{equation}
 q_n=\bm v\cdot\bm n,\qquad q_\ell=\bm v\cdot\bm\ell,\qquad q_m=\bm v\cdot\bm m. \tag{1.3.20}
\end{equation}
则
\begin{gather}
 q_n\bm n+q_\ell\bm\ell+q_m\bm m=\bm v,\qquad
 q_n^2+q_\ell^2+q_m^2=|\bm v|^2, \tag{1.3.21--22}\\
 q_nn_x+q_\ell\ell_x+q_mm_x=u,\quad
 q_nn_y+q_\ell\ell_y+q_mm_y=v,\quad
 q_nn_z+q_\ell\ell_z+q_mm_z=w. \tag{1.3.23--25}
\end{gather}
对这些关系微分即可得到 $dq_n,dq_\ell,dq_m$ 与 $du,dv,dw$ 的对应关系，并有
\begin{equation}
 q_n\,dq_n+q_\ell\,dq_\ell+q_m\,dq_m=u\,du+v\,dv+w\,dw. \tag{1.3.29}
\end{equation}
在三维 CFD 代码中，可把 $\bm v$ 看成速度，把 $\bm n$ 看成单元面法向，把 $\bm\ell,\bm m$ 看成两个切向。上述关系可以从只依赖法向投影的 Jacobian 绝对值等表达式中消去切向矢量，例如 Roe 求解器中就会用到这一技巧。

再考虑对称二阶张量
\begin{equation}
 \bm\tau=\begin{bmatrix}\tau_{xx}&\tau_{xy}&\tau_{xz}\\\tau_{yx}&\tau_{yy}&\tau_{yz}\\\tau_{zx}&\tau_{zy}&\tau_{zz}\end{bmatrix},
 \quad \tau_{yx}=\tau_{xy},\;\tau_{zx}=\tau_{xz},\;\tau_{zy}=\tau_{yz}. \tag{1.3.30}
\end{equation}
定义 $\tau_{nx}=\tau_{xx}n_x+\tau_{xy}n_y+\tau_{xz}n_z$，并类似定义 $\tau_{ny},\tau_{nz}$；再令
\begin{gather}
 \tau_{nn}=\tau_{nx}n_x+\tau_{ny}n_y+\tau_{nz}n_z,\quad
 \tau_{n\ell}=\tau_{nx}\ell_x+\tau_{ny}\ell_y+\tau_{nz}\ell_z,\\
 \tau_{nm}=\tau_{nx}m_x+\tau_{ny}m_y+\tau_{nz}m_z,\quad
 \tau_{nv}=\tau_{nx}u+\tau_{ny}v+\tau_{nz}w. \tag{1.3.31--37}
\end{gather}
利用正交基的完备性可得
\begin{gather}
 \tau_{nn}n_i+\tau_{nm}m_i+\tau_{n\ell}\ell_i=\tau_{ni},\qquad i=x,y,z, \tag{1.3.38--40}\\
 \tau_{nn}q_n+\tau_{nm}q_m+\tau_{n\ell}q_\ell=\tau_{nv}, \tag{1.3.41}\\
 \tau_{nn}d\tau_{nn}+\tau_{nm}d\tau_{nm}+\tau_{n\ell}d\tau_{n\ell}
 =\tau_{nx}d\tau_{nx}+\tau_{ny}d\tau_{ny}+\tau_{nz}d\tau_{nz}, \tag{1.3.42}\\
 \tau_{nn}dq_n+\tau_{nm}dq_m+\tau_{n\ell}dq_\ell
 =\tau_{nx}du+\tau_{ny}dv+\tau_{nz}dw. \tag{1.3.43}
\end{gather}
把第一组关系微分，还可得到 $d\tau_{nn}n_i+d\tau_{nm}m_i+d\tau_{n\ell}\ell_i=d\tau_{ni}$。这些恒等式同样可从法向投影的 Jacobian 绝对值中消去切向矢量，例如上风黏性通量中就会遇到。

\section{指标记号}

指标记号既紧凑，又特别适合证明矢量恒等式。它还省去了反复书写求和符号 $\sum$ 的麻烦。

\subsection{Einstein 求和约定}

Einstein 求和约定规定：同一项中某指标重复出现两次，就对该指标求和。例如
\begin{equation}
 a_ib_i=a_1b_1+a_2b_2+a_3b_3. \tag{1.4.1}
\end{equation}
于是点积可简写为 $\bm a\cdot\bm b=a_ib_i$。矩阵--矢量乘积的第 $i$ 个分量为
\begin{equation}
 (\bm A\bm b)_i=A_{ij}b_j=A_{i1}b_1+A_{i2}b_2+A_{i3}b_3. \tag{1.4.2}
\end{equation}
求和指标是哑指标，换成任何字母都不改变结果：$a_ib_i=a_jb_j=a_kb_k$。

\subsection{Kronecker 符号}

二阶张量 Kronecker delta 定义为
\begin{equation}
 \delta_{ij}=\begin{cases}1,&i=j,\\0,&i\ne j.\end{cases} \tag{1.4.5}
\end{equation}
它正是单位矩阵的分量。由于
\begin{equation}
 \delta_{ij}a_j=a_i, \tag{1.4.7}
\end{equation}
它也常被称为代换算子，好像把 $i$ 代入了 $j$。进一步有
\begin{gather}
 \delta_{ij}a_ib_j=a_ib_i,\qquad \bm I:(\bm a\otimes\bm b)=\bm a^t\bm I\bm b=\bm a\cdot\bm b, \tag{1.4.8--9}\\
 \delta_{ii}=3,\qquad \delta_{ij}\delta_{jk}=\delta_{ik},\qquad \bm I\bm I=\bm I. \tag{1.4.10--12}
\end{gather}

\subsection{Eddington 符号}

三阶 Eddington epsilon 定义为
\begin{equation}
 \varepsilon_{ijk}=\begin{cases}
 0,&i,j,k\text{ 中任意两个相同},\\
 1,&(i,j,k)\text{ 为偶排列},\\
 -1,&(i,j,k)\text{ 为奇排列}.
 \end{cases} \tag{1.4.13}
\end{equation}
因此
\begin{equation}
 \varepsilon_{ijk}=-\varepsilon_{jik}=-\varepsilon_{ikj},\qquad
 \varepsilon_{ijk}=\varepsilon_{jki}=\varepsilon_{kij},\qquad
 \delta_{ij}\varepsilon_{ijk}=0. \tag{1.4.14--16}
\end{equation}
若以三阶张量 $\bm E$ 表示这些分量，其三个切片就像魔方的三层。叉积和行列式可分别写成
\begin{equation}
 \bm a\times\bm b=\bm E(\bm a\otimes\bm b),\qquad
 \det\bm A=\varepsilon_{ijk}A_{1i}A_{2j}A_{3k}. \tag{1.4.20--21}
\end{equation}
并有关键恒等式
\begin{gather}
 \varepsilon_{ijk}=
 \begin{vmatrix}\delta_{i1}&\delta_{i2}&\delta_{i3}\\
 \delta_{j1}&\delta_{j2}&\delta_{j3}\\
 \delta_{k1}&\delta_{k2}&\delta_{k3}\end{vmatrix}, \tag{1.4.24}\\
 \varepsilon_{ijk}\varepsilon_{lmn}=
 \begin{vmatrix}\delta_{il}&\delta_{im}&\delta_{in}\\
 \delta_{jl}&\delta_{jm}&\delta_{jn}\\
 \delta_{kl}&\delta_{km}&\delta_{kn}\end{vmatrix}, \tag{1.4.27}\\
 \varepsilon_{ijk}\varepsilon_{kmn}=\delta_{im}\delta_{jn}-\delta_{in}\delta_{jm},\quad
 \varepsilon_{ijk}\varepsilon_{ijn}=2\delta_{kn},\quad
 \varepsilon_{ijk}\varepsilon_{ijk}=6. \tag{1.4.29--31}
\end{gather}
这些关系可以机械地证明矢量恒等式。例如，写出 $[\bm a\times(\bm b\times\bm c)]_i=\varepsilon_{ijk}a_j\varepsilon_{klm}b_lc_m$，再用 $\varepsilon$--$\delta$ 恒等式收缩指标，即得 Lagrange 公式 $\bm a\times(\bm b\times\bm c)=\bm b(\bm a\cdot\bm c)-\bm c(\bm a\cdot\bm b)$。

\section{散度、梯度与旋度}

\subsection{坐标系}

矢量微分算子在笛卡尔坐标 $(x,y,z)$ 中最简洁，在柱坐标 $(r,\theta,z)$ 和球坐标 $(r,\theta,\phi)$ 中则稍复杂。坐标关系为
\begin{align}
 &\text{柱坐标：}&x&=r\cos\theta,&y&=r\sin\theta,&z&=z, \tag{1.5.1--3}\\
 &\text{球坐标：}&x&=r\sin\phi\cos\theta,&y&=r\sin\phi\sin\theta,&z&=r\cos\phi. \tag{1.5.4--6}
\end{align}
相应单位基矢分别记为 $(\bm e_x,\bm e_y,\bm e_z)$、$(\bm e_r,\bm e_\theta,\bm e_z)$ 和 $(\bm e_r,\bm e_\theta,\bm e_\phi)$；矢量 $\bm a$ 的分量也依次写作 $(a_1,a_2,a_3)$、$(a_r,a_\theta,a_z)$ 和 $(a_r,a_\theta,a_\phi)$。

\subsection{标量函数的梯度}

标量函数 $\alpha$ 的梯度可定义为
\begin{equation}
 \operatorname{grad}\alpha=\lim_{\Delta V\to0}\frac1{\Delta V}\oint\alpha\bm n\,dS,
 \qquad d\alpha=\operatorname{grad}\alpha\cdot d\bm x. \tag{1.5.13--14}
\end{equation}
三个坐标系中的线元分别为
\begin{equation}
 d\bm x=(dx,dy,dz),\quad (dr,r\,d\theta,dz),\quad
 (dr,r\sin\phi\,d\theta,r\,d\phi), \tag{1.5.15--17}
\end{equation}
从而
\begin{align}
 \operatorname{grad}\alpha&=\alpha_x\bm e_x+\alpha_y\bm e_y+\alpha_z\bm e_z, \tag{1.5.18}\\
 &=\alpha_r\bm e_r+\frac1r\alpha_\theta\bm e_\theta+\alpha_z\bm e_z, \tag{1.5.19}\\
 &=\alpha_r\bm e_r+\frac1{r\sin\phi}\alpha_\theta\bm e_\theta+\frac1r\alpha_\phi\bm e_\phi. \tag{1.5.20}
\end{align}
这些式子相当直观。要推导它们及以下公式，可以把积分定义用于一个很小的体积元，或直接展开微分定义。

\subsection{矢量函数的散度}

散度定义为
\begin{equation}
 \operatorname{div}\bm a=\lim_{\Delta V\to0}\frac1{\Delta V}\oint\bm a\cdot\bm n\,dS. \tag{1.5.21}
\end{equation}
三个坐标系中分别为
\begin{align}
 \operatorname{div}\bm a&=\frac{\partial a_1}{\partial x}+\frac{\partial a_2}{\partial y}+\frac{\partial a_3}{\partial z}, \tag{1.5.22}\\
 &=\frac1r\frac{\partial(ra_r)}{\partial r}+\frac1r\frac{\partial a_\theta}{\partial\theta}+\frac{\partial a_z}{\partial z}, \tag{1.5.23}\\
 &=\frac1{r^2}\frac{\partial(r^2a_r)}{\partial r}+\frac1{r\sin\phi}\frac{\partial a_\theta}{\partial\theta}
 +\frac1{r\sin\phi}\frac{\partial(a_\phi\sin\phi)}{\partial\phi}. \tag{1.5.24}
\end{align}

\subsection{矢量函数的旋度}

旋度定义为
\begin{equation}
 \operatorname{curl}\bm a=-\lim_{\Delta V\to0}\frac1{\Delta V}\oint\bm a\times\bm n\,dS. \tag{1.5.25}
\end{equation}
笛卡尔坐标下
\begin{equation}
 \operatorname{curl}\bm a=
 \begin{vmatrix}\bm e_x&\bm e_y&\bm e_z\\\partial_x&\partial_y&\partial_z\\a_1&a_2&a_3\end{vmatrix}. \tag{1.5.26}
\end{equation}
柱坐标和球坐标的完整行列式形式及逐分量展开如下，保留原式以确保各尺度因子位置无误：
\sourcecrop{25}{50}{150}{562}{620}{0.62}
\eqnote{$a_r,a_\theta,a_z$ 或 $a_r,a_\theta,a_\phi$ 是对应曲线坐标下的矢量分量；$r\sin\phi$ 等因子来自坐标尺度因子。}

\subsection{矢量函数的梯度}

与标量情形类似，定义
\begin{equation}
 \operatorname{grad}\bm a=\lim_{\Delta V\to0}\frac1{\Delta V}\oint\bm a\otimes\bm n\,dS,
 \qquad d\bm a=(\operatorname{grad}\bm a)d\bm x. \tag{1.5.32--33}
\end{equation}
笛卡尔坐标中，$\operatorname{grad}\bm a$ 就是分量对坐标的 Jacobian 矩阵。在柱坐标和球坐标中还会出现几何附加项；三种坐标下的完整矩阵为：
\sourcecrop{26}{48}{330}{564}{700}{0.55}
\eqnote{曲线坐标中的“附加项”源于单位基矢随空间变化。例如柱坐标有 $\partial_\theta\bm e_r=\bm e_\theta$、$\partial_\theta\bm e_\theta=-\bm e_r$；球坐标的各基矢还同时随 $\theta$ 与 $\phi$ 变化。}

以柱坐标为例，把 $\bm a=a_r\bm e_r+a_\theta\bm e_\theta+a_z\bm e_z$ 微分；对分量和基矢同时使用乘积法则，再利用上述基矢导数，可整理成
\begin{equation}
 d\bm a=
 \begin{bmatrix}
 \partial_ra_r&r^{-1}\partial_\theta a_r-a_\theta/r&\partial_za_r\\
 \partial_ra_\theta&r^{-1}\partial_\theta a_\theta+a_r/r&\partial_za_\theta\\
 \partial_ra_z&r^{-1}\partial_\theta a_z&\partial_za_z
 \end{bmatrix}
 \begin{bmatrix}dr\\r\,d\theta\\dz\end{bmatrix}. \tag{1.5.41}
\end{equation}
矩阵就是柱坐标下的矢量梯度。球坐标可按完全相同的步骤推导。下面的张量散度和矢量 Laplace 算子也会出现同类附加项。

\subsection{张量函数的散度}

二阶张量 $\bm A$ 的散度定义为
\begin{equation}
 \operatorname{div}\bm A=\lim_{\Delta V\to0}\frac1{\Delta V}\oint\bm A\bm n\,dS. \tag{1.5.42}
\end{equation}
笛卡尔坐标中，第 $i$ 个分量为 $\partial_xA_{i1}+\partial_yA_{i2}+\partial_zA_{i3}$。柱坐标和球坐标中的完整式为：
\sourcecrop{28}{48}{360}{564}{700}{0.52}
\eqnote{$A_{rr},A_{r\theta},\ldots$ 是张量在局部正交基上的分量；含 $1/r$、$1/(r\tan\phi)$ 的项都是基矢变化产生的几何项。}

柱坐标推导可从微小体积 $V=dr\,(r\,d\theta)\,dz$ 开始。把相对两个面上的张量通量作 Taylor 展开，并把随坐标变化的面矢量一同展开。三个方向的面积分别为 $dS_r=r\,d\theta\,dz$、$dS_\theta=dr\,dz$、$dS_z=r\,d\theta\,dr$。径向项给出 $r^{-1}\partial_r(rA_{ir})$；周向项除 $r^{-1}\partial_\theta A_{i\theta}$ 外，还因 $\partial_\theta\bm e_r=\bm e_\theta$、$\partial_\theta\bm e_\theta=-\bm e_r$ 产生 $-A_{\theta\theta}/r$ 和 $+A_{r\theta}/r$；轴向项只是普通的 $\partial_zA_{iz}$。三部分相加便得到上面的柱坐标公式。可见，所有附加项都来自空间变化的单位基矢。

\subsection{标量函数的 Laplace 算子}

标量 Laplace 算子是梯度的散度：
\begin{align}
 \operatorname{div}\operatorname{grad}\alpha&=\alpha_{xx}+\alpha_{yy}+\alpha_{zz}, \tag{1.5.56}\\
 &=\frac1r\frac{\partial}{\partial r}\left(r\alpha_r\right)+\frac1{r^2}\alpha_{\theta\theta}+\alpha_{zz}, \tag{1.5.57}\\
 &=\frac1{r^2}\frac{\partial}{\partial r}\left(r^2\alpha_r\right)+\frac1{r^2\sin^2\phi}\alpha_{\theta\theta}
 +\frac1{r^2\sin\phi}\frac{\partial}{\partial\phi}\left(\sin\phi\,\alpha_\phi\right). \tag{1.5.58}
\end{align}
三式依次对应笛卡尔、柱和球坐标。常用经济记号为 $\Delta\alpha=\operatorname{div}\operatorname{grad}\alpha$。

\subsection{矢量函数的 Laplace 算子}

矢量 Laplace 算子同样定义为梯度的散度。笛卡尔坐标中，只需分别对三个分量取标量 Laplace；柱坐标和球坐标中则有附加耦合项：
\sourcecrop{30}{47}{300}{565}{680}{0.55}
\eqnote{式中的 $\operatorname{div}\operatorname{grad}a_r$ 等表示先把相应分量当作标量，再使用该坐标系的标量 Laplace 公式。}

\subsection{物质导数}

标量的物质导数在三个坐标系中分别为
\begin{align}
 \frac{D\alpha}{Dt}&=\alpha_t+u\alpha_x+v\alpha_y+w\alpha_z, \tag{1.5.63}\\
 &=\alpha_t+v_r\alpha_r+\frac{v_\theta}{r}\alpha_\theta+v_z\alpha_z, \tag{1.5.64}\\
 &=\alpha_t+v_r\alpha_r+\frac{v_\theta}{r\sin\phi}\alpha_\theta+\frac{v_\phi}{r}\alpha_\phi. \tag{1.5.65}
\end{align}
笛卡尔坐标中矢量物质导数只需逐分量计算；柱坐标和球坐标中还会出现基矢随运动变化产生的项：
\sourcecrop{31}{48}{475}{564}{700}{0.34}
\eqnote{$D a_r/Dt$ 等只表示把该分量作为标量按上一组公式求物质导数；矩阵中的其余项是曲线坐标基矢的几何贡献。}

\subsection{复杂片状流与 Beltrami 流}

满足
\begin{equation} \bm a\cdot(\operatorname{curl}\bm a)=0 \tag{1.5.69}\end{equation}
的矢量场称为\textbf{复杂片状场}（complex lamellar field）。若 $\bm a$ 是流速，这意味着速度与涡量相互正交，相应流动称为复杂片状流。另一方面，满足
\begin{equation} \bm a\times(\operatorname{curl}\bm a)=0 \tag{1.5.70}\end{equation}
的矢量场称为 \textbf{Beltrami 场}；若 $\bm a$ 是速度，则速度与涡量平行，称为 Beltrami 流。前者让作者联想到升力翼型绕流，后者则让人想到飞机尾涡。矢量微分算子的更多细节可参见专门教材。

\section{del 算子}

笛卡尔坐标中的 del（符号名称为 nabla）定义为
\begin{equation}
 \nabla=[\partial_x,\partial_y,\partial_z]. \tag{1.6.1}
\end{equation}
有人不喜欢它，理由是 $\nabla$ 依赖坐标系、物理含义不如散度、梯度、旋度清楚；不过它依然非常有用：
\begin{equation}
 \operatorname{grad}\alpha=\nabla\alpha,\quad
 \operatorname{div}\bm a=\nabla\cdot\bm a,\quad
 \operatorname{curl}\bm a=\nabla\times\bm a,\quad
 \operatorname{grad}\bm a=\nabla\bm a,\quad
 \operatorname{div}\bm A=\nabla\cdot\bm A. \tag{1.6.2--6}
\end{equation}
把它像矢量一样形式运算，可以写出
\begin{gather}
 \nabla\cdot\nabla\alpha=\nabla^2\alpha=\Delta\alpha,\quad
 \nabla\times\nabla\alpha,\quad \nabla(\nabla\cdot\bm a),\\
 \nabla\cdot(\nabla\times\bm a),\quad
 \nabla\times(\nabla\times\bm a),\quad
 (\operatorname{grad}\bm a)\bm b=(\bm b\cdot\nabla)\bm a. \tag{1.6.7--12}
\end{gather}
但必须牢记 $\nabla$ 本质上是算子：虽然 $\bm a\cdot\bm b=\bm b\cdot\bm a$，却一般有 $\bm a\cdot\nabla\ne\nabla\cdot\bm a$。

\section{矢量恒等式}

与其死记，不如用指标记号证明。常用恒等式如下：
\sourcecrop{32}{43}{90}{569}{610}{0.72}
\eqnote{$\alpha,\beta$ 为标量函数，$\bm a,\bm b,\bm c$ 为矢量函数，$\bm A$ 为二阶张量，$\bm E$ 为 Eddington epsilon 张量。原式中的 grad、div、curl 分别就是梯度、散度和旋度。}

例如，乘积散度的分量展开为
\begin{equation}
 [\operatorname{div}(\alpha\bm a)]_j=\frac{\partial(\alpha a_j)}{\partial x_j}
 =\alpha\frac{\partial a_j}{\partial x_j}+\frac{\partial\alpha}{\partial x_j}a_j,
\end{equation}
故 $\operatorname{div}(\alpha\bm a)=\alpha\operatorname{div}\bm a+\operatorname{grad}\alpha\cdot\bm a$。再如，把 $[\bm a\times\operatorname{curl}\bm b]_i$ 写成 $\varepsilon_{ijk}a_j\varepsilon_{kmn}\partial_mb_n$，用 $\varepsilon$--$\delta$ 恒等式收缩，就得到
\begin{equation}
 \bm a\times\operatorname{curl}\bm b=(\operatorname{grad}\bm b)^t\bm a-(\operatorname{grad}\bm b)\bm a.
\end{equation}
自己用同一方法验证其余恒等式，会比记忆它们更有趣。

\section{矢量的坐标变换}

在三个坐标系之间来回切换很常见；流体速度也经常需要从一个坐标系转换到另一个坐标系。令速度矢量在笛卡尔、柱和球坐标中的分量依次为
\begin{equation}
 \bm v=[u,v,w]^t,\qquad [u_r,u_\theta,u_z]^t,\qquad [u_r,u_\theta,u_\phi]^t. \tag{1.8.1--3}
\end{equation}
柱坐标与笛卡尔坐标之间的变换为
\begin{gather}
 \begin{bmatrix}u\\v\\w\end{bmatrix}=
 \begin{bmatrix}\cos\theta&-\sin\theta&0\\\sin\theta&\cos\theta&0\\0&0&1\end{bmatrix}
 \begin{bmatrix}u_r\\u_\theta\\u_z\end{bmatrix}, \tag{1.8.4}\\
 \begin{bmatrix}u_r\\u_\theta\\u_z\end{bmatrix}=
 \begin{bmatrix}\cos\theta&\sin\theta&0\\-\sin\theta&\cos\theta&0\\0&0&1\end{bmatrix}
 \begin{bmatrix}u\\v\\w\end{bmatrix}. \tag{1.8.5}
\end{gather}
球坐标与笛卡尔坐标之间则有
\sourcecrop{34}{54}{500}{558}{700}{0.30}
\eqnote{正变换矩阵的列分别是 $\bm e_r,\bm e_\theta,\bm e_\phi$ 在笛卡尔基底中的分量；由于它是正交矩阵，逆变换就是其转置。}

\section{偏导数的坐标变换}

柱坐标满足 $x=r\cos\theta$、$y=r\sin\theta$，而 $r=(x^2+y^2)^{1/2}$、$\cos\theta=x/r$、$\sin\theta=y/r$。链式法则给出
\begin{equation}
 \begin{bmatrix}\partial_x\\\partial_y\\\partial_z\end{bmatrix}=
 \begin{bmatrix}
 \cos\theta&-\sin\theta/r&0\\
 \sin\theta&\cos\theta/r&0\\0&0&1
 \end{bmatrix}
 \begin{bmatrix}\partial_r\\\partial_\theta\\\partial_z\end{bmatrix}, \tag{1.9.11}
\end{equation}
其逆变换为
\begin{equation}
 \begin{bmatrix}\partial_r\\\partial_\theta\\\partial_z\end{bmatrix}=
 \begin{bmatrix}
 \cos\theta&\sin\theta&0\\-r\sin\theta&r\cos\theta&0\\0&0&1
 \end{bmatrix}
 \begin{bmatrix}\partial_x\\\partial_y\\\partial_z\end{bmatrix}. \tag{1.9.12}
\end{equation}

球坐标满足 $x=r\sin\phi\cos\theta$、$y=r\sin\phi\sin\theta$、$z=r\cos\phi$。逐项使用链式法则，可得完整正、逆变换：
\sourcecrop{35}{46}{55}{566}{190}{0.28}
\sourcecrop{36}{50}{600}{562}{712}{0.22}
\eqnote{源式 (1.9.24) 的左端应理解为 $\partial\phi/\partial x$、$\partial\phi/\partial y$、$\partial\phi/\partial z$；文字抽取时容易误成 $\partial\theta$。}
借助这些变换，可以由笛卡尔坐标公式推导柱、球坐标中的梯度等矢量算子。理解不同坐标系之间的关系总是有益的。

\section{特征值与特征矢量}

对方阵 $\bm A$，若标量 $\lambda$ 与非零列矢量 $\bm r$ 满足
\begin{equation}
 (\bm A-\lambda\bm I)\bm r=0, \tag{1.10.1}
\end{equation}
则 $\lambda$ 称为 $\bm A$ 的特征值，$\bm r$ 是相应的右特征矢量。若行矢量 $\bm\ell^t$ 满足
\begin{equation}
 \bm\ell^t(\bm A-\lambda\bm I)=0, \tag{1.10.2}
\end{equation}
则称其为左特征矢量。特征矢量并不唯一，任意非零标量倍数仍是同一特征值的特征矢量，因此可以按需要归一化。

$n\times n$ 矩阵有 $n$ 个特征值（计代数重数）$\lambda_1,\ldots,\lambda_n$。若它们互异，相应特征矢量必线性无关；即便存在重特征值，也仍可能得到完备的线性无关特征矢量。当右特征矢量完备时，矩阵可对角化：
\begin{equation}
 \bm\Lambda=\bm R^{-1}\bm A\bm R,\qquad
 \bm R=[\bm r_1,\bm r_2,\ldots,\bm r_n],\qquad
 \bm\Lambda=\operatorname{diag}(\lambda_1,\ldots,\lambda_n). \tag{1.10.4--6}
\end{equation}
并且 $\bm L=\bm R^{-1}$ 正是左特征矢量矩阵，其第 $k$ 行是 $\bm\ell_k^t$。给定 $\bm R$ 后，$\bm L$ 被唯一确定。若 $\bm A$ 对称，则所有特征值均为实数，不同特征值对应的特征矢量彼此正交。

\section{相似矩阵}

若同阶方阵 $\bm A,\bm B,\bm M$ 满足
\begin{equation}
 \bm B=\bm M^{-1}\bm A\bm M, \tag{1.11.1}
\end{equation}
则 $\bm A$ 与 $\bm B$ 相似，上式称为相似变换。相似矩阵具有相同特征值；对某一特征值，其特征矢量满足
\begin{equation}
 \bm r_B=\bm M^{-1}\bm r_A. \tag{1.11.2}
\end{equation}
因此，只要已知 $\bm A$ 的特征矢量，就能直接得到任何相似矩阵的特征矢量。分析不同变量形式的守恒律时，这一点特别有用：变量变换总会使 Jacobian 矩阵发生相似变换，因此只需完整分析一种方程形式。

\section{散度定理}

散度定理为
\begin{equation}
 \int_V\operatorname{div}\bm a\,dV=\oint_S\bm a\cdot\bm n\,dS, \tag{1.12.1}
\end{equation}
其中 $\bm n$ 是体积 $V$ 的边界面 $S$ 上的单位外法向。在二维中，它常以 Green--Gauss 或 Green 定理的形式写成
\begin{equation}
 \int_V\left(\frac{\partial a_x}{\partial x}+\frac{\partial a_y}{\partial y}\right)dx\,dy
 =\oint_S(a_x\,dy-a_y\,dx), \tag{1.12.2}
\end{equation}
其中 $(dx,dy)$ 沿边界逆时针取向。另一常见形式
\begin{equation}
 \int_V\operatorname{grad}\alpha\,dV=\oint_S\alpha\bm n\,dS \tag{1.12.3}
\end{equation}
更准确地说应称为梯度定理；它的每个分量都可看成散度定理的一部分。

散度定理也可理解为分部积分。由 $\operatorname{div}(\alpha\bm a)=\alpha\operatorname{div}\bm a+\operatorname{grad}\alpha\cdot\bm a$ 得
\begin{equation}
 \int_V\alpha\operatorname{div}\bm a\,dV=\oint_S\alpha\bm a\cdot\bm n\,dS
 -\int_V\operatorname{grad}\alpha\cdot\bm a\,dV. \tag{1.12.5}
\end{equation}
对二阶张量同样有 $\int_V\operatorname{div}\bm A\,dV=\oint_S\bm A\bm n\,dS$。把散度定理中的 $\bm a$ 换成常矢量 $\bm c$ 与 $\bm b$ 的叉积，结合矢量恒等式，还可推得
\begin{equation}
 \int_V\operatorname{curl}\bm b\,dV=-\oint_S(\bm b\times\bm n)\,dS. \tag{1.12.10}
\end{equation}
这些公式在 CFD 中经常用于计算导数、降低微分方程阶数（例如 Laplace 方程的 Galerkin 形式），以及计算网格单元体积。

\section{Stokes 定理}

旋度定理更常以 Stokes 定理之名出现：
\begin{equation}
 \int_S\operatorname{curl}\bm a\cdot\bm n\,dS=\oint_{\partial S}\bm a\cdot d\bm\ell, \tag{1.13.1}
\end{equation}
其中 $S$ 是有界曲面，$d\bm\ell=(dx,dy,dz)$ 沿边界 $\partial S$ 逆时针取向。笛卡尔分量展开为
\begin{multline}
 \int_S\left[\left(a_{z,y}-a_{y,z}\right)dy\,dz+
 \left(a_{x,z}-a_{z,x}\right)dz\,dx+
 \left(a_{y,x}-a_{x,y}\right)dx\,dy\right]\\
 =\oint_{\partial S}(a_x\,dx+a_y\,dy+a_z\,dz). \tag{1.13.2}
\end{multline}
令 $a_z=0$，并令其余分量与 $z$ 无关，就立即得到 Green 定理。再把 $\bm a$ 换成 $\bm c\times\bm b$，其中 $\bm c$ 为常矢量，可推得
\begin{equation}
 \int_S\left[\bm n\operatorname{div}\bm b-(\operatorname{grad}\bm b)^t\bm n\right]dS
 =\oint_{\partial S}(\bm b\times d\bm\ell). \tag{1.13.7}
\end{equation}
在流体力学中，Stokes 定理自然出现在 Kelvin 环量定理和三维流函数理论中。更一般地说，旋度定理和散度定理都可视为广义 Stokes（Kelvin--Stokes）定理的特例。

\section{守恒律}

许多流动方程具有守恒律形式：控制体内某个量的时间变化，由穿过控制体边界的净通量来平衡；质量守恒是典型例子。设 $V$ 为控制体，$S$ 为其边界，$\bm U$ 为守恒变量矢量，$\bm F_n$ 为沿边界法向的通量矢量，则
\begin{equation}
 \frac{d}{dt}\int_V\bm U\,dV+\oint_S\bm F_n\,dS=0. \tag{1.14.1}
\end{equation}
这就是守恒律的积分形式：$\bm U$ 在控制体内的积分值，只能因边界通量的净效应而改变。

笛卡尔坐标中，若 $\bm n=(n_x,n_y,n_z)$，则
\begin{equation}
 \bm F_n=\bm F n_x+\bm G n_y+\bm H n_z=\boldsymbol{\mathsf F}\bm n,
 \qquad \boldsymbol{\mathsf F}=[\bm F,\bm G,\bm H]. \tag{1.14.2--3}
\end{equation}
在通量可微、控制体固定的条件下，利用散度定理可得
\begin{equation}
 \int_V\left(\frac{\partial\bm U}{\partial t}+\operatorname{div}\boldsymbol{\mathsf F}\right)dV=0.
\end{equation}
由于它对任意控制体成立，得到微分形式
\begin{equation}
 \frac{\partial\bm U}{\partial t}+\operatorname{div}\boldsymbol{\mathsf F}=0,\qquad
 \frac{\partial\bm U}{\partial t}+\frac{\partial\bm F}{\partial x}
 +\frac{\partial\bm G}{\partial y}+\frac{\partial\bm H}{\partial z}=0. \tag{1.14.6--7}
\end{equation}

积分形式允许不连续解，更一般地允许弱解，而经典微分形式要求可微，因此积分形式更基础、更一般，尤其适合发展激波捕捉方法。它还适用于随时间变形的控制体和任意形状控制体，是有限体积法的基础。

微分形式也可由散度定义写成守恒积分近似：
\begin{equation}
 \frac{\partial\bm U}{\partial t}+\frac1V\oint_S\boldsymbol{\mathsf F}\bm n\,dS\approx0. \tag{1.14.9}
\end{equation}
有限 $V$ 时它是近似，但在节点中心有限体积法中广泛使用，可看成在节点处以守恒方式求解微分形式。对于二阶方法，这通常是很好的近似。

若希望得到精确的表面积分表示，可引入物理通量的原函数 $\boldsymbol{\mathsf G}$：
\begin{equation}
 \boldsymbol{\mathsf F}=\frac1V\int_V\boldsymbol{\mathsf G}(\xi,\eta,\zeta)\,d\xi\,d\eta\,d\zeta,\qquad
 \operatorname{div}\boldsymbol{\mathsf F}=\frac1V\oint_S\boldsymbol{\mathsf G}\bm n\,dS. \tag{1.14.10--11}
\end{equation}
于是对固定的有限体积，
\begin{equation}
 \frac{\partial\bm U}{\partial t}+\frac1V\oint_S\boldsymbol{\mathsf G}\bm n\,dS=0 \tag{1.14.12}
\end{equation}
是精确的。边界上需要的不是物理通量本身，而是其原函数。这一形式是规则网格高阶有限差分方法的基础。

在一维区间 $[x_j-\Delta x/2,x_j+\Delta x/2]$ 上，三种形式可清楚比较。

\subsection*{积分形式}
\begin{equation}
 \frac{d\overline{\bm U}_j}{dt}+\frac{\bm F_{j+1/2}-\bm F_{j-1/2}}{\Delta x}=0,\qquad
 \overline{\bm U}_j=\frac1{\Delta x}\int_{x_j-\Delta x/2}^{x_j+\Delta x/2}\bm U\,dx. \tag{1.14.13--14}
\end{equation}
界面通量由界面解计算：$\bm F_{j\pm1/2}=\bm F(\bm U_{j\pm1/2})$。这一形式本身是精确的，高阶格式可通过提高界面通量以及二维、三维面积分求积的精度来构造，因此它在 CFD 中极为常用。

\subsection*{近似守恒微分形式}
在点 $x_j$ 处应用上面的守恒近似，得到相同的通量差形式，但 $\bm U_j$ 是点值而非体积平均值。它等价于 $\partial_x\bm F$ 的中心差分，所以本质上仅二阶；界面通量本身做得更高阶，并不自动使整条方程高阶。若要更高阶，时间导数和源项也必须作相容的高阶离散。

\subsection*{精确守恒微分形式}
采用通量原函数，则
\begin{equation}
 \frac{d\bm U_j}{dt}+\frac{\bm F^G_{j+1/2}-\bm F^G_{j-1/2}}{\Delta x}=0, \tag{1.14.17}
\end{equation}
其中
\begin{equation}
 \bm F(x)=\frac1{\Delta x}\int_{x-\Delta x/2}^{x+\Delta x/2}\bm F^G(\xi)\,d\xi,\qquad
 \bm F_x=\frac{\bm F^G_{j+1/2}-\bm F^G_{j-1/2}}{\Delta x}. \tag{1.14.18}
\end{equation}
这是精确形式，因此提高 $\bm F^G$ 的界面精度会直接提高整条方程的精度。二阶情况下，三种形式彼此等价，因为体积平均值与中心点值只差 $O(\Delta x^2)$。这也解释了为什么二阶 CFD 代码既可说在计算体积平均，也可说在计算点值，取决于初值以哪一种量给定。

\section{广义坐标中的守恒律}

坐标变换
\begin{equation}
 \xi=\xi(x,y,z),\qquad \eta=\eta(x,y,z),\qquad \zeta=\zeta(x,y,z) \tag{1.15.1}
\end{equation}
把物理空间中的曲线网格映射为计算空间 $(\xi,\eta,\zeta)$ 中的笛卡尔网格。链式法则为
\begin{equation}
 \partial_x=\xi_x\partial_\xi+\eta_x\partial_\eta+\zeta_x\partial_\zeta,
\quad \partial_y=\xi_y\partial_\xi+\eta_y\partial_\eta+\zeta_y\partial_\zeta,
\quad \partial_z=\xi_z\partial_\xi+\eta_z\partial_\eta+\zeta_z\partial_\zeta. \tag{1.15.2--4}
\end{equation}
当逆映射 $x=x(\xi,\eta,\zeta)$、$y=y(\xi,\eta,\zeta)$、$z=z(\xi,\eta,\zeta)$ 已知时，所有度量系数和 Jacobian $J$ 可由下式计算：
\sourcecrop{42}{46}{390}{566}{600}{0.38}
\eqnote{$J$ 是从物理坐标到计算坐标变换的 Jacobian；$x_\xi,y_\eta$ 等表示物理坐标对计算坐标的偏导。}
守恒律变换为
\begin{equation}
 \frac{\partial\bm U^\star}{\partial t}+\frac{\partial\bm F^\star}{\partial\xi}
 +\frac{\partial\bm G^\star}{\partial\eta}+\frac{\partial\bm H^\star}{\partial\zeta}=0, \tag{1.15.10}
\end{equation}
其中
\begin{align}
 \bm U^\star&=\bm U/J,\\
 \bm F^\star&=(\bm F\xi_x+\bm G\xi_y+\bm H\xi_z)/J,\\
 \bm G^\star&=(\bm F\eta_x+\bm G\eta_y+\bm H\eta_z)/J,\\
 \bm H^\star&=(\bm F\zeta_x+\bm G\zeta_y+\bm H\zeta_z)/J. \tag{1.15.11--14}
\end{align}
若通量还依赖 $\bm U$ 的导数，这些导数也按链式法则变换。这样便可在规则计算网格上直接使用标准有限差分公式。

\section{微分方程组的分类}

\subsection{一般方程组}

考虑 $n$ 个拟线性偏微分方程组成的系统
\begin{equation}
 \bm U_t+\bm A\bm U_x+\bm B\bm U_y+\bm C\bm U_z=0, \tag{1.16.1}
\end{equation}
其中 $\bm U$ 含 $n$ 个未知量，$\bm A,\bm B,\bm C$ 为 $n\times n$ 矩阵。该系统可能存在平面波解
\begin{equation}
 \bm U=f(\bm x\cdot\bm n-\lambda_nt)\bm U_0, \tag{1.16.2}
\end{equation}
其中 $f$ 为标量函数，$\bm n=(n_x,n_y,n_z)$ 是波面单位法向，$\lambda_n$ 是沿 $\bm n$ 的波速，$\bm U_0$ 为常矢量。一维时，$f(x-\lambda t)$ 在 $(x,t)$ 平面上沿斜率 $dx/dt=\lambda$ 的曲线保持常值；这条曲线称为特征曲线，$f$ 称为 Riemann 不变量。二维时可写成 $f(x\cos\theta+y\sin\theta-\lambda(\theta)t)$，不同方向的波速共同决定波前包络。

将平面波代入方程得到
\begin{equation}
 (\bm A n_x+\bm B n_y+\bm C n_z-\lambda_n\bm I)\bm U_0=0. \tag{1.16.5}
\end{equation}
因此，$\lambda_n$ 是投影 Jacobian
\begin{equation}
 \bm A_n=\bm A n_x+\bm B n_y+\bm C n_z \tag{1.16.6}
\end{equation}
的特征值，$\bm U_0$ 是相应特征矢量。如果对任意方向 $\bm n$，$\bm A_n$ 的全部特征值均为实数且特征矢量完备，则系统关于时间是双曲型的；若特征值互异，则称严格双曲。若全部特征值为复数，则为椭圆型；既有实数又有复数时为混合型；若实特征值不足 $n$ 个，通常称为抛物型。

平面波的存在意味着：空间某点的解只依赖于波到达该点以前经过的区域，这一区域称为\textbf{依赖域}。典型一维 $2\times2$ 双曲系统有左行波和右行波，因而一点的依赖域是在 $(x,t)$ 平面中由两条特征线围成的下三角区域。反过来，从一点出发能够影响的区域称为\textbf{影响域}，也由穿过该点的特征线界定。双曲问题在无限域或周期域上只给定初值即可，称初值问题；有界域上若某些特征线无法回溯到初始面，还须在相应边界给定部分数据，成为初边值问题。

椭圆系统没有平面波，通常来自稳态问题。既然不存在限定信息传播范围的波，其任一点解便依赖整个区域，因而必须在整个边界给出适当条件，这就是边值问题。抛物问题通常同时需要初值和边界值；某点不仅依赖全部早期时刻的解，还依赖当前时刻整个空间的解。

若 $\bm A,\bm B,\bm C$ 能通过变量变换同时对称化，系统就一定是双曲型；物理上合理的系统通常具有这种可对称化性质，例如 Euler 方程。若它们还能同时对角化，系统便化为互不耦合的一组标量平流方程。这在只有一个系数矩阵的一维问题中可行，但在二维和三维中一般非常困难。

\subsection{一维方程}

考虑一维双曲系统
\begin{equation}
 \bm U_t+\bm A\bm U_x=0. \tag{1.16.7}
\end{equation}
左乘 $\bm A$ 的左特征矢量矩阵 $\bm L$，并用 $\bm L\bm A\bm R=\bm\Lambda$，得到
\begin{equation}
 \bm L\bm U_t+\bm\Lambda\bm L\bm U_x=0. \tag{1.16.9}
\end{equation}
线性系统中 $\bm L$ 为常矩阵，因此可写成
\begin{equation}
 (\bm L\bm U)_t+\bm\Lambda(\bm L\bm U)_x=0. \tag{1.16.10}
\end{equation}
$\bm V=\bm L\bm U$ 是 Riemann 不变量或特征变量，第 $k$ 个分量满足
\begin{equation}
 \frac{\partial v_k}{\partial t}+\lambda_k\frac{\partial v_k}{\partial x}=0. \tag{1.16.11}
\end{equation}
也就是说，它只是以第 $k$ 个特征值给定的速度被平流。把方程组分解成标量平流方程，会使分析和数值处理容易得多。

非线性系统中，Riemann 不变量通常只能以微分形式 $\bm L\,d\bm U=d\bm V$ 表示，很难解析积分。任意解的变化可分解成
\begin{equation}
 d\bm U=\bm R\,d\bm V=\sum_{k=1}^n\bm r_k\,dV_k. \tag{1.16.13}
\end{equation}
因此第 $k$ 个右特征矢量指出第 $k$ 个特征变量（或第 $k$ 支波）变化时，哪些物理变量会随之变化。例如一维 Euler 方程以原始变量 $\bm U=[\rho,u,p]^t$ 表示时，有一支波的特征矢量为 $[1,0,0]^t$，说明它只改变密度；这就是熵波。

定义投影矩阵
\begin{equation}
 \bm\Pi_k=\bm r_k\bm\ell_k^t, \tag{1.16.16}
\end{equation}
则
\begin{equation}
 d\bm U=\sum_{k=1}^n\bm\Pi_kd\bm U,\qquad
 \sum_{k=1}^n\bm\Pi_k=\bm I,\qquad
 \bm A=\sum_{k=1}^n\lambda_k\bm\Pi_k. \tag{1.16.15--18}
\end{equation}
$\bm\Pi_k$ 把解的变化投影到第 $k$ 个特征矢量张成的子空间。于是 Jacobian 的负、正特征部分可写成
\begin{equation}
 \bm A^- =\sum_{\lambda_k\le0}\lambda_k\bm\Pi_k,\qquad
 \bm A^+ =\sum_{\lambda_k>0}\lambda_k\bm\Pi_k. \tag{1.16.19}
\end{equation}
这类分裂常用于构造朝上游一侧取值的上风格式。

守恒律还可写成
\begin{equation}
 \bm U_t+\sum_{k=1}^n\lambda_k\frac{\partial v_k}{\partial x}\bm r_k=0. \tag{1.16.22}
\end{equation}
它清楚显示每个标量平流子问题如何贡献于守恒变量的时间变化，并且每一贡献都沿相应右特征矢量方向。

\subsection{二维稳态方程}

考虑二维稳态系统
\begin{equation}
 \bm A\bm U_x+\bm B\bm U_y=0. \tag{1.16.23}
\end{equation}
若
\begin{equation}
 (\bm A n_x+\bm B n_y)\bm U_0=0, \tag{1.16.25}
\end{equation}
它允许稳态平面波 $\bm U=f(xn_x+yn_y)\bm U_0$。等价地，广义特征值问题
\begin{equation}
 \det(\bm B-\lambda\bm A)=0,\qquad \lambda=-n_x/n_y \tag{1.16.26}
\end{equation}
须有实根。对 $n$ 方程系统，若 $n$ 个特征值全为实数且特征矢量完备，则为双曲型；全为复数则为椭圆型；二者兼有则为混合型；实特征值不足 $n$ 个则为抛物型。

以一般二阶标量方程
\begin{equation}
 a\phi_{xx}+b\phi_{xy}+c\phi_{yy}+d\phi_x+e\phi_y+f=0 \tag{1.16.27}
\end{equation}
为例，令 $u=\phi_x$、$v=\phi_y$，并把混合导数对称地拆分，可写成 $\bm A\bm U_x+\bm B\bm U_y=\bm S$，其中
\begin{equation}
 \bm U=\begin{bmatrix}u\\v\end{bmatrix},\quad
 \bm A=\begin{bmatrix}a&b/2\\0&1\end{bmatrix},\quad
 \bm B=\begin{bmatrix}b/2&c\\-1&0\end{bmatrix},\quad
 \bm S=\begin{bmatrix}-du-ev-f\\0\end{bmatrix}. \tag{1.16.30}
\end{equation}
第二个方程 $v_x-u_y=0$ 是新变量定义的相容条件。特征方程为
\begin{equation}
 a\lambda^2-b\lambda+c=0, \tag{1.16.33}
\end{equation}
故 $b^2-4ac>0$ 时为双曲型，$b^2-4ac<0$ 时为椭圆型，等于零时为抛物型。抛物情形仅有一条退化特征，必须结合低阶项才能描述解的行为；例如 $u_y=u_{xx}$ 中的一阶项 $u_y$ 正是时间演化项。

若 $\bm A^{-1}$ 存在，可写成 $\bm U_x+\bm A^{-1}\bm B\bm U_y=0$。若所有特征都朝 $x$ 增大方向传播，$x$ 就可视作时间型方向，$y$ 为类空间方向；二维稳态 Euler 方程是重要例子。给稳态系统加入时间导数后，整个系统还能关于时间变为双曲型：实空间特征对应定向传播的波，而复空间特征对应各向传播。稳态系统的特征结构还可用于把方程分解为双曲与椭圆子系统。

\section{标量微分方程的分类}

对二阶方程
\begin{equation}
 a\phi_{xx}+b\phi_{xy}+c\phi_{yy}+d\phi_x+e\phi_y+f=0, \tag{1.17.1}
\end{equation}
只需考察判别式：
\begin{equation}
 b^2-4ac\begin{cases}<0,&\text{椭圆型},\\=0,&\text{抛物型},\\>0,&\text{双曲型}.\end{cases} \tag{1.17.2}
\end{equation}
在求解之前先判断方程类型非常重要，因为相应数值策略可能完全不同。

典型椭圆方程是 Laplace 方程 $u_{xx}+u_{yy}=0$。它要求在整个边界上给条件：指定解的 Dirichlet 条件、指定法向导数的 Neumann 条件，或指定二者线性组合的 Robin 条件。内部解光滑，最大值和最小值只会出现在边界，这就是最大值原理。

典型抛物方程是热方程 $u_y-u_{xx}=0$，其中 $y$ 类似时间。$y=0$ 上的初始温度分布随 $y$ 增大沿 $x$ 扩散；还需在左右边界给定温度。即使初值不光滑，扩散也会随时间把它抹平。

典型双曲方程是波动方程 $u_{xx}-u_{yy}=0$，仍可把 $y$ 看成时间。它描述沿 $x$ 传播的波，要求在 $y=0$ 给初始数据；每一点只依赖初始数据的一部分。若采用周期域，则无需额外边界条件，成为纯初值问题。解可以不光滑甚至间断；这时经典微分方程不再逐点成立，但相应积分形式仍可定义弱解。

决定方程类型的仅是最高阶导数项，即方程的\textbf{主部}。即使源项含有很复杂的 $\phi_x\phi_y$ 或 $\phi_x^2$，仍只需考察 $a\phi_{xx}+b\phi_{xy}+c\phi_{yy}$，就能由判别式确定类型并选择适当算法。

\section{伪瞬态方程}

CFD 中常为了计算优势或便利而改写偏微分方程，甚至改变其类型。例如椭圆方程
\begin{equation}u_{xx}+u_{yy}=0\end{equation}
可以作为伪时间 $\tau$ 上的抛物方程
\begin{equation}u_\tau=u_{xx}+u_{yy} \tag{1.18.2}\end{equation}
推进，其稳态解就是原椭圆问题的解。这称为\textbf{伪瞬态法}或延拓法，使得为初值问题开发的时间推进方法可以直接用于稳态问题。

甚至可以引入 $u_{\tau\tau}$，得到双曲波动方程 $u_{\tau\tau}=u_{xx}+u_{yy}$，再趋近稳态。更实用的一阶双曲系统是
\begin{equation}
 u_\tau=\nu(p_x+q_y),\qquad p_\tau=\frac{u_x-p}{T_r},\qquad q_\tau=\frac{u_y-q}{T_r}. \tag{1.18.4}
\end{equation}
无论正参数 $T_r$ 取何值，稳态时它都退化为原 Laplace 方程，因此可用双曲系统的算法求解椭圆问题。

对二维双曲守恒律 $\bm U_t+\bm F_x+\bm G_y=0$，如果只关心稳态，$t$ 就可视为伪时间。还可用正定矩阵 $\bm P$ 改变瞬态波速：
\begin{equation}
 \bm U_\tau+\bm P(\bm F_x+\bm G_y)=0. \tag{1.18.7}
\end{equation}
只要 $\bm P$ 非奇异，稳态方程不变；却可以设计 $\bm P$ 来减小最快与最慢波速之比，加速所有瞬态波衰减。这就是局部预处理技术。瞬态解可能不具物理意义，但会更快收敛到原方程的正确稳态解。它还能改善低 Mach 数极限的精度，并可把控制方程分裂为双曲和椭圆部分，以便优化离散与多重网格收敛。不可压 Euler、Navier--Stokes 方程的人工可压缩形式也可理解为一种伪瞬态法。

伪瞬态思想也常用于非线性方程组的 Newton 求解：Newton 法在初值不够接近解时可能很慢，早期加入逐步增大的伪时间步可扩大收敛域。双时间方法则在每个物理时间步内，用伪瞬态迭代求解隐式时间离散产生的非线性系统。因此，一个快速稳态求解器也是隐式非定常算法的重要构件。

\section{算子分裂}

CFD 中常把偏微分方程逐项求解。考虑
\begin{equation}
 u_t=(\mathcal A+\mathcal B)u, \tag{1.19.1}
\end{equation}
其中 $\mathcal A,\mathcal B$ 为不同的线性微分算子，例如平流与扩散算子，或 $x$ 与 $y$ 方向导数算子。线性使得 $u_{tt}=(\mathcal A+\mathcal B)^2u$，更高时间导数同理，因此精确一步演化为
\begin{equation}
 u(t+\Delta t)=\exp[(\mathcal A+\mathcal B)\Delta t]u(t). \tag{1.19.6}
\end{equation}

若先解 $u_t^\star=\mathcal A u$，再解 $u_t=\mathcal B u^\star$，就得到算子分裂或分步法；这种顺序分裂通常称为 Godunov 分裂。它允许针对每个物理项选择专门算法，也使代码更简单，精度和稳定性可在各步分别控制。例如平流步可用显式高阶上风格式，扩散步可用隐式格式以容许较大时间步。

但分裂并不自动等价于原方程。两步展开给出
\begin{equation}
 u(t+\Delta t)=\left[1+\Delta t(\mathcal A+\mathcal B)
 +\frac{\Delta t^2}{2}(\mathcal A^2+2\mathcal B\mathcal A+\mathcal B^2)+O(\Delta t^3)\right]u(t). \tag{1.19.11}
\end{equation}
而精确展开的二阶项为 $(\mathcal A+\mathcal B)^2=\mathcal A^2+\mathcal A\mathcal B+\mathcal B\mathcal A+\mathcal B^2$。只有当
\begin{equation}\mathcal B\mathcal A=\mathcal A\mathcal B \tag{1.19.13}\end{equation}
即两算子可交换时，Godunov 分裂才无这一误差；否则整体时间精度至多一阶，无论各子步本身多高阶。

Strang 分裂通过对称顺序恢复二阶精度：
\begin{equation}
 \mathcal A(\Delta t/2)\;\longrightarrow\;\mathcal B(\Delta t)\;\longrightarrow\;\mathcal A(\Delta t/2). \tag{1.19.14--16}
\end{equation}
三段演化算子的乘积在 $O(\Delta t^2)$ 前与精确解一致。若以 $\mathcal A_{\Delta t/2}$、$\mathcal B_{\Delta t}$ 表示相应子问题的演化算子，一步为
\begin{equation}
 u(t+\Delta t)=\mathcal A_{\Delta t/2}\mathcal B_{\Delta t}\mathcal A_{\Delta t/2}u(t). \tag{1.19.18}
\end{equation}
相邻时间步之间的两个半步 $\mathcal A$ 可合并为一个整步，因此半步只在计算开始和真正需要输出解的时刻出现，额外代价可以降低。更高阶分裂也可以构造。使用算子分裂时仍须警惕不同物理算子的强耦合，以及稳态下分裂迭代可能收敛到错误解的情形，例如 $(\mathcal A+\mathcal B)u=0$ 且 $\mathcal A=-\mathcal B$。

\section{旋转不变性}

对守恒律
\begin{equation}
 \bm U_t+\bm F_x+\bm G_y+\bm H_z=0, \tag{1.20.1}
\end{equation}
若对任意非零方向 $\bm n=[n_x,n_y,n_z]^t$，法向通量满足
\begin{equation}
 \bm F_n=n_x\bm F+n_y\bm G+n_z\bm H=\bm T^{-1}\bm F(\bm T\bm U), \tag{1.20.2}
\end{equation}
则称该守恒律具有旋转不变性。Euler 方程具有此性质，其旋转矩阵由三个正交单位矢量 $\bm n,\bm\ell,\bm m$ 构造：密度和总能分量不变，三个动量分量旋转到法向和两个切向，且 $\bm T^{-1}=\bm T^t$。

这一性质意味着只实现 $x$ 方向通量函数 $\bm F$ 就够了。选择相应的置换旋转矩阵即可得到 $\bm G$ 和 $\bm H$；更一般地，在有限体积程序中，把单元面局部法向旋转到 $x$ 方向，用同一个界面通量子程序计算，再旋转回去即可。它既简化实现，也确保不同方向的通量处理一致。

\section{双曲方程的平面波}

依据 Huygens 原理，点扰动产生的波前是同时通过该点的全部平面波面的包络。若二维平面波沿极角 $\theta$ 的速度为 $\lambda(\theta)$，则波前参数式为
\begin{equation}
 \begin{bmatrix}X(\theta)\\Y(\theta)\end{bmatrix}=
 \begin{bmatrix}\cos\theta&-\sin\theta\\\sin\theta&\cos\theta\end{bmatrix}
 \begin{bmatrix}\lambda(\theta)\\\lambda'(\theta)\end{bmatrix}. \tag{1.21.1}
\end{equation}

\subsection*{线性平流}
方程 $u_t+a u_x+b u_y=0$ 的平面波速度为 $\lambda(\theta)=a\cos\theta+b\sin\theta$。代入包络公式得到 $(X,Y)=(a,b)$，说明点扰动只沿固定矢量 $(a,b)$ 传播。

\subsection*{各向同性平流}
若平流速度为 $R(\cos\theta,\sin\theta)$，则 $\lambda=R$、$\lambda'=0$，波前为
\begin{equation}X=R\cos\theta,\qquad Y=R\sin\theta,\qquad X^2+Y^2=R^2. \tag{1.21.7--8}\end{equation}
即扰动以相同速度向所有方向传播。

\subsection*{方程组}
系统 $\bm U_t+\bm A\bm U_x+\bm B\bm U_y=0$ 的平面波速度，是 $\bm A_n=\bm A\cos\theta+\bm B\sin\theta$ 的特征值。二维 Euler 方程有 $q_n-a,q_n,q_n,q_n+a$，其中 $q_n=(u,v)\cdot(\cos\theta,\sin\theta)$，$a$ 为声速。两支 $q_n$ 波随流以速度 $(u,v)$ 平移；两支声学波的包络满足
\begin{equation}
 X=u\pm a\cos\theta,\qquad Y=v\pm a\sin\theta,\qquad
 (X-u)^2+(Y-v)^2=a^2. \tag{1.21.13--14}
\end{equation}
即以流速点 $(u,v)$ 为圆心、声速 $a$ 为半径的圆。更复杂系统（例如磁流体力学）的波前会具有更有趣的形状。

\section{Rankine--Hugoniot 关系}

设不连续面（激波）的单位法向为 $\bm n$，法向传播速度为 $V_D$。把守恒律积分形式用于跨越不连续面的无限小控制体，得到
\begin{equation}
 [\bm F_n]=V_D[\bm U], \tag{1.22.1}
\end{equation}
方括号表示跨越不连续面的跳跃。若跳跃很小，则
\begin{equation}
 d\bm F_n=V_Dd\bm U,\qquad
 \frac{\partial\bm F_n}{\partial\bm U}d\bm U=V_Dd\bm U. \tag{1.22.2--3}
\end{equation}
因此，小振幅极限下激波速度是通量 Jacobian 的特征值，解跳跃沿相应特征矢量。线性方程中这一点对任意跳跃都成立；非线性方程中只对小跳跃成立。大振幅激波的速度\textbf{不是}局部 Jacobian 的特征值。

一维 Euler 方程中，一支声学激波的速度为
\begin{equation}
 V_D=u_L+a_L\sqrt{1+\frac{\gamma+1}{2\gamma}\frac{p_R-p_L}{p_L}}, \tag{1.22.4}
\end{equation}
其中通常 $\gamma=1.4$，$u_L,p_L,a_L$ 是激波一侧的速度、压力和声速，$p_R$ 是另一侧压力。当 $p_R-p_L\to0$ 时，$V_D\to u_L+a_L$，这正是相应声学场的 Jacobian 特征值。另一支激波速度
\begin{equation}
 V_D'=u_L-a_L\sqrt{1+\frac{\gamma+1}{2\gamma}\frac{p_R-p_L}{p_L}} \tag{1.22.6}
\end{equation}
在小跳跃极限趋于 $u_L-a_L$。因此可借极限识别激波属于哪一个特征场。

\section{线性退化场与真正非线性场}

对双曲系统 $\bm U_t+\bm A\bm U_x=0$，第 $k$ 个特征值 $\lambda_k(\bm U)$ 及其右特征矢量 $\bm r_k(\bm U)$ 定义第 $k$ 个特征场。若对所有状态
\begin{equation}
 \frac{\partial\lambda_k}{\partial\bm U}\cdot\bm r_k=0, \tag{1.23.2}
\end{equation}
则该场是\textbf{线性退化}的。线性双曲系统的特征值均为常数，所有场都线性退化。接触间断就是线性退化场中的间断，其两侧特征线彼此平行，如同标量线性平流。

若对所有状态
\begin{equation}
 \frac{\partial\lambda_k}{\partial\bm U}\cdot\bm r_k\ne0, \tag{1.23.3}
\end{equation}
则称该场\textbf{真正非线性}。此时特征速度会随波内状态改变，可产生特征线会聚的非线性激波和特征线发散的非线性膨胀波。具有非零特征速度的非线性标量守恒律属于这种情形，其通量称为凸通量。

一维 Euler 方程以原始变量 $\bm U=[\rho,u,p]^t$ 表示时，
\begin{equation}
 \lambda_1=u-a,\qquad\lambda_2=u,\qquad\lambda_3=u+a,
 \quad a=\sqrt{\gamma p/\rho}, \tag{1.23.4}
\end{equation}
相应右特征矢量可取
\begin{equation}
 \bm r_1=\begin{bmatrix}-\rho/a\\1\\-\rho a\end{bmatrix},\quad
 \bm r_2=\begin{bmatrix}1\\0\\0\end{bmatrix},\quad
 \bm r_3=\begin{bmatrix}\rho/a\\1\\\rho a\end{bmatrix}. \tag{1.23.5}
\end{equation}
直接计算得到
\begin{equation}
 \nabla_{\bm U}\lambda_1\cdot\bm r_1=\frac{\gamma+1}{2}\ne0,\qquad
 \nabla_{\bm U}\lambda_2\cdot\bm r_2=0,\qquad
 \nabla_{\bm U}\lambda_3\cdot\bm r_3=\frac{\gamma+1}{2}\ne0. \tag{1.23.6--8}
\end{equation}
因此两支声学场 $\lambda_1,\lambda_3$ 真正非线性，而熵波场 $\lambda_2$ 线性退化。

\chapter{模型方程}
\markboth{第 2 章\quad 模型方程}{第 2 章\quad 模型方程}

\section{线性平流}

线性平流方程可写成
\begin{equation}
 u_t+(\operatorname{grad}u)\cdot\bm\lambda=0, \tag{2.1.1}
\end{equation}
或守恒形式
\begin{equation}
 u_t+\operatorname{div}(u\bm\lambda)=0, \tag{2.1.2}
\end{equation}
也可简写为 $Du/Dt=0$。这里 $\bm\lambda$ 是常平流速度，解 $u$ 以该速度传播。若 $\bm\lambda=(a,b,c)$，则笛卡尔坐标下
\begin{align}
 &\text{一维：}&u_t+a u_x&=0, \tag{2.1.4}\\
 &\text{二维：}&u_t+a u_x+b u_y&=0, \tag{2.1.5}\\
 &\text{三维：}&u_t+a u_x+b u_y+c u_z&=0. \tag{2.1.6}
\end{align}

线性平流是一般拟线性系统的模型，例如
\begin{equation}
 \bm U_t+\bm A\bm U_x=0, \tag{2.1.7}
\end{equation}
其中 $\bm A$ 具有实特征值；气体动力学 Euler 方程属于这一类。开发求解方程组的算法时，先研究保留关键特征的简单模型方程是明智的。由于系统可以对角化为一组线性平流方程，为标量平流开发的方法可逐个用于各特征分量，再推广到完整系统。

平流方程的积分形式
\begin{equation}
 \frac{d}{dt}\int_Vu\,dV+\oint_Su\bm\lambda\cdot\bm n\,dS=0 \tag{2.1.8}
\end{equation}
允许不连续解：任何连续或不连续的初始分布都会以 $\bm\lambda$ 原样平移。这使它成为开发激波流数值方法的良好模型；许多稳定、准确计算间断的技术，后来都成功推广到 Euler 方程。

一维线性平流与波动方程关系密切。一维波动方程 $u_{tt}-a^2u_{xx}=0$ 可因式分解为 $(\partial_t-a\partial_x)(\partial_tu+a\partial_xu)=0$，所以有两支一般解，分别以 $a$ 和 $-a$ 传播；平流方程只控制其中一支。二维波动方程
\begin{equation}u_{tt}-a^2(u_{xx}+u_{yy})=0 \tag{2.1.11}\end{equation}
则在所有方向都有速度 $a$ 的平面波。常系数二维平流只沿固定 $(a,b)$ 传播；若要描述各向同性传播，速度必须随空间方向变化为 $a(\cos\theta,\sin\theta)$。

若仍希望保留常系数，可用系统表示各向同性波：
\begin{equation}
 u_t=a(p_x+q_y),\qquad p_t=a u_x,\qquad q_t=a u_y. \tag{2.1.13--15}
\end{equation}
直接微分即可证明它等价于二维波动方程，但额外引入了约束 $(p_y-q_x)_t=0$，即任意初始“涡量”都被永久保持。这类模型曾用于发展涡量保持格式；三维情形类似，不过涡量有三个分量，约束也有三个。

有人用 convection（对流）代替 advection（平流）。作者因习惯而偏爱后者，并不打算争论哪一个更好。

\section{圆周平流}

即使 $\bm\lambda=(a,b,c)$ 依赖空间，平流方程仍可能是线性的。例如二维方程
\begin{equation}
 u_t+a u_x+b u_y=0,\qquad (a,b)=(y,-x) \tag{2.2.1--2}
\end{equation}
表示绕原点顺时针旋转的速度场，称为圆周平流。三维可取 $(a,b,c)=(y,-x,0)$，表示绕 $z$ 轴旋转。取 $(0,z,-y)$ 或 $(z,0,-x)$ 可分别绕 $x$ 或 $y$ 轴旋转；取 $(y,-x,V)$ 则在绕 $z$ 轴旋转的同时，以常速度 $V$ 沿 $z$ 方向推进，形成更具三维性的螺旋运动。由于这些速度场散度为零，仍可写成守恒形式 $u_t+\operatorname{div}(u\bm\lambda)=0$。

\section{扩散方程}

扩散方程为
\begin{equation}
 u_t=\operatorname{div}(\nu\operatorname{grad}u), \tag{2.3.1}
\end{equation}
其中 $\nu>0$ 为扩散系数。只要 $\nu$ 不依赖解本身，方程就是线性的。笛卡尔坐标下
\begin{align}
 &\text{一维：}&u_t&=(\nu u_x)_x, \tag{2.3.2}\\
 &\text{二维：}&u_t&=(\nu u_x)_x+(\nu u_y)_y, \tag{2.3.3}\\
 &\text{三维：}&u_t&=(\nu u_x)_x+(\nu u_y)_y+(\nu u_z)_z. \tag{2.3.4}
\end{align}
扩散方程关于时间为抛物型，需要初值和边界条件。其解通常很光滑；即使初值不光滑，二阶扩散项的平滑作用也会逐渐把它抹平。若边界条件和源项不随时间变化，解最终趋于稳态，此时等价于 Laplace 方程。

积分形式为
\begin{equation}
 \frac{d}{dt}\int_Vu\,dV+\oint_S\nu\operatorname{grad}u\cdot\bm n\,dS=0,
 \qquad
 \frac{d}{dt}\int_Vu\,dV+\oint_S\nu\frac{\partial u}{\partial n}\,dS=0. \tag{2.3.5--6}
\end{equation}
这是构造有限体积扩散格式的自然起点。在规则网格上发展二阶扩散格式相对容易，因为只处理光滑解，而且界面法向常与相邻两个数据点的连线重合，法向梯度可用简单差分计算。困难更多在效率：显式格式的时间步受严重限制，通常更偏好隐式格式。这些平滑性和效率特征，也正是 Navier--Stokes 方程黏性部分的特征，因此扩散方程是黏性项的典型模型。

非规则网格上，扩散离散并不一定容易；若法向梯度评价不慎，精度和迭代收敛都会显著恶化。压缩 Navier--Stokes 黏性项还不只包含法向梯度，也包含沿控制体边界的切向梯度。有些薄层近似忽略切向梯度，这类格式往往稳健，但因不相容而不能保证正确精度；不过它们仍可用于构造压缩 Navier--Stokes 方程的隐式近似算子。

\section{扩散方程：混合形式}

扩散方程常写成一阶系统
\begin{equation}
 u_t=\nu\operatorname{div}\bm p,\qquad \bm p=\operatorname{grad}u, \tag{2.4.1--2}
\end{equation}
其中 $\bm p=(p,q,r)$ 代表梯度变量。笛卡尔坐标下
\begin{align}
 &\text{一维：}&u_t&=\nu p_x,&p&=u_x, \tag{2.4.3}\\
 &\text{二维：}&u_t&=\nu(p_x+q_y),&p&=u_x,&q&=u_y, \tag{2.4.4}\\
 &\text{三维：}&u_t&=\nu(p_x+q_y+r_z),&p&=u_x,&q&=u_y,&r&=u_z. \tag{2.4.5}
\end{align}
这在有限元领域称为混合形式，也广泛用于高阶方法。表面上它只含一阶导数，在非结构网格上通常比直接离散二阶导数方便，由此发展出了许多成功的扩散格式。

不过，这种改写并未真正消除二阶导数：该系统与原扩散方程完全等价，显式时间步仍受 $O(h^2)$ 限制。它关于时间仍是抛物型，因此没有改变扩散物理；降阶只是为了离散方便。

\section{扩散方程：一阶双曲系统}

与扩散方程相关的一个有趣系统是
\begin{align}
 &\text{一维：}&u_t&=\nu p_x,&p_t&=(u_x-p)/T_r, \tag{2.5.1}\\
 &\text{二维：}&u_t&=\nu(p_x+q_y),&p_t&=(u_x-p)/T_r,&q_t&=(u_y-q)/T_r, \tag{2.5.2}\\
 &\text{三维：}&u_t&=\nu(p_x+q_y+r_z),&p_t&=(u_x-p)/T_r,&q_t&=(u_y-q)/T_r,&r_t&=(u_z-r)/T_r. \tag{2.5.3}
\end{align}
在 $T_r\to0$ 极限下，它与非稳态扩散方程等价，常称双曲热方程。令 $\bm U=[u,p,q,r]^t$，系统可写成 $\bm U_t+\bm F_x+\bm G_y+\bm H_z=\bm S$；通量、源项以及三个通量 Jacobian 的完整矩阵如下：
\sourcecrop{60}{46}{52}{566}{300}{0.42}
\eqnote{$T_r$ 是松弛时间；$p,q,r$ 分别逼近 $u_x,u_y,u_z$。源项以时间尺度 $T_r$ 把梯度变量松弛到真实梯度。原书 (2.5.4) 中 $\partial_yH$ 应按三维上下文理解为 $\partial_zH$。}

沿任意单位方向 $\bm n=[n_x,n_y,n_z]^t$ 的投影 Jacobian 为
\begin{equation}
 \bm A_n=
 \begin{bmatrix}
 0&-\nu n_x&-\nu n_y&-\nu n_z\\
 -n_x/T_r&0&0&0\\-n_y/T_r&0&0&0\\-n_z/T_r&0&0&0
 \end{bmatrix}, \tag{2.5.7}
\end{equation}
其特征值为
\begin{equation}
 \lambda_{1,2}=\mp\sqrt{\nu/T_r},\qquad \lambda_{3,4}=0. \tag{2.5.8}
\end{equation}
特征值全为实数，特征矢量线性无关，因此系统为双曲型。波速与方向 $\bm n$ 无关，说明传播各向同性。两个零特征值对应“不相容性阻尼模态”，作用于理论上应为零的 $q_x-p_y$、$r_x-p_z$。系统只在 $O(T_r)$ 的短时间尺度上显现双曲行为；更长时间后，源项把波阻尼掉，系统表现为扩散。用它代替扩散方程可以避免二阶导数并复用成熟的双曲格式，但很小的 $T_r$ 会使源项刚性，数值求解仍然困难。

若只关心稳态，情况大不相同：$T_r$ 可取任意正值。时间导数消失后，梯度方程强制 $p=u_x$、$q=u_y$、$r=u_z$，主方程于是分别化为
\begin{equation}
 \nu u_{xx}=0,\qquad \nu(u_{xx}+u_{yy})=0,\qquad
 \nu(u_{xx}+u_{yy}+u_{zz})=0. \tag{2.5.9--11}
\end{equation}
即任意 $T_r>0$ 时都与稳态扩散方程（Laplace 方程）等价。为区别于必须取很小 $T_r$ 才逼近非稳态扩散的双曲热方程，这种专为稳态计算设计的系统称为\textbf{双曲扩散系统}。

等价性也可由能量估计看出。定义
\begin{equation}
 E=\frac{u^2+L_r^2(p^2+q^2+r^2)}{2}, \tag{2.5.12}
\end{equation}
则系统给出能量方程
\begin{equation}
 E_t+(f^E)_x+(g^E)_y+(h^E)_z=-\nu(p^2+q^2+r^2),
 \quad (f^E,g^E,h^E)=-\nu u(p,q,r). \tag{2.5.13--14}
\end{equation}
在区域 $\Omega$ 上积分得
\begin{equation}
 \frac{d}{dt}\int_\Omega E\,dV=-\oint_{\partial\Omega}\bm f^E\cdot\bm n\,dA
 -\nu\int_\Omega(p^2+q^2+r^2)\,dV. \tag{2.5.15}
\end{equation}
总能量由解梯度的大小耗散，只能通过边界通量改变。稳态时 $(p,q,r)=\nabla u$，能量估计退化为
\begin{equation}
 \int_\Omega\nabla u\cdot\nabla u\,dV-
 \oint_{\partial\Omega}u\frac{\partial u}{\partial n}\,dA=0, \tag{2.5.16}
\end{equation}
这正是 Laplace 方程的经典能量估计。

稳态双曲扩散方法不再受刚性松弛限制，可以选择 $T_r$ 使系统始终强双曲并优化收敛。它允许 $O(h)$ 而不是 $O(h^2)$ 的显式时间步；可直接使用上风格式等双曲系统技术；解 $u$ 与梯度变量 $(p,q,r)$ 可获得相同阶精度；Neumann 条件也可通过梯度变量像 Dirichlet 条件一样直接施加。它在稳态下虽与混合形式相同，但离散方法和数值性质截然不同。

\section{平流--扩散方程}

平流--扩散方程为
\begin{equation}
 u_t+(\operatorname{grad}u)\cdot\bm\lambda=
 \operatorname{div}(\nu\operatorname{grad}u), \tag{2.6.1}
\end{equation}
守恒形式为
\begin{equation}
 u_t+\operatorname{div}(u\bm\lambda-\nu\operatorname{grad}u)=0. \tag{2.6.2}
\end{equation}
若 $\bm\lambda=(a,b,c)$，则一至三维分别为
\begin{align}
 u_t+a u_x&=(\nu u_x)_x, \tag{2.6.3}\\
 u_t+a u_x+b u_y&=(\nu u_x)_x+(\nu u_y)_y, \tag{2.6.4}\\
 u_t+a u_x+b u_y+c u_z&=(\nu u_x)_x+(\nu u_y)_y+(\nu u_z)_z. \tag{2.6.5}
\end{align}
把各方向的平流和扩散合并为总通量，就得到 $(au-\nu u_x)_x$ 等项。它是拟线性系统
\begin{equation}
 \bm U_t+\bm A\bm U_x=\bm D\bm U_{xx} \tag{2.6.9}
\end{equation}
的标量模型，其中 $\bm A$ 有实特征值，$\bm D$ 为正系数矩阵；Navier--Stokes 方程的拟线性形式属于这一类。

它虽只是平流与扩散之和，数值上却不能草率地把任意平流格式与任意扩散格式相加。在两种效应同等重要的中间区域，不相容离散会使某些高阶格式丧失形式精度。保持精度的重要原则，是统一处理平流和扩散；守恒总通量形式尤其适合这一目的，其积分形式为
\begin{equation}
 \frac{d}{dt}\int_Vu\,dV+\oint_S\left(u\bm\lambda\cdot\bm n-\nu\frac{\partial u}{\partial n}\right)dS=0. \tag{2.6.10}
\end{equation}

另一个困难是方程性质会随参数剧烈变化。考虑 $x\in(0,L)$ 上
\begin{equation}
 u_t+a u_x=\nu u_{xx},\qquad a>0. \tag{2.6.11}
\end{equation}
当 $\nu>0$ 时是抛物问题，必须在 $x=0,L$ 两端给值；非平凡解要求两值不同。当 $\nu=0$ 时却变成右行平流，只能在左端给入流值 $u(0)$，右端值必须由从左传播而来的波决定。一般 $\nu>0$ 时，平流试图把左端值带到右边，而扩散在靠近右边界处调和这一值与规定的 $u(L)$。平流占优时，调和只发生在狭窄的右边界层内。

衡量两种效应的无量纲参数是 Peclet 数
\begin{equation}
 \mathrm{Pe}=\frac{aL}{\nu}. \tag{2.6.12}
\end{equation}
$\mathrm{Pe}\to0$ 时扩散占优，过渡层变宽；$\mathrm{Pe}\to\infty$ 时平流占优，过渡层极窄，宏观上近似纯平流。Navier--Stokes 方程中的对应参数是 Reynolds 数，物体表面的窄层就是边界层。若要计算黏性阻力或表面传热，就不能丢弃扩散项，必须准确、高效地解析这一薄层。

所以，平流--扩散方程虽然只是简单线性标量方程，却包含实际 Navier--Stokes 计算会遇到的核心困难，既适合说明这些困难，也适合在简化环境中发展新方法。

\section{平流--反应方程}

平流--反应方程为
\begin{equation}
 u_t+(\operatorname{grad}u)\cdot\bm\lambda=-ku. \tag{2.7.1}
\end{equation}
笛卡尔坐标中的一、二、三维形式，分别在平流方程左端加入同一衰减源项 $-ku$。它描述物质以常速度 $\bm\lambda$ 随流输运，同时以速率 $k$ 衰减，是含源项守恒律 $u_t+\operatorname{div}\bm f=s$ 的典型模型，其中 $\bm f=u\bm\lambda$、$s=-ku$。当 $k$ 很大时，源项产生刚性，数值求解会非常困难。

\section{Poisson 与 Laplace 方程}

Poisson 方程为
\begin{equation}
 \operatorname{div}\operatorname{grad}u=f, \tag{2.8.1}
\end{equation}
$f=0$ 时称为 Laplace 方程。笛卡尔坐标下
\begin{align}
 u_{xx}&=f(x), &&\text{一维},\\
 u_{xx}+u_{yy}&=f(x,y), &&\text{二维},\\
 u_{xx}+u_{yy}+u_{zz}&=f(x,y,z), &&\text{三维}. \tag{2.8.2--4}
\end{align}
Laplace 解非常光滑，最大值和最小值只能出现在边界。Poisson 解的光滑性则取决于源项。两者都含容易在规则网格上用中心差分离散的 Laplace 算子。直接法往往代价很高，普通迭代也可能需要很多步，但多重网格对这类方程非常有效；非规则网格或高阶精度要求则会显著增加难度。

\section{Cauchy--Riemann 方程}

\subsection{速度分量}

二维不可压无旋流中，速度分量 $(u,v)$ 满足
\begin{equation}
 u_x+v_y=0\quad\text{（不可压条件）},\qquad
 v_x-u_y=0\quad\text{（无旋条件）}. \tag{2.9.1--2}
\end{equation}
写成 $\bm A\bm U_x+\bm B\bm U_y=0$，其中
\begin{equation}
 \bm U=\begin{bmatrix}u\\v\end{bmatrix},\quad
 \bm A=\begin{bmatrix}1&0\\0&-1\end{bmatrix},\quad
 \bm B=\begin{bmatrix}0&1\\1&0\end{bmatrix}. \tag{2.9.4}
\end{equation}
$\bm A^{-1}\bm B$ 的特征值为 $\pm i$，相应特征矢量可取 $[1,\pm i]^t$，故该系统是椭圆型。若 $u,v$ 二阶可微，它们各自满足 Laplace 方程。并且
\begin{equation}
 \operatorname{grad}u\cdot\operatorname{grad}v=u_xv_x+u_yv_y=(-v_y)v_x+v_xv_y=0, \tag{2.9.8}
\end{equation}
所以 $u$ 与 $v$ 的等值线彼此正交。

Cauchy--Riemann 系统与一对 Laplace 方程密切相关，很多离散方法实际会隐式求解相应 Laplace 方程。但反过来，Laplace 方程的解未必满足 Cauchy--Riemann 约束。例如固体转动 $(u,v)=r(\sin\theta,-\cos\theta)$ 满足分量 Laplace 方程，却有非零涡量。因此，若目标确实是 Cauchy--Riemann 解，就必须使用真正保持这两个一阶约束的算法。

该系统是 Euler 方程椭圆子系统（亚声速声学子系统）的模型。更一般地，任何二维偏微分方程组都可分解成若干平流方程和/或若干 Cauchy--Riemann 系统。

\subsection{流函数与速度势}

流函数 $\psi$ 和速度势 $\phi$ 定义为
\begin{equation}
 u=\phi_x=\psi_y,\qquad v=\phi_y=-\psi_x. \tag{2.9.9--10}
\end{equation}
于是
\begin{equation}
 \psi_x+\phi_y=0,\qquad \phi_x-\psi_y=0 \tag{2.9.11--12}
\end{equation}
正是 Cauchy--Riemann 方程；若二阶可微，$\phi$ 与 $\psi$ 各自满足 Laplace 方程。流函数的固壁边界条件非常简单：沿不可穿透固壁取常数即可。速度势在多连通域中则可能不连续，例如有升力翼型绕流。流线（$\psi$ 等值线）与等势线（$\phi$ 等值线）在交点处正交。

\subsection{涡量与压力}

在黏性项与压力梯度平衡的 Stokes 流中，压力 $p$ 与涡量 $\omega=v_x-u_y$ 满足
\begin{equation}
 p_x+\mu\omega_y=0,\qquad p_y-\mu\omega_x=0, \tag{2.9.15--16}
\end{equation}
其中 $\mu$ 为动力黏度。因此可以直接应用复变函数理论，像势流一样构造许多黏性流精确解。

\subsection{含时间系统}

给 Cauchy--Riemann 方程加入时间导数：
\begin{equation}
 u_t+u_x+v_y=0,\qquad v_t-v_x+u_y=0. \tag{2.9.17--18}
\end{equation}
它不再是原 Cauchy--Riemann 系统，但稳态时完全等价，因此可以时间推进到不再变化，从而求得稳态解。时间项的符号经过选择，使演化稳定。任意方向投影 Jacobian 的特征值为 $\pm1$，特征矢量完备，所以系统关于时间是双曲型，尽管其空间稳态问题为椭圆型。

若改为给两个 Laplace 方程加入一阶时间项，就得到两条抛物方程；若想得到双曲方程，必须加入二阶时间导数，形成波动方程。无论抛物还是双曲，这些伪时间系统的稳态解都是 Laplace 解。作者更喜欢双曲形式，因为简单显式格式往往可以用较大时间步快速达到稳态。

\section{双调和方程}

双调和方程是四阶偏微分方程
\begin{equation}
 \operatorname{div}\operatorname{grad}(\operatorname{div}\operatorname{grad}u)=f,
 \qquad \Delta^2u=f. \tag{2.10.1--2}
\end{equation}
笛卡尔坐标下
\begin{align}
 u_{xxxx}&=f(x),\\
 u_{xxxx}+u_{yyyy}+2u_{xxyy}&=f(x,y),\\
 u_{xxxx}+u_{yyyy}+u_{zzzz}+2u_{xxyy}+2u_{yyzz}+2u_{zzxx}&=f(x,y,z). \tag{2.10.3--5}
\end{align}
它确实出现在 CFD 中：极低 Reynolds 数的 Stokes 方程，用流函数--涡量形式表示时，流函数满足双调和方程。典型边界条件为
\begin{equation}
 u=g_1(x,y,z),\qquad \frac{\partial u}{\partial n}=g_2(x,y,z). \tag{2.10.6}
\end{equation}
二维 Stokes 流中 $u$ 是流函数；固壁可取 $g_1=0$，表示壁面为流线；第二个条件指定切向速度，无滑移壁面可取 $g_2=0$。

规则网格上直接离散四阶导数并不困难，非规则网格上却不简单。一个有趣的降阶办法是引入新变量 $v$，写成两条 Laplace 型方程
\begin{equation}
 \operatorname{div}\operatorname{grad}u=v,\qquad
 \operatorname{div}\operatorname{grad}v=f. \tag{2.10.7--8}
\end{equation}
它们与 $u$ 的双调和方程等价，因而可以复用 Laplace 方程的成熟求解器，例如多重网格。困难在于新变量 $v$ 的边界条件可能未知；不过总的独立边界条件数仍是两个，Stokes 流中前述流函数条件通常已经足够，无需另外指定 $v$。

\section{Euler--Tricomi 方程}

Euler--Tricomi 方程
\begin{equation}
 y u_{xx}+u_{yy}+u_{zz}=0 \tag{2.11.1}
\end{equation}
常用于研究跨声速流数值方法。它的类型随 $y$ 改变：$y<0$ 时为双曲型，$y>0$ 时为椭圆型。跨声速流也有类似性质：局部流速超过声速时为双曲型，低于声速时为空间椭圆型。该模型在双曲区形成陡峭波，在椭圆区产生光滑解，因此常用于测试跨声速算法。处理混合性质可以直接采用中心--上风混合离散，也可以在伪时间上把方程改造成双曲系统，再用上风法推进；后者与稳态 Euler 计算中常见做法相似。

\section{松弛模型}

一维松弛模型为
\begin{equation}
 u_t+v_x=0,\qquad v_t+a^2u_x=(cu-v)/\tau, \tag{2.12.1--2}
\end{equation}
其中 $\tau>0$、$a^2>c^2$。当 $\tau\ll1$ 时，经 Chapman--Enskog 展开可知它表现为
\begin{equation}
 u_t+c u_x=\tau(a^2-c^2)u_{xx}. \tag{2.12.3}
\end{equation}
这是从 Boltzmann 方程导出的矩方程的简单模型：小松弛时间极限下，它们表现得像 Navier--Stokes 方程。

二维版本为
\begin{equation}
 u_t+v_x+w_y=0,\qquad
 v_t+a^2u_x=-(v-ru)/\tau,\qquad
 w_t+b^2u_y=-(w-su)/\tau, \tag{2.12.4--6}
\end{equation}
其中 $a^2>r^2$、$b^2>s^2$。当 $\tau\ll1$ 时，它逼近平流--扩散方程
\begin{equation}
 u_t+r u_x+s u_y=\tau\left[(a^2-r^2)u_{xx}-2rsu_{xy}+(b^2-s^2)u_{yy}\right]. \tag{2.12.8}
\end{equation}
直接求解一阶松弛系统有时比求解右端含混合导数的平流--扩散方程更有利，因为可避开混合导数离散的困难。

更一般地，非线性守恒律 $\bm U_t+\bm F(\bm U)_x=0$ 可通过新变量 $\bm V$ 改写为带源项的线性系统
\begin{equation}
 \bm U_t+\bm V_x=0,\qquad
 \bm V_t+\bm A\bm U_x=\frac{\bm F(\bm U)-\bm V}{\tau}. \tag{2.12.10--11}
\end{equation}
常矩阵 $\bm A$ 的特征值必须满足次特征条件，$\tau$ 必须足够小，才能保证与原守恒律等价。这样避免了通量 Jacobian 的非线性，却换来了刚性源项，所以困难并未消失，只是转移了形式。这类方法称为松弛方法。

\section{Burgers 方程}

无黏 Burgers 方程为
\begin{equation}
 u_t+u u_x=0, \tag{2.13.1}
\end{equation}
守恒形式是
\begin{equation}
 u_t+\left(\frac{u^2}{2}\right)_x=0. \tag{2.13.2}
\end{equation}
跨越左右状态 $u_L,u_R$ 的间断应用 Rankine--Hugoniot 关系，得到激波速度
\begin{equation}
 V_s=\frac{u_L+u_R}{2}. \tag{2.13.3}
\end{equation}
若把原方程乘以 $2u$，还可得到另一守恒形式
\begin{equation}
 (u^2)_t+\left(\frac23u^3\right)_x=0, \tag{2.13.4}
\end{equation}
它给出的激波速度却是
\begin{equation}
 V_s=\frac23\left(u_L+u_R-\frac{u_Lu_R}{u_L+u_R}\right). \tag{2.13.5}
\end{equation}
两个守恒形式具有相同的强解，却有不同弱解，这是因为守恒量不同：前者守恒 $u$，后者守恒 $u^2$。这个例子说明，必须选择物理上正确的守恒变量，才能得到有意义的间断解。

二维矢量版本为
\begin{equation}
 u_t+u u_x+v u_y=0,\qquad v_t+u v_x+v v_y=0, \tag{2.13.6--7}
\end{equation}
也常使用 $u_t+u u_x+u_y=0$；稳态时把 $y$ 看作类时间方向，它与一维方程本质相同。三维版本为速度三个分量都随自身速度场输运：
\begin{equation}
 \bm v_t+(\bm v\cdot\nabla)\bm v=0. \tag{2.13.9--11}
\end{equation}
顺便说明，英文名是 Burgers' equation，来自人名 Burgers，而不是 Burger。

\section{交通流方程}

交通流模型为
\begin{equation}
 \rho_t+(\rho v)_x=0, \tag{2.14.1}
\end{equation}
其中 $\rho$ 是道路上单位长度的车辆数，$v$ 是局部车速。一个合理闭合关系是
\begin{equation}
 \frac{v}{v_{\max}}=1-\frac{\rho}{\rho_{\max}}, \tag{2.14.2}
\end{equation}
即车辆少时速度高，拥挤时速度低。通量 $f=\rho v$ 的特征速度为
\begin{equation}
 \frac{\partial f}{\partial\rho}=v_{\max}\left(1-\frac{2\rho}{\rho_{\max}}\right), \tag{2.14.3}
\end{equation}
而 Rankine--Hugoniot 关系给出的激波速度为
\begin{equation}
 V_s=\frac{f_R-f_L}{\rho_R-\rho_L}
 =v_{\max}\left(1-\frac{\rho_R+\rho_L}{\rho_{\max}}\right). \tag{2.14.4}
\end{equation}
交通流方程需要一个密度--速度关系来闭合，这个关系可针对真实交通数据提出不同模型；但每换一种闭合关系，都必须重新推导相应的特征速度和激波速度。例如可尝试 $v/v_{\max}=1-\rho^2/\rho_{\max}^2$。

\section{黏性 Burgers 方程}

黏性 Burgers 方程为
\begin{equation}
 u_t+u u_x=\nu u_{xx},\qquad \nu>0. \tag{2.15.1}
\end{equation}
它是 Navier--Stokes 方程的简单模型。二维版本为
\begin{equation}
 u_t+u u_x+v u_y=\nu(u_{xx}+u_{yy}),\qquad
 v_t+u v_x+v v_y=\nu(v_{xx}+v_{yy}), \tag{2.15.2--3}
\end{equation}
实际上就是去掉压力梯度后的二维不可压 Navier--Stokes 动量方程。另一常用形式
\begin{equation}
 u_t+u u_x+u_y=\nu u_{xx} \tag{2.15.4}
\end{equation}
在稳态下把 $y$ 看作类时间变量，表现得像一维黏性 Burgers 方程。三维版本可紧凑写成
\begin{equation}
 \bm v_t+(\bm v\cdot\nabla)\bm v=\nu\Delta\bm v, \tag{2.15.5--7}
\end{equation}
即三个速度分量均由同一速度场输运并作各向同性扩散。黏性 Burgers 方程最有趣的性质之一，是其解可以由扩散方程的解构造；后文将通过 Cole--Hopf 型变换详细说明。


% ---- 合并分片 2 ----
\chapter{欧拉方程}

\section{热力学关系}

本书主要考虑服从理想气体定律（也称热状态方程）的气体：
\begin{equation}
  p=\rho RT,
\end{equation}
其中，$p$ 为压强，$\rho$ 为密度，$T$ 为温度，$R$ 为气体常数；海平面空气约有
$R=287\,\mathrm{N\,m/(kg\,K)}$。服从该关系的气体称为理想气体或热完全气体。
它表明，在同一点处，$p$、$\rho$、$T$ 三者中的任意两个即可确定气体状态，因此
它们都是状态变量。虽然流场变量会随时间变化，但该关系假定气体在每一时刻、每一
位置都处于局部热力学平衡。这个假设虽非在所有真实条件下都严格成立，却能很好地
覆盖从低速到超声速的广泛流动，是热力学在流体力学和 CFD 中极为有用的原因。

气体内能是分子平动、转动、振动等微观能量之和。以 $e$ 表示单位质量内能。一般地，
$e$ 由两个状态变量决定；对于理想气体则只依赖温度，$e=e(T)$。定义比焓
\begin{equation}
  h\equiv e+\frac{p}{\rho},
\end{equation}
由于 $p/\rho=RT$，理想气体的焓同样只依赖于温度，即 $h=h(T)$。

令比容 $v_s=1/\rho$，热力学第一定律可写成
\begin{equation}
  dq=de+p\,dv_s,
  \qquad
  dq=dh-\frac{1}{\rho}\,dp .
\end{equation}
这里 $de$、$dh$ 只取决于状态，而所加比热量 $dq$ 以及功取决于过程；上式用可逆的
压力功 $p\,dv_s$ 代替了过程功，因此未包含黏性力。这两式也适用于非理想气体。
定容比热与定压比热定义为
\begin{equation}
 c_v=\left(\frac{\partial q}{\partial T}\right)_{v_s},\qquad
 c_p=\left(\frac{\partial q}{\partial T}\right)_p,
\end{equation}
于是立刻得到
\begin{equation}
 c_v=\left(\frac{\partial e}{\partial T}\right)_{v_s},\qquad
 c_p=\left(\frac{\partial h}{\partial T}\right)_p.
\end{equation}
也就是说，$c_v$ 表示定容条件下内能对温度的变化率，$c_p$ 表示定压条件下焓对温度
的变化率。

若分子具有 $\alpha$ 个平动与转动自由度，动理学给出
\begin{equation}
 c_v=\frac{\alpha}{2}R,\qquad
 c_p=\left(1+\frac{\alpha}{2}\right)R,
 \qquad
 \gamma\equiv\frac{c_p}{c_v}=\frac{\alpha+2}{\alpha}.
\end{equation}
$\gamma$ 称为比热比或绝热指数。单原子气体 $\alpha=3$，故 $\gamma=5/3$；空气等双
原子气体 $\alpha=5$，故 $\gamma=7/5$。如果温度高到足以激发分子振动能，这些常数
会随温度变化。对理想气体，由第一定律和状态方程可得重要关系
\begin{equation}
 c_p-c_v=R,
 \qquad c_v=\frac{R}{\gamma-1},
 \qquad c_p=\frac{\gamma R}{\gamma-1}.
\end{equation}
若比热随温度变化，则 $e(T)=\int c_v\,dT+C_e$、$h(T)=\int c_p\,dT+C_h$。若比热为
常数，气体称为量热完全气体，并有
\begin{equation}
 e=c_v T=\frac{RT}{\gamma-1},\qquad
 h=c_p T=\frac{\gamma RT}{\gamma-1}.
\end{equation}
下文简称这种气体为完全气体。

定义比总能和比总焓
\begin{equation}
 E=e+\frac{\bm v\cdot\bm v}{2}
   =\frac{p}{(\gamma-1)\rho}+\frac{\lvert\bm v\rvert^2}{2},
 \qquad H=E+\frac{p}{\rho},
\end{equation}
则量热完全气体的状态方程也可写为
\begin{equation}
 p=(\gamma-1)\rho e
  =(\gamma-1)\left(\rho E-\frac{\rho\lvert\bm v\rvert^2}{2}\right).
\end{equation}
文献中有时用 $E,H$ 表示单位体积总能和总焓；本书统一用它们表示单位质量量。

将第一定律除以温度，并利用理想气体关系，可把 $dq/T$ 写成全微分。因此定义熵
\begin{equation}
 ds\equiv\frac{dq}{T}
  =d\!\left(c_v\ln T-R\ln\rho\right).
\end{equation}
该等式以可逆过程为前提。对含传热或黏性耗散的不可逆过程，热力学第二定律给出
\begin{equation}
 ds=\frac{dq}{T}+ds',\qquad ds'\ge 0.
\end{equation}
绝热过程满足 $ds=ds'\ge0$；若又可逆，则 $ds=0$，即等熵过程。对完全气体，
\begin{equation}
 \frac{\gamma-1}{R}ds
 =\frac{dp}{p}-\gamma\frac{d\rho}{\rho}
 =\gamma\frac{dT}{T}-(\gamma-1)\frac{dp}{p}.
\end{equation}
因此，等熵过程中
\begin{equation}
 \frac{T}{\rho^{\gamma-1}}=\text{常数},\qquad
 \frac{p}{\rho^\gamma}=\text{常数},\qquad
 \frac{T^\gamma}{p^{\gamma-1}}=\text{常数}.
\end{equation}
这些等熵（或绝热）关系适用于许多实际流动。例如，翼型超声速绕流除激波内部及
薄黏性/热边界层外，常可视为绝热可逆。曲激波会使不同流线获得不同的熵增，但跨过
激波以后，每条流线上仍可分别等熵；此时上式中的常数随流线而异。

\section{声速}

声速是微小压力扰动在气体中的传播速度。声波远弱于激波，过程可视为等熵，因而
\begin{equation}
 a^2=\left(\frac{\partial p}{\partial\rho}\right)_s.
\end{equation}
对完全气体，由 $p/\rho^\gamma=\text{常数}$ 得
\begin{equation}
 a^2=\gamma\frac{p}{\rho}=\gamma RT.
\end{equation}
本书通常用 $a$ 表示声速；有时为避免与其他变量冲突而用 $c$。若声速在整个流场中
恒定（如等温流），用 $c$ 尤其方便。

\section{欧拉方程}

忽略黏性和热传导以后，质量、动量和总能守恒分别为
\begin{align}
 \partial_t\rho+\nabla\cdot(\rho\bm v)&=0,\\
 \partial_t(\rho\bm v)+\nabla\cdot(\rho\bm v\otimes\bm v)+\nabla p&=0,\\
 \partial_t(\rho E)+\nabla\cdot(\rho\bm v H)&=0.
\end{align}
虽然 Euler 本人最初只推导了动量方程，通常把连续方程、动量方程和能量方程的整体
都称为欧拉方程。

采用物质导数 $D/Dt=\partial_t+\bm v\cdot\nabla$，原始变量形式为
\begin{align}
 \frac{D\rho}{Dt}+\rho\nabla\cdot\bm v&=0,\\
 \rho\frac{D\bm v}{Dt}+\nabla p&=0,\\
 \frac{Dp}{Dt}+\rho c^2\nabla\cdot\bm v&=0.
\end{align}
总能方程可等价替换为总焓、比总能、内能、焓、比总焓或熵方程：
\begin{align}
 \partial_t(\rho H)-\partial_t p+\nabla\cdot(\rho\bm vH)&=0,\\
 \rho\frac{DE}{Dt}+\nabla\cdot(p\bm v)&=0,\\
 \rho\frac{De}{Dt}+p\nabla\cdot\bm v&=0,\\
 \rho\frac{Dh}{Dt}-\frac{Dp}{Dt}&=0,\\
 \rho\frac{DH}{Dt}-\partial_t p&=0.
\end{align}
由 $ds=de/T-p\,d\rho/(T\rho^2)$、内能方程和连续方程得到
\begin{equation}
 \rho\frac{Ds}{Dt}=0.
\end{equation}
对含不可逆过程的一般情形，根据热力学第二定律应理解为
$\rho Ds/Dt\ge0$。

\section{一维欧拉方程}

\subsection{守恒形式}

一维方程写为
\begin{equation}
 \bm U_t+\bm F_x=0,
 \qquad
 \bm U=\begin{bmatrix}\rho\\ \rho u\\ \rho E\end{bmatrix},
 \quad
 \bm F=\begin{bmatrix}\rho u\\ \rho u^2+p\\ \rho uH\end{bmatrix}.
\end{equation}
通量雅可比 $\bm A=\partial\bm F/\partial\bm U$ 可分解为
$\bm A=\bm R\bm\Lambda\bm L$，其特征速度为 $u-c,u,u+c$。完整雅可比、左右特征
向量以及波强度如下保留；它们也是后续迎风通量与特征边界条件的基础。

\sourcecrop{76}{55}{92}{558}{673}{0.52}
\sourcecrop{77}{65}{405}{548}{716}{0.31}

其中 $\bm L\,d\bm U$ 的中间分量
$(dp-c^2d\rho)/c^2$ 与熵变化成正比。相应右特征向量说明熵波使守恒变量按
$[1,u,u^2/2]^\mathsf T$ 的比例变化。该波以局部流速 $u$ 传播，故称熵波；在激波管
问题中，它表现为携带密度跃变的接触间断，也常称接触波。

\subsection{原始变量形式}

令 $\bm W=[\rho,u,p]^\mathsf T$，则
\begin{equation}
 \bm W_t+\bm A_w\bm W_x=0,
 \qquad
 \bm A_w=\begin{bmatrix}
 u&\rho&0\\0&u&1/\rho\\0&\rho c^2&u
 \end{bmatrix}.
\end{equation}
$\bm A_w$ 并非某一通量对 $\bm W$ 的雅可比，所以更准确的名称是原始变量方程的
系数矩阵。它与守恒雅可比具有相同的三个特征值。原始变量下，熵波的变化方向为
$[d\rho,du,dp]^\mathsf T\propto[1,0,0]^\mathsf T$：接触波只改变密度，不改变速度
和压力。

变量变换满足
\begin{equation}
 d\bm U=\frac{\partial\bm U}{\partial\bm W}d\bm W,
 \qquad
 d\bm W=\frac{\partial\bm W}{\partial\bm U}d\bm U,
 \qquad
 \bm A_w=\frac{\partial\bm W}{\partial\bm U}\bm A
          \frac{\partial\bm U}{\partial\bm W}.
\end{equation}
因此两矩阵相似，特征值相同，且左右特征向量由相应变量变换矩阵互相转换。

\subsection{特征变量}

左乘 $\bm L_w$ 可把方程完全对角化：
\begin{equation}
 \bm W_{c,t}+\bm\Lambda\bm W_{c,x}=0,
 \qquad d\bm W_c=\bm L_w\,d\bm W.
\end{equation}
三个微分特征变量为
\begin{equation}
 du-\frac{dp}{\rho c},\qquad
 d\rho-\frac{dp}{c^2},\qquad
 du+\frac{dp}{\rho c},
\end{equation}
分别沿 $u-c$、$u$、$u+c$ 传播。中间方程是熵的对流，另外两个是声波。这说明熵
沿流线输运，不受弱声波影响；只有当声波强到形成激波时才会产生熵。

对绝热流，特征变量可积分为黎曼不变量
\begin{equation}
 R^- =u-\frac{2c}{\gamma-1},\qquad s,\qquad
 R^+=u+\frac{2c}{\gamma-1},
\end{equation}
并分别满足沿对应特征速度的对流方程。这组关系常用于 CFD 的特征边界条件。
当 $\gamma=3$ 时，三个方程化为两个关于 $u\mp c$ 的 Burgers 方程和一个熵输运
方程。若假定简单波 $R^-=$ 常数且 $s=$ 常数，则整个系统进一步退化为单个
$u+c$ 的 Burgers 方程；书中从 $R^-_\infty$ 解出 $c$，再代入 $R^+$ 方程验证了这一点。

\subsection{参数向量}

定义 Roe 参数向量
\begin{equation}
 \bm Z=\sqrt\rho\begin{bmatrix}1\\u\\H\end{bmatrix}
      =\begin{bmatrix}z_1\\z_2\\z_3\end{bmatrix}.
\end{equation}
守恒变量和物理通量的每个分量都可写成 $z_i$ 的二次式，例如
\begin{equation}
 \rho=z_1^2,\quad \rho u=z_1z_2,\quad \rho uH=z_2z_3,\quad
 p=\frac{\gamma-1}{\gamma}\left(z_1z_3-\frac{z_2^2}{2}\right).
\end{equation}
这种二次齐次结构特别适合构造欧拉方程的线性化。

\section{二维欧拉方程}

\subsection{守恒形式}

二维系统为
\begin{equation}
 \bm U_t+\bm F_x+\bm G_y=0,
\end{equation}
其中
\begin{equation}
 \bm U=\begin{bmatrix}\rho\\\rho u\\\rho v\\\rho E\end{bmatrix},\quad
 \bm F=\begin{bmatrix}\rho u\\\rho u^2+p\\\rho uv\\\rho uH\end{bmatrix},\quad
 \bm G=\begin{bmatrix}\rho v\\\rho uv\\\rho v^2+p\\\rho vH\end{bmatrix}.
\end{equation}
沿单位法向 $\bm n=[n_x,n_y]^\mathsf T$ 的法向速度为
$q_n=un_x+vn_y$，法向通量为
\begin{equation}
 \bm F_n=\bm F n_x+\bm G n_y=
 \begin{bmatrix}\rho q_n\\\rho q_n u+p n_x\\\rho q_n v+p n_y\\\rho q_n H\end{bmatrix}.
\end{equation}
记 $\bm A=\partial\bm F/\partial\bm U$、$\bm B=\partial\bm G/\partial\bm U$，则
$\bm A_n=\bm A n_x+\bm Bn_y$。取 $\bm n=(1,0)$ 或 $(0,1)$ 即分别恢复 $\bm A$、
$\bm B$，所以法向雅可比是一种特别方便的统一表示。

\sourcecrop{81}{48}{78}{564}{707}{0.58}
\sourcecrop{82}{52}{315}{557}{715}{0.40}

$\bm A_n=\bm R_n\bm\Lambda_n\bm L_n$ 的特征速度为
$q_n-c,q_n,q_n+c,q_n$。令切向单位向量
$\bm\ell=[-n_y,n_x]^\mathsf T$，$q_\ell=\bm v\cdot\bm\ell$；四个波强分别对应
左行声波、熵/接触波、右行声波和切向速度波。最后一个波携带切向速度跃变，因而称
剪切波或涡量波。法向雅可比虽可对角化，却无法同时对角化 $\bm A$ 与 $\bm B$；这
正是一维系统可完全分解为标量对流方程，而二维非定常系统难以做到的原因。

\subsection{原始变量形式与对称形式}

取 $\bm W=[\rho,u,v,p]^\mathsf T$，有
\begin{equation}
 \bm W_t+\bm A_w\bm W_x+\bm B_w\bm W_y=0.
\end{equation}
$\bm A_w,\bm B_w$ 是系数矩阵而不是通量雅可比。法向系数矩阵
$\bm A_{wn}=\bm A_w n_x+\bm B_w n_y$ 仍具有 $q_n-c,q_n,q_n+c,q_n$ 四个特征速度；
其第四个特征幅值是 $\rho\,dq_\ell$，明确表示切向剪切/涡量波。守恒变量与原始变量
的雅可比变换、左右特征向量关系以及完整矩阵见下式。

\sourcecrop{83}{52}{75}{560}{705}{0.57}
\sourcecrop{84}{54}{85}{558}{690}{0.53}

对称双曲形式可取变量
\begin{equation}
 d\bm V=\begin{bmatrix}dp/(\rho c)\\du\\dv\\dp-c^2d\rho\end{bmatrix},
\end{equation}
从而把时间系数化为单位阵，并得到对称的空间系数矩阵。这种形式在稳定性与能量
估计中很有价值。进一步把局部坐标旋转到流线方向 $\xi$ 和法向 $\eta$，其中
\begin{equation}
 q=u\cos\theta+v\sin\theta,\qquad
 q_\ell=-u\sin\theta+v\cos\theta,\qquad
 \theta=\tan^{-1}(v/u),
\end{equation}
可显著简化横向系数矩阵。按定义流线法向速度 $q_\ell=0$，但它沿流线的变化率并
不必为零。

\subsection{参数向量}

二维参数向量为
\begin{equation}
 \bm Z=\sqrt\rho[1,u,v,H]^\mathsf T.
\end{equation}
守恒变量和两方向通量均为 $\bm Z$ 的二次式，而参数变量通量雅可比
$\bm A^z=\partial\bm F/\partial\bm Z$、$\bm B^z=\partial\bm G/\partial\bm Z$ 对
$\bm Z$ 是线性的。这使单元内的 Roe 型线性化特别自然。齐次二次结构还给出
\begin{equation}
 \bm Z=2\frac{\partial\bm Z}{\partial\bm U}\bm U,\qquad
 \bm U=\frac12\frac{\partial\bm U}{\partial\bm Z}\bm Z,\qquad
 \bm F=\frac12\bm A^z\bm Z,\qquad
 \bm G=\frac12\bm B^z\bm Z.
\end{equation}
因此若 CFD 程序直接以参数向量作为工作变量，无需先恢复全部原始变量即可计算
物理通量。

\sourcecrop{86}{55}{80}{557}{690}{0.52}
\sourcecrop{87}{55}{80}{558}{696}{0.53}

\section{三维欧拉方程}

\subsection{守恒形式}

三维系统为
\begin{equation}
 \bm U_t+\bm F_x+\bm G_y+\bm H_z=0,
\end{equation}
其中 $\bm U=[\rho,\rho u,\rho v,\rho w,\rho E]^\mathsf T$；三方向通量分别由质量
通量、三个动量通量和能量通量构成。沿
$\bm n=[n_x,n_y,n_z]^\mathsf T$ 的速度 $q_n=un_x+vn_y+wn_z$，法向通量为
\begin{equation}
 \bm F_n=
 \begin{bmatrix}
 \rho q_n\\ \rho q_n u+p n_x\\ \rho q_n v+p n_y\\
 \rho q_n w+p n_z\\ \rho q_n H
 \end{bmatrix}.
\end{equation}
法向雅可比 $\bm A_n=\bm A n_x+\bm Bn_y+\bm Cn_z$ 统一包含三个坐标方向；取
坐标基向量即可恢复 $\bm A,\bm B,\bm C$。

\sourcecrop{89}{45}{80}{567}{710}{0.60}
\sourcecrop{90}{48}{80}{565}{710}{0.60}

取与 $\bm n$ 正交的两个单位切向量 $\bm\ell,\bm m$，并定义
$q_\ell=\bm v\cdot\bm\ell$、$q_m=\bm v\cdot\bm m$。三维法向雅可比的五个特征速度
为 $q_n-c,q_n,q_n+c,q_n,q_n$；后三个 $q_n$ 波分别携带熵及两个切向速度分量。
完整左右特征矩阵如下。

\sourcecrop{91}{48}{75}{565}{710}{0.61}

三维中切向基并不唯一。绝对雅可比
$|\bm A_n|=\bm R_n|\bm\Lambda_n|\bm L_n$ 表面上含 $\bm\ell,\bm m$，但把两个剪切
波的外积相加，并用
$\bm\ell\otimes\bm\ell+\bm m\otimes\bm m=\bm I-\bm n\otimes\bm n$，即可完全消去
切向量。这样，迎风耗散项只依赖唯一的法向量，避免三维切向基选择的歧义。相同技巧
也表明 $|\bm A_n|d\bm U$ 必能写成只含 $\bm n$ 的形式；这正是 Roe 通量等迎风格式
实际需要的量。

\sourcecrop{92}{45}{72}{567}{710}{0.61}
\sourcecrop{93}{48}{80}{563}{710}{0.59}

\subsection{展开形式与原始变量形式}

把守恒形式按 $\rho,u,v,w,p$ 的导数展开，可逐项看出连续方程、三个动量方程和能量
方程如何组合。利用连续方程消去重复项，即得到熟悉的物质导数形式。取
$\bm W=[\rho,u,v,w,p]^\mathsf T$ 后，原始变量系统写为
\begin{equation}
 \bm W_t+\bm A_w\bm W_x+\bm B_w\bm W_y+\bm C_w\bm W_z=0.
\end{equation}
法向系数矩阵有相同的五个特征速度和相同的声、熵、双剪切波物理解释。守恒/原始
变量之间的变换仍为相似变换，因此不改变特征值。

\sourcecrop{94}{47}{80}{566}{705}{0.59}
\sourcecrop{95}{48}{80}{565}{710}{0.60}

\subsection{对称形式}

三维对称变量可取
\begin{equation}
 d\bm V=\begin{bmatrix}dp/(\rho c) & du & dv & dw & dp-c^2d\rho\end{bmatrix}^{\!\mathsf T}.
\end{equation}
对原始变量方程施加相应变量变换，可得到实对称的三个方向系数矩阵。旋转到局部
流线坐标后，纵向速度为 $q$，另外两个为 $q_\ell,q_m$；在流线上
$q_\ell=q_m=0$，横向矩阵因而显著简化。完整变换与矩阵保留如下。

\sourcecrop{96}{48}{78}{565}{710}{0.60}
\sourcecrop{97}{48}{78}{565}{710}{0.60}
\sourcecrop{98}{48}{80}{565}{710}{0.60}
\sourcecrop{99}{48}{80}{565}{710}{0.60}

\subsection{参数向量}

三维参数向量为
\begin{equation}
 \bm Z=\sqrt\rho[1,u,v,w,H]^\mathsf T.
\end{equation}
与一、二维完全类似，$\bm U,\bm F,\bm G,\bm H$ 的各分量都是参数分量的二次式，
参数空间雅可比对 $\bm Z$ 线性，且通量可直接由雅可比乘参数向量得到一半。书中列出
了所有正、逆变量变换矩阵和三个方向的参数雅可比。

\sourcecrop{100}{48}{78}{565}{710}{0.60}
\sourcecrop{101}{48}{78}{565}{710}{0.60}
\sourcecrop{102}{48}{78}{565}{710}{0.60}

\subsection{特征值的导数}

一些隐式方法、线性化通量和通量差分格式需要特征值对状态变量的导数。由于
$\lambda=q_n\pm c$ 或 $q_n$，先对法向速度和声速求微分，再通过
$d\bm W=(\partial\bm W/\partial\bm U)d\bm U$ 转为守恒变量导数。对 Roe 平均状态，
则需同时考虑左右状态通过参数平均对 $\hat u,\hat v,\hat w,\hat H,\hat c$ 的影响。
原书依次给出单状态导数、Roe 平均量及其对左右守恒变量的完整链式法则；为避免长
矩阵转写失真，公式完整保留如下。

\sourcecrop{103}{47}{75}{566}{710}{0.61}
\sourcecrop{104}{47}{75}{566}{710}{0.61}
\sourcecrop{105}{47}{75}{566}{710}{0.61}
\sourcecrop{106}{47}{75}{566}{710}{0.61}
\sourcecrop{107}{47}{75}{566}{710}{0.61}
\sourcecrop{108}{47}{75}{566}{710}{0.61}

这里帽号表示 Roe 平均量，下标 $L,R$ 分别表示界面左右状态；$\lambda_k$ 是第 $k$
个特征值，$\bm r_k$、$\bm\ell_k^\mathsf T$ 是对应右、左特征向量。切向基的导数可
通过前述消去技巧避免进入最终法向耗散表达式。

\section{轴对称欧拉方程}

轴对称方程形式上是二维的，却能描述三维流动，因而计算效率很高。采用柱坐标
$(r,\theta,z)$，速度为 $[u_r,u_\theta,u_z]^\mathsf T$。若关于 $z$ 轴轴对称，所有
变量对 $\theta$ 无变化且 $u_\theta=0$，于是周向动量方程消失。乘以半径 $r$ 后，方程
成为守恒形式；把 $(z,r)$ 重命名为 $(x,y)$、把 $(u_z,u_r)$ 重命名为 $(u,v)$，可写成
\begin{equation}
 (y\bm U)_t+(y\bm F)_x+(y\bm G)_y=\bm S,
 \qquad
 \bm S=[0,0,p,0]^\mathsf T.
\end{equation}
它与二维欧拉方程的差别仅是守恒量及通量多乘了 $y$，并在径向动量方程中出现压力
几何源项。可引入开关 $\omega$ 统一表示二维和平面轴对称形式：
\begin{equation}
 (\omega_a\bm U)_t+(\omega_a\bm F)_x+(\omega_a\bm G)_y=\omega\bm S,
 \qquad \omega_a=(1-\omega)+\omega y,
\end{equation}
其中 $\omega=0$ 为二维平面流，$\omega=1$ 为轴对称流。这意味着已有二维程序只需
很少改动便可计算本质上的三维轴对称问题。

\sourcecrop{109}{48}{80}{565}{710}{0.60}
\sourcecrop{110}{48}{80}{565}{710}{0.60}

\section{旋转不变性}

欧拉通量具有旋转不变性：
\begin{equation}
 \bm F_n=n_x\bm F+n_y\bm G+n_z\bm H
 =\bm T^{-1}\bm F(\bm T\bm U),
\end{equation}
其中 $\bm T$ 把速度分量从全局坐标旋转到由 $\bm n,\bm\ell,\bm m$ 构成的正交局部
坐标。先计算 $\bm T\bm U=[\rho,\rho q_n,\rho q_\ell,\rho q_m,\rho E]^\mathsf T$，用
单一方向通量函数 $\bm F$ 求值，再旋回全局坐标，就能得到任意方向的法向通量。
有限体积程序因此只需实现一个一维方向的界面通量例程；也可直接用法向通量显式式，
后者在程序中往往更直观。

\sourcecrop{111}{50}{82}{563}{710}{0.59}

\section{便于降维的分量排序}

若把守恒变量排为 $[\rho,\rho E,\rho u,\rho v,\rho w]^\mathsf T$，并以相同原则重排
通量，则删去第五行及 $z$ 向通量即可得到二维方程，再删去第四行及 $y$ 向通量即可
得到一维方程。雅可比和特征向量同样可以直接截取。这种排序使同一三维代码仅通过
限制数组大小便能运行一、二维问题。另一种做法是用一列三维单元模拟一维、用单位
厚度棱柱层模拟二维，但分量排序仍很有用。例如在 MHD 扩展系统中，把磁场分量附加
到纯流体状态向量末尾，就可通过限制状态维数运行纯气动力学算例。

\sourcecrop{112}{47}{80}{566}{710}{0.60}

\section{欧拉通量的齐次性质}

若 $p/\rho$ 只依赖内能 $e$，即 $p=\rho f(e)$，则欧拉通量关于守恒变量是一阶齐次的：
\begin{equation}
 \bm F(\lambda\bm U)=\lambda\bm F(\bm U),\quad
 \bm G(\lambda\bm U)=\lambda\bm G(\bm U),\quad
 \bm H(\lambda\bm U)=\lambda\bm H(\bm U).
\end{equation}
对 $\lambda$ 求导并令 $\lambda=1$，Euler 齐次定理给出
\begin{equation}
 \bm A\bm U=\bm F,\qquad \bm B\bm U=\bm G,\qquad \bm C\bm U=\bm H.
\end{equation}
所以守恒方程既可写成 $\bm U_t+\bm A\bm U_x+\bm B\bm U_y+\bm C\bm U_z=0$，也可
写成 $\bm U_t+(\bm A\bm U)_x+(\bm B\bm U)_y+(\bm C\bm U)_z=0$，尽管雅可比本身
并非常矩阵。Steger--Warming 通量矢量分裂正利用了该性质。比较
$d\bm F=d(\bm A\bm U)=\bm A\,d\bm U+d\bm A\,\bm U$ 与
$d\bm F=\bm A\,d\bm U$，还得到 $d\bm A\,\bm U=0$。

\section{无量纲化}

\subsection{无量纲欧拉方程}

无量纲化把自由来流速度等基准量缩放为一，避免程序处理过大或过小数值。取参考
密度 $\rho_r$、速度 $v_r$、长度 $L_r$，并令
\begin{equation}
 \rho^*=\rho/\rho_r,\quad \bm v^*=\bm v/v_r,\quad
 p^*=p/(\rho_r v_r^2),\quad E^*=E/v_r^2,\quad
 t^*=tv_r/L_r,
\end{equation}
则代入有量纲方程后所有公共尺度因子消去，无量纲欧拉方程保持完全相同的守恒形式。
温度和气体常数的缩放必须协调，使 $p^*=\rho^*R^*T^*$ 仍成立；声速也满足
$a^{*2}=\gamma p^*/\rho^*$。

\subsection{自由来流、声速与滞止量}

若以自由来流密度、声速和特征长度作参考量，则
$\rho_\infty^*=1$、$a_\infty^*=1$，自由来流速度的模为马赫数
$M_\infty$。因此入口只需给定 $M_\infty$ 与流向单位向量。若改以自由来流速度作
速度尺度，则 $\lvert\bm v_\infty^*\rvert=1$，而
$a_\infty^*=1/M_\infty$。压强尺度选择不同会使无量纲自由来流压强分别为
$1/\gamma$、$1/(\gamma M_\infty^2)$ 等，但方程形式不变。

滞止量来自等熵减速到零速度的状态。对完全气体，
\begin{equation}
 \frac{T_0}{T}=1+\frac{\gamma-1}{2}M^2,\quad
 \frac{p_0}{p}=\left(1+\frac{\gamma-1}{2}M^2\right)^{\gamma/(\gamma-1)},
\end{equation}
密度比为同一括号的 $1/(\gamma-1)$ 次幂。以滞止量作参考时，入口同样只需给定来流
马赫数和方向，并便于直接比较计算的滞止点压强。

\sourcecrop{113}{48}{80}{565}{710}{0.60}
\sourcecrop{114}{48}{80}{565}{710}{0.60}
\sourcecrop{115}{48}{405}{565}{710}{0.40}

\section{变量转换}

本节汇总守恒变量 $\bm U$、原始变量 $\bm W$ 与参数向量 $\bm Z$ 之间的显式变换，
并且这些关系对上一节各种一致的无量纲定义均成立。

\subsection{一维}

令 $\bm U=[u_1,u_2,u_3]^\mathsf T=[\rho,\rho u,\rho E]^\mathsf T$，则
\begin{equation}
 \rho=u_1,\quad u=u_2/u_1,\quad
 p=(\gamma-1)\left(u_3-\frac{u_2^2}{2u_1}\right),
\end{equation}
且 $z_1=\sqrt{u_1}$、$z_2=u_2/\sqrt{u_1}$，
$z_3=\sqrt{u_1}\,[\gamma u_3/u_1-(\gamma-1)(u_2/u_1)^2/2]$。
从原始变量 $[w_1,w_2,w_3]=[\rho,u,p]$ 出发，
$\rho E=w_3/(\gamma-1)+w_1w_2^2/2$；从参数向量出发，
$\rho=z_1^2$、$u=z_2/z_1$、$H=z_3/z_1$，其余量由前述二次关系得到。

\subsection{二维与三维}

二维时，守恒变量到原始变量的转换为
\begin{equation}
 \rho=u_1,\quad u=u_2/u_1,\quad v=u_3/u_1,\quad
 p=(\gamma-1)\left[u_4-\frac{u_2^2+u_3^2}{2u_1}\right].
\end{equation}
三维时再加入 $w=u_4/u_1$，并把压力式中的动能改为
$(u_2^2+u_3^2+u_4^2)/(2u_1)$。原始变量到守恒变量的通式是
\begin{equation}
 \rho E=\frac{p}{\gamma-1}+\frac{\rho}{2}(u^2+v^2+w^2).
\end{equation}
参数向量始终为 $\sqrt{\rho}\,[1,u,v,w,H]^\mathsf T$；删去不用的速度分量即得到低维
版本。原书把一、二、三维的全部逐分量公式列于下列页面，便于直接照搬到程序中。

\sourcecrop{115}{48}{75}{565}{405}{0.33}
\sourcecrop{116}{48}{75}{565}{710}{0.60}
\sourcecrop{117}{48}{75}{565}{710}{0.60}
\sourcecrop{118}{48}{75}{565}{710}{0.60}

\section{不可压缩/伪可压缩欧拉方程}

不可压缩流体的密度变化可以忽略，例如常温水。等价地说，流动马赫数很小；通常
$M<0.3$ 即可近似视为不可压缩。令 $\rho$ 为常数，欧拉方程化为
\begin{align}
 \nabla\cdot(\rho\bm v)&=0,\\
 \partial_t(\rho\bm v)+\nabla\cdot(\rho\bm v\otimes\bm v)+\nabla p&=0,
\end{align}
或者
\begin{equation}
 \nabla\cdot\bm v=0,\qquad
 \rho\frac{D\bm v}{Dt}+\nabla p=0.
\end{equation}
该系统无需能量方程即可闭合，温度可在求出流场后另行计算。不过，即使流动非定常，
连续方程也不含时间导数，因而不能直接把整个系统作为普通双曲初值问题推进。

Chorin 的伪可压缩方法引入常数人工声速 $a^*$ 和人工密度
$\rho^*=p/a^{*2}$，并在人工时间 $t^*$ 中写成
\begin{equation}
 \frac{\partial P}{\partial t^*}+\nabla\cdot(a^{*2}\bm v)=0,
 \qquad P=p/\rho,
\end{equation}
其中 $P$ 称运动压强。这样系统恢复双曲性质，可用时间推进法求稳态解；只有人工时间
达到稳态时，原连续条件 $\nabla\cdot\bm v=0$ 才真正满足。$a^*$ 取值依问题而定，通常
在 $0.1$ 至 $10$ 之间；也有文献用 $\beta=\rho a^{*2}$ 表示人工可压缩参数。

伪可压缩欧拉方程的守恒形式为
\begin{equation}
 \bm U_{t^*}+\bm F_x+\bm G_y+\bm H_z=0,
 \quad \bm U=[P,u,v,w]^\mathsf T,
\end{equation}
其法向声学特征速度为
$q_n-c,q_n,q_n,q_n+c$，这里 $c=\sqrt{q_n^2+a^{*2}}$，并非真实热力学声速。
完整通量、雅可比、左右特征向量与波幅如下。切向量同样可在
$|\bm A_n|d\bm U$ 中消去，最终迎风耗散只依赖法向量。

\sourcecrop{119}{47}{75}{566}{710}{0.61}
\sourcecrop{120}{47}{75}{566}{710}{0.61}
\sourcecrop{121}{47}{310}{566}{710}{0.40}

该守恒伪时间系统只在稳态时与不可压缩原始方程等价。动量方程中的物理时间导数在
这里也统一写作人工时间导数，因为与伪连续方程同时求解时，未收敛的中间状态没有
真实非定常意义。

\section{均熵/等温欧拉方程}

任意维欧拉方程都包含熵输运
\begin{equation}
 s_t+us_x+vs_y+ws_z=0,
\end{equation}
所以熵沿流线不变，但不同流线可以具有不同的熵。这类流动称等熵流。曲激波后的
典型流动即如此：每条流线离开激波后熵保持不变，而不同位置的激波强度不同，使各
流线熵值不相同。若初始熵在整个流场均匀，且以后也不产生激波，则熵处处为同一常数，
称均熵流（homentropic flow）。

对均熵流，
\begin{equation}
 p/\rho^\gamma=K
\end{equation}
中的 $K$ 是全局常数，而非仅沿流线为常数。因此总能可由密度与速度显式恢复，能量
方程不再需要。更一般地，只要 $p=p(\rho)$，即流体为正压流体，系统便可在没有能量
方程时闭合，声速为 $c=\sqrt{p'(\rho)}$；均熵完全气体则
$c=\sqrt{\gamma K\rho^{\gamma-1}}$。

$\gamma=1$ 对应等温模型。形式上这意味着分子自由度和比热趋于无穷，温度极难改变；
此时
\begin{equation}
 p=c^2\rho,\qquad c^2=RT=\text{常数}.
\end{equation}
真实气体并不严格具有 $\gamma=1$，但该人工模型很有用。均熵或等温系统都只含质量、
动量方程。一维系统的两个特征速度是 $u\mp c$；二维法向系统为
$q_n-c,q_n,q_n+c$；三维再多一个 $q_n$ 剪切波。系统双曲性的条件是
$p'(\rho)>0$。各维守恒量、通量、雅可比及特征矩阵完整列示如下。

\sourcecrop{122}{48}{75}{565}{710}{0.61}
\sourcecrop{123}{48}{75}{565}{710}{0.61}
\sourcecrop{124}{48}{75}{565}{710}{0.61}

其中 $\bm n,\bm\ell,\bm m$ 为两两正交的单位向量，$q_n,q_\ell,q_m$ 为速度在这些
方向上的分量。没有能量方程以后，熵/接触波也随之消失；保留下来的 $q_n$ 重根只
对应横向速度或涡量扰动。

\section{线性声学方程（线性化欧拉方程）}

把变量分解为均匀基态和小扰动：
\begin{equation}
 \rho=\rho_0+\rho',\quad p=p_0+p',\quad
 (u,v,w)=(u_0,v_0,w_0)+(u',v',w').
\end{equation}
代入欧拉方程，舍去扰动量的二次及更高阶项，即得线性声学方程。若基态静止，原始
变量形式为
\begin{equation}
 p'_t+\rho_0c_0^2\nabla\cdot\bm v'=0,\qquad
 \bm v'_t+\frac1{\rho_0}\nabla p'=0.
\end{equation}
对第一式再求时间导数，并用第二式消去速度，得到经典波动方程
\begin{equation}
 p'_{tt}-c_0^2\nabla^2p'=0.
\end{equation}
速度各分量也满足同型波动方程。若均匀基流非零，只需把 $\partial_t$ 替换为随基流
对流的导数 $\partial_t+\bm v_0\cdot\nabla$，声波相对地面以
$q_{n,0}\mp c_0$ 传播。

一维系统只有左右传播的两支声波；二维多一个以基流法向速度传播的切向涡量模态；
三维有两个切向模态。完整的一、二、三维常系数矩阵和特征分解如下。

\sourcecrop{125}{48}{75}{565}{710}{0.61}
\sourcecrop{126}{48}{75}{565}{710}{0.61}

符号下标 $0$ 表示均匀基态，撇号表示一阶小扰动；$c_0$ 是基态声速。因为系数常数，
这些方程是检验数值色散、耗散和边界反射的理想模型。

\section{准一维欧拉方程}

准一维系统描述变截面喷管。截面积 $A(x)$ 随轴向变化，但流动变量只依赖 $x$：
\begin{equation}
 \frac{\partial\bm U}{\partial t}+\frac{\partial\bm F}{\partial x}=\bm S,
\end{equation}
其中可取
\begin{equation}
 \bm U=\begin{bmatrix}\rho A\\\rho uA\\\rho EA\end{bmatrix},
 \quad
 \bm F=\begin{bmatrix}\rho uA\\\rho u^2A\\\rho uHA\end{bmatrix},
 \quad
 \bm S=\begin{bmatrix}0\\-A p_x\\0\end{bmatrix}.
\end{equation}
也可把面积只放入几何源项，或者把压力通量改写成
$(\rho u^2+p)A$，使源项成为 $pA_x$。这些形式连续等价，但离散时的平衡性质可能
不同。若把 $\rho A$ 当作新的密度变量，除源项外系统与一维欧拉方程完全同型，故可
直接复用其特征结构。

令 $A=2\pi r$ 可得到柱对称波动方程；令 $A=4\pi r^2$ 可得到球对称波动方程。几何
扩张分别产生与 $u/r$、$2u/r$ 成正比的源项。这说明准一维模型不仅适用于喷管，也
统一包含柱面波和球面波。

\sourcecrop{127}{48}{75}{565}{710}{0.61}
\sourcecrop{128}{48}{430}{565}{710}{0.29}

\section{二维定常欧拉方程}

从二维原始变量方程删去时间导数，得到
\begin{equation}
 \bm A_w\bm W_x+\bm B_w\bm W_y=0.
\end{equation}
若 $\bm A_w$ 可逆（$u^2\ne c^2$），可把 $x$ 看作类似时间的推进方向：
\begin{equation}
 \bm W_x+\bm A_w^{-1}\bm B_w\bm W_y=0.
\end{equation}
矩阵 $\bm A_w^{-1}\bm B_w$ 的两个重特征值为 $v/u$，另外两个为
\begin{equation}
 \lambda_{3,4}=\frac{uv\mp c\sqrt{q^2-c^2}}{u^2-c^2},\qquad
 \beta=\sqrt{M^2-1}.
\end{equation}
因此定常二维欧拉方程的类型由马赫数决定：$M>1$ 时特征值为实数，系统为双曲型；
$M<1$ 时出现共轭复特征值，系统为椭圆型。一个流场中可同时存在超声速双曲区和
亚声速椭圆区，故它是混合型方程。

超声速时系统可完全特征化，四个特征方向分别对应两条流线方向及两族马赫线。令
$\zeta=v_x-u_y$ 为涡量标量，书中得到两族声学特征速度向量的散度
\begin{equation}
 \nabla\cdot\bm\lambda_3=\nabla\cdot(\beta\bm v)+\zeta,\qquad
 \nabla\cdot\bm\lambda_4=\nabla\cdot(\beta\bm v)-\zeta.
\end{equation}
特征线会聚/发散分别指示激波/膨胀波；不仅涡量或速度散度单独起作用，而是二者的
和或差决定波族。非负标量
$[\nabla\cdot\bm\lambda_3]^2+[\nabla\cdot\bm\lambda_4]^2
=2\{[\nabla\cdot(\beta\bm v)]^2+\zeta^2\}$ 可用于检测任意一族非线性波，曾用于
自适应欧拉格式。

\sourcecrop{128}{48}{75}{565}{430}{0.36}
\sourcecrop{129}{48}{75}{565}{710}{0.61}
\sourcecrop{130}{48}{430}{565}{710}{0.29}

有些方法尝试对定常方程的椭圆部分使用中心格式、对双曲部分使用迎风格式。更多 CFD
方法则基于非定常欧拉方程始终关于时间双曲这一事实构造统一迎风离散；两种观点各有
价值。

\section{伯努利方程}

伯努利方程本质上是动量方程的积分。利用恒等式
$(\nabla\bm v)\bm v=\nabla(\lvert\bm v\rvert^2/2)-\bm v\times(\nabla\times\bm v)$，
欧拉动量方程写为
\begin{equation}
 \rho\bm v_t+\rho\nabla\frac{q^2}{2}
 -\rho\bm v\times\bm\omega+\nabla p=0,
 \qquad \bm\omega=\nabla\times\bm v.
\end{equation}

对非定常、不可压缩、无旋流，$\rho$ 为常数且
$\bm v=\nabla\phi$，故
\begin{equation}
 \phi_t+\frac{q^2}{2}+\frac{p}{\rho}=f(t).
\end{equation}
稳态时即著名的 $q^2/2+p/\rho=$ 常数。对稳态不可压缩有旋流，该量不再全场恒定，
但点乘速度后涡量项消失，说明它沿每条流线恒定；停滞压
$p_0=p+\rho q^2/2$ 可在不同流线间取不同常数。其法向变化为
\begin{equation}
 \frac{\partial p_0}{\partial n}=\rho\bm n\cdot(\bm v\times\bm\omega),
\end{equation}
所以涡量是停滞压横向变化的来源。若为 Beltrami 流
$\bm v\times\bm\omega=0$，即使涡量非零，$p_0$ 仍可处处相同。

对绝热完全气体，$dp/\rho=dh=d(a^2)/(\gamma-1)$。因此非定常可压缩无旋流满足
\begin{equation}
 \phi_t+\frac{q^2}{2}+\frac{a^2}{\gamma-1}=g(t),
\end{equation}
稳态有旋流则把该量理解为沿流线常数，其法向导数仍为
$\bm n\cdot(\bm v\times\bm\omega)$。

\sourcecrop{130}{48}{75}{565}{430}{0.36}
\sourcecrop{131}{48}{75}{565}{710}{0.61}
\sourcecrop{132}{48}{430}{565}{710}{0.29}

\section{气体动力学方程（非线性势方程）}

若流动处处无旋，可定义速度势 $\bm v=\nabla\phi$。连续方程变成
\begin{equation}
 \nabla\cdot(\rho\nabla\phi)=0.
\end{equation}
它需与能量关系闭合。对定常、绝热、量热完全气体，
\begin{equation}
 \frac{a^2}{\gamma-1}+\frac{u^2+v^2+w^2}{2}=H=\text{常数},
\end{equation}
再结合等熵关系即可把 $\rho$ 表示为速度和自由来流量的函数。该模型要求全场均熵，
严格说不适用于含激波流动；不过在跨声速附近，它作为近似仍曾成功使用。二维时也可
写成由质量守恒与无旋条件构成的 Cauchy--Riemann 型系统。

展开连续方程，并将动量方程点乘速度，可消去密度梯度，得到气体动力学方程
\begin{equation}
 \nabla\cdot\bm v-\frac1{a^2}\bm v\cdot[(\nabla\bm v)\bm v]=0.
\end{equation}
代入 $\bm v=\nabla\phi$ 得到完整势方程；其三维展开式包含三个纯二阶导数和三个混合
导数，系数由局部速度与声速决定。该标量方程配合密度关系即可求解速度势，比直接
求解多变量系统容易。无旋势分解原则甚至可扩展到某些黏性流或有旋流，但需要额外的
旋转速度分量。

\sourcecrop{132}{48}{75}{565}{430}{0.36}
\sourcecrop{133}{48}{75}{565}{710}{0.61}

\section{线性势方程}

考虑自由来流沿 $x$ 方向掠过细长物体，令
$u=U_\infty+u'$、$v=v'$、$w=w'$，且扰动远小于 $U_\infty$。写
$\phi=U_\infty x+\phi'$，对完整势方程线性化。二维一般得到
\begin{equation}
 (1-M_\infty^2)\phi'_{xx}+\phi'_{yy}=0,
\end{equation}
三维则为
\begin{equation}
 (1-M_\infty^2)\phi'_{xx}+\phi'_{yy}+\phi'_{zz}=0.
\end{equation}
跨声速 $M_\infty\approx1$ 时，$x$ 向线性项很小，必须保留原展开中的主要非线性项
$M_\infty^2(\gamma+1)\phi'_x\phi'_{xx}/U_\infty$，否则方程会错误地近似退化为一维。
求得 $\phi'$ 后，可由扰动速度计算压强系数。小马赫数时该方程接近 Laplace 方程，
容易求解；大马赫数时它转为双曲型。

\sourcecrop{134}{48}{75}{565}{710}{0.61}

\section{小扰动方程}

把线性势方程改写成速度分量形式，即 Prandtl--Glauert 小扰动方程。超声速时令
$\beta=\sqrt{M_\infty^2-1}$，二维系统为
\begin{equation}
 \beta^2u_x-v_y=0,\qquad v_x-u_y=0.
\end{equation}
其特征组合 $\beta u-v$ 和 $\beta u+v$ 分别沿
$dy/dx=1/\beta$ 与 $dy/dx=-1/\beta$ 保持不变，只有 $M_\infty>1$ 时该实特征分解
才成立。若 $M_\infty=\sqrt2$，则 $\beta=1$，表达式尤其简单。

亚声速时，方程为 $(1-M_\infty^2)u_x+v_y=0$ 与 $v_x-u_y=0$，特征值为复数，系统
为椭圆型。采用 Prandtl--Glauert 变换
\begin{equation}
 \xi=x,\quad \eta=y\sqrt{1-M_\infty^2},\quad
 U=u\sqrt{1-M_\infty^2},\quad V=v,
\end{equation}
可把任意 $|M_\infty|<1$ 的系统化为标准 Cauchy--Riemann 方程，说明亚声速小扰动
流可映射为不可压缩势流。

\sourcecrop{135}{48}{430}{565}{710}{0.29}

\section{不可压缩势方程}

定常不可压缩流满足 $\nabla\cdot\bm v=0$。若又无旋，令
$\bm v=\nabla\phi$，则
\begin{equation}
 \nabla^2\phi=0.
\end{equation}
即速度势在任意维均满足 Laplace 方程，因此可以利用大量精确解。二维还可同时使用
速度势和流函数；二者满足 Cauchy--Riemann 方程，可用复变函数理论构造解。该理论
一般不能直接用于三维或轴对称流，因为相应流函数不再满足 Laplace 方程。

\section{Crocco 方程}

利用
$(\nabla\bm v)\bm v=\nabla(q^2/2)-\bm v\times\bm\omega$、
$T\nabla s=\nabla h-\nabla p/\rho$ 以及 $H=h+q^2/2$，动量方程可改写为
\begin{equation}
 \frac{\partial\bm v}{\partial t}-\bm v\times\bm\omega
 =T\nabla s-\nabla H.
\end{equation}
这就是 Crocco 方程。对定常绝热流，总焓满足
$\bm v\cdot\nabla H=0$，即沿流线守恒。再将速度点乘 Crocco 方程，左侧叉乘项消失，
得到 $\bm v\cdot\nabla s=0$，证明熵沿流线不变。反过来，若 $\bm\omega=0$ 且 $H$
全场均匀，则立即得到 $\nabla s=0$，即流动均熵。Crocco 方程只是欧拉动量方程的另一
种形式，却清楚揭示了涡量、熵梯度和总焓梯度之间的关系。

\section{Kelvin 环量定理}

沿闭合曲线 $C$ 的环量定义为
\begin{equation}
 \Gamma=\oint_C\bm v\cdot d\bm l
 =\iint_S\bm\omega\cdot\bm n\,dS,
\end{equation}
即环量等于任意以 $C$ 为边界的曲面上涡量通量。对随流体运动的闭合曲线，欧拉方程
给出
\begin{equation}
 \frac{D\Gamma}{Dt}=-\oint_C\frac{\nabla p}{\rho}\cdot d\bm l.
\end{equation}
不可压缩无黏流中 $\rho$ 常数，右端是压力全微分的闭合积分，故为零。可压缩流中用
$dp/\rho=dh-Tds$，若熵均匀，右端同样为零。因此随体闭曲线的环量保持不变，这就是
Kelvin 环量定理。黏性流还会多出黏性应力散度的闭合积分贡献。

涡量积分曲线称涡线；穿过闭合曲线的一束涡线形成涡管（或涡丝），其强度等于截面
涡量通量，也等于环量。环量定理直接导出 Helmholtz 三定理：
\begin{enumerate}
 \item 若无旋转体力，初始无旋的流动保持无旋；
 \item 涡管强度沿管长保持不变；
 \item 涡管不能终止于流体内部，只能延伸至边界或闭合成环。
\end{enumerate}
第二条还说明，当涡管截面收缩/扩张时涡量增强/减弱，这就是与涡管拉伸相联系的涡
伸展机制。

\sourcecrop{136}{48}{75}{565}{710}{0.61}
\sourcecrop{137}{48}{430}{565}{710}{0.29}

\section{涡量方程}

对原始变量动量方程取旋度，并用 $\bm\omega=\nabla\times\bm v$ 及向量恒等式，可得
\begin{equation}
 \frac{D\bm\omega}{Dt}
 =(\nabla\bm v)\bm\omega-(\nabla\cdot\bm v)\bm\omega
 +\frac{\nabla\rho\times\nabla p}{\rho^2}.
\end{equation}
左侧是涡量沿流线的变化率。右侧第一项描述涡伸展与扭转（倾斜），第二项是可压缩性
引起的体积伸缩，第三项称斜压项，表示密度梯度和压强梯度不共线时产生涡量。对
正压流体 $p=p(\rho)$，两梯度平行，斜压项为零；不可压缩流中 $\nabla\rho=0$，它也
为零。

结合连续方程还可写成更紧凑的形式
\begin{equation}
 \frac{D}{Dt}\left(\frac{\bm\omega}{\rho}\right)
 =(\nabla\bm v)\frac{\bm\omega}{\rho}
 +\frac{\nabla\rho\times\nabla p}{\rho^3}.
\end{equation}
显式的压缩伸缩项在这里消失，被吸收到 $\bm\omega/\rho$ 的物质变化中。这说明在可
压缩流里，真正以类似不可压缩涡量方式演化的是比涡量 $\bm\omega/\rho$。涡量方程是
理解湍流机理和建立涡方法的基本工具。

\sourcecrop{137}{48}{75}{565}{430}{0.36}
\sourcecrop{138}{48}{430}{565}{710}{0.29}

\section{浅水（Saint--Venant）方程}

当竖直速度和自由表面竖直位移相对于水平尺度足够小时，水深 $h$ 与深度平均水平速度
$\bm v=[u,v]^\mathsf T$ 满足
\begin{align}
 \partial_t h+\nabla\cdot(h\bm v)&=0,\\
 \partial_t(h\bm v)+\nabla\cdot(h\bm v\otimes\bm v)
 +\nabla\left(\frac{gh^2}{2}\right)&=-gh\nabla z_0(x,y),
\end{align}
其中 $g\simeq9.8\,\mathrm{m/s^2}$，$z_0$ 描述底床高程。它可由含重力的三维欧拉
方程忽略竖向速度与竖向变化后，在底床至自由表面之间积分得到。竖向动量平衡给出
熟悉的静水压关系
\begin{equation}
 p=p_0+\rho g(h+z_0-z).
\end{equation}

若令等效密度为 $h$、等效压强为 $p=gh^2/2$，且底床平坦，则浅水方程就是删去能量
方程的欧拉系统。重力波速
\begin{equation}
 c=\sqrt{gh}=\sqrt{2p/h}
\end{equation}
对应均熵欧拉方程的声速，等效比热比为 $\gamma=2$。因此浅水方程可视为欧拉方程的
一个子集。

\subsection{一维浅水方程}

一维形式为
\begin{equation}
 h_t+(hu)_x=0,\qquad
 (hu)_t+\left(hu^2+\frac12gh^2\right)_x=-gh(z_0)_x.
\end{equation}
守恒变量为 $[h,hu]^\mathsf T$，两个特征速度为 $u\mp\sqrt{gh}$。左右特征变量分别
混合水深与速度扰动，其结构与一维均熵欧拉方程完全一致。

\subsection{二维浅水方程}

二维守恒变量为 $[h,hu,hv]^\mathsf T$，源项为
$[0,-gh(z_0)_x,-gh(z_0)_y]^\mathsf T$。沿法向的三个特征速度为
$q_n-\sqrt{gh},q_n,q_n+\sqrt{gh}$；中间波携带切向速度扰动。完整雅可比与特征向量
如下。

\sourcecrop{138}{48}{75}{565}{430}{0.36}
\sourcecrop{139}{48}{75}{565}{710}{0.61}
\sourcecrop{140}{48}{75}{565}{710}{0.61}
\sourcecrop{141}{48}{75}{565}{710}{0.61}

\subsection{三维浅水方程}

浅水近似本身已对竖直方向积分，因而通常只存在一维或二维水平形式。原作者以幽默的
方式指出，他无法写出真正的“三维浅水方程”，因为在这种意义下它并不存在。


% ---- 合并分片 3 ----
\chapter{Navier--Stokes 方程}

\section{Navier--Stokes 方程}

为了比 Euler 方程更真实地描述流动，需要把其中缺少的黏性应力张量 $\bm\tau$ 和热流矢量 $\bm q$ 补入；由此得到 Navier--Stokes 方程。其守恒形式为
\begin{align}
 \partial_t\rho+\nabla\cdot(\rho\bm v)&=0, \tag{4.1.1}\\
 \partial_t(\rho\bm v)+\nabla\cdot(\rho\bm v\otimes\bm v)&=\nabla\cdot\bm\sigma, \tag{4.1.2}\\
 \partial_t(\rho E)+\nabla\cdot(\rho\bm vH)&=\nabla\cdot(\bm\tau\bm v)-\nabla\cdot\bm q, \tag{4.1.3}
\end{align}
其中总表面应力张量由静压和黏性应力组成：
\begin{equation}
 \bm\sigma=-p\bm I+\bm\tau . \tag{4.1.4}
\end{equation}
$\bm I$ 为单位张量。相应的原始变量形式为
\begin{align}
 \frac{D\rho}{Dt}+\rho\nabla\cdot\bm v&=0, \tag{4.1.5}\\
 \rho\frac{D\bm v}{Dt}+\nabla p&=\nabla\cdot\bm\tau, \tag{4.1.6}\\
 \frac{Dp}{Dt}+\gamma p\nabla\cdot\bm v&=(\gamma-1)\left(\bm\tau:\nabla\bm v-\nabla\cdot\bm q\right). \tag{4.1.7}
\end{align}

把式（4.1.1）的散度项展开即可得到连续方程（4.1.5）；展开式（4.1.2）左端并代入连续方程，即可得到动量方程（4.1.6）。这里还用到了
\begin{equation}
 \nabla\cdot\bm\sigma=-\nabla\cdot(p\bm I)+\nabla\cdot\bm\tau=-\nabla p+\nabla\cdot\bm\tau. \tag{4.1.8}
\end{equation}
压力方程（4.1.7）的推导分三步：先用速度点乘动量方程，得到动能方程
\begin{equation}
 \rho\frac{D}{Dt}\left(\frac{v^2}{2}\right)+(\nabla p)\cdot\bm v
 =-\bm\tau:\nabla\bm v+\nabla\cdot(\bm\tau\bm v), \tag{4.1.9}
\end{equation}
再从总能方程中减去它，最后结合完全气体状态方程 $p=(\gamma-1)\rho e$ 与连续方程，把内能的时间导数换成压力的时间导数。守恒形式更适合数值离散：守恒性对于获得正确的激波传播速度尤其重要，而且直接离散守恒式更容易构造守恒格式。

\section{黏性应力与热流}

对于 Newton 流体（例如海平面附近的空气或水），黏性应力与应变率线性相关：
\begin{equation}
 \bm\tau=\lambda(\nabla\cdot\bm v)\bm I+
 \mu\left[\nabla\bm v+(\nabla\bm v)^{\mathsf T}\right]. \tag{4.2.1}
\end{equation}
其中 $\nabla\bm v+(\nabla\bm v)^{\mathsf T}$ 是应变率张量（亦称变形张量），$\mu$ 为动力黏度，$\lambda$ 为第二黏性系数。体积黏度定义为
\begin{equation}
 \mu_B=\lambda+\frac{2}{3}\mu. \tag{4.2.2}
\end{equation}
它通常很小，故常取 $\mu_B=0$，即 $\lambda=-2\mu/3$；这就是 Stokes 假设。

平均法向应力（也称平均压力）定义为 $P=-\sigma_{ii}/3$，代入本构关系可得
\begin{equation}
 P=p-\left(\lambda+\frac{2}{3}\mu\right)\nabla\cdot\bm v
   =p-\mu_B\nabla\cdot\bm v. \tag{4.2.4}
\end{equation}
这里采用负号，是因为作用在流体微元上的 $\sigma_{ii}$ 为负。$P$ 给出了与坐标无关的平均法向力，一般应当等同于热力学压力。热力学平衡下，作为状态量的压力不应依赖体积变化率 $\nabla\cdot\bm v$，因此令 $\mu_B=0$ 并得到 $P=p$ 是合理的。它毕竟只是假设，但在很广的工况范围内表现良好。黏度一般是温度的函数。

热流由 Fourier 定律给出：
\begin{equation}
 \bm q=-\kappa\nabla T, \tag{4.2.5}
\end{equation}
其中 $\kappa$ 为热导率，也随温度变化。负号表示热量沿温度降低的方向传递。Newton 黏性定律和 Fourier 导热定律具有相同的便利之处：黏性应力和热流都与原始变量的一阶导数呈线性关系。

\section{人工体积黏度}

尽管物理体积黏度通常取零，一些 CFD 方法会把正的人工体积黏度 $\mu_B^{\mathrm{artificial}}$ 重新引入控制方程：
\begin{equation}
 \bm\tau=\left(\mu_B^{\mathrm{artificial}}-\frac{2}{3}\mu\right)
 (\nabla\cdot\bm v)\bm I+
 \mu\left[\nabla\bm v+(\nabla\bm v)^{\mathsf T}\right]. \tag{4.3.1}
\end{equation}
文献 [124] 用它加速伪可压缩不可压 Navier--Stokes 求解器的收敛。伪可压缩形式会产生伪声波，稳态收敛速度取决于这些波经传播和阻尼被消除的快慢；人工体积黏度正是用来增强阻尼。对于稳态不可压流，$\nabla\cdot\bm v=0$，所以该项不改变稳态解，只改变数值迭代的收敛过程。这一思想也已推广到可压 Euler 和 Navier--Stokes 方程：把人工项写成与 $\nabla\cdot(\rho\bm v)$ 有关的形式，使其在可压稳态解上同样消失。方法虽面向稳态问题，但也可结合隐式时间推进：每个物理时间步内，用快速稳态求解器求解隐式残差方程，从而进行时间精确计算。

另一个例子是局部人工扩散率方法 [27]。为在不损害湍流结构分辨率的前提下捕捉激波，人工体积黏度只在激波附近局部激活，故称“局部”。该方法已用于采用甚高阶格式的可压缩激波--湍流相互作用计算。

令膨胀率 $\delta=\nabla\cdot\bm v$、涡量 $\bm\omega=\nabla\times\bm v$。把除以密度后的动量方程分别取散度和旋度，可见人工体积黏度对二者的作用：
\begin{align}
 \frac{\partial\delta}{\partial t}+\cdots &=\nabla\cdot\left[
 \frac{1}{\rho}\nabla(\mu_B^{\mathrm{artificial}}\delta)\right]+\cdots, \tag{4.3.2}\\
 \frac{\partial\bm\omega}{\partial t}+\cdots &=-\frac{1}{\rho^2}\nabla\rho\times
 \nabla(\mu_B^{\mathrm{artificial}}\delta)+\cdots. \tag{4.3.3}
\end{align}
第一式表明人工体积黏度扩散膨胀率，因而能增强收敛阻尼并帮助捕捉激波；第二式表明它原则上也可能产生涡量。前述收敛加速方法中，该项在稳态消失，故最终不影响涡量。在局部人工扩散率方法中，密度梯度和人工体积黏度梯度都主要沿激波法向，二者叉积预期很小，因此额外涡量也应很小。这一点很重要，因为准确预测穿越激波时的涡量生成对湍流计算可能至关重要。这类方法的共同思路是先把期望的数值性质嵌入控制方程，再对改进后的方程作高精度离散。

\section{黏度与热导率}

利用 Prandtl 数定义，可由黏度计算热导率：
\begin{equation}
 \Pr=\frac{c_p\mu}{\kappa}=\frac{\gamma R\mu}{\kappa(\gamma-1)},
 \qquad \kappa=\frac{\gamma R\mu}{\Pr(\gamma-1)}. \tag{4.4.1}
\end{equation}
因此热流可写成等价形式
\begin{equation}
 \bm q=-\kappa\nabla T
 =-\frac{\mu}{\Pr(\gamma-1)}\nabla(a^2)
 =-\frac{\gamma\mu}{\Pr}\nabla e
 =-\frac{\mu}{\Pr}\nabla h, \tag{4.4.2}
\end{equation}
这里使用了 $a^2=\gamma p/\rho=\gamma RT$、$e=RT/(\gamma-1)$ 和 $h=\gamma RT/(\gamma-1)$。海平面空气的 Prandtl 数在很宽的温度范围内近似恒定，$200\text{ K}\le T\le1000\text{ K}$ 时取 $\Pr=0.72$ 是有效近似。一般气体还可使用
\begin{equation}
 \Pr=\frac{4\gamma}{9\gamma-5}. \tag{4.4.3}
\end{equation}
于是计算热流所需的主要物性只剩 $\mu$。空气黏度可由 Sutherland 定律估计：
\begin{equation}
 \mu=\mu_0\frac{T_0+C}{T+C}\left(\frac{T}{T_0}\right)^{3/2},
 \quad C=110.5\text{ K},\quad T_0=273.1\text{ K},\quad
 \mu_0=1.716\times10^{-5}\text{ Pa}\cdot\text{s}. \tag{4.4.4}
\end{equation}
它在约 $170$--$1900\text{ K}$ 的范围内也是良好近似。

\section{其他能量方程}

从总能方程减去动能方程，得到内能方程
\begin{equation}
 \rho\frac{De}{Dt}=-p\nabla\cdot\bm v-\nabla\cdot\bm q+\Phi, \tag{4.5.1}
\end{equation}
其中耗散函数为
\begin{equation}
 \Phi=\bm\tau:\nabla\bm v
 =\mu[\nabla\bm v+(\nabla\bm v)^{\mathsf T}]:\nabla\bm v
 +\lambda(\nabla\cdot\bm v)^2. \tag{4.5.2}
\end{equation}
把应变率分解为无迹部分和体积膨胀部分，可将 $\Phi$ 写成平方和。完整的张量展开如下。
\sourcecrop{146}{45}{120}{567}{690}{0.30}
因此只要 $\mu\ge0$ 且 $\lambda+2\mu/3\ge0$，耗散函数就非负。Stokes 假设 $\lambda=-2\mu/3$ 恰好保证这一性质。推导时先利用对称张量与反对称张量的双点积为零，以及单位张量与反对称速度梯度的双点积为零；再把 $\nabla\bm v$ 分解成对称、反对称两部分，代入本构关系，最后加减 $2\mu(\nabla\cdot\bm v)^2/3$ 并重新因式分解。

能量方程还可等价地写成总焓、总能、内能或静焓形式：
\begin{align}
 \frac{\partial(\rho H)}{\partial t}-\frac{\partial p}{\partial t}
 +\nabla\cdot(\rho\bm vH)&=\nabla\cdot(\bm\tau\bm v)-\nabla\cdot\bm q, \tag{4.5.9}\\
 \rho\frac{DH}{Dt}-\frac{\partial p}{\partial t}&=\nabla\cdot(\bm\tau\bm v)-\nabla\cdot\bm q, \tag{4.5.10}\\
 \rho\frac{DE}{Dt}+\nabla\cdot(p\bm v)&=\nabla\cdot(\bm\tau\bm v)-\nabla\cdot\bm q, \tag{4.5.11}\\
 \rho\frac{De}{Dt}+p\nabla\cdot\bm v&=\Phi-\nabla\cdot\bm q, \tag{4.5.12}\\
 \rho\frac{Dh}{Dt}-\frac{Dp}{Dt}&=\Phi-\nabla\cdot\bm q. \tag{4.5.13}
\end{align}
也可使用熵方程（或熵不等式）
\begin{equation}
 \rho\frac{Ds}{Dt}+\nabla\cdot\left(\frac{\bm q}{T}\right)
 =\frac{\Phi}{T}+\frac{\kappa}{T^2}\nabla T\cdot\nabla T\ge0. \tag{4.5.14}
\end{equation}

\section{总压方程}

对稳态不可压流，容易推导总压 $p_t=p+\rho v^2/2$ 的控制方程。用速度点乘稳态动量方程，并利用矢量恒等式与连续方程，得到
\begin{equation}
 \nabla\cdot(p_t\bm v)=\nabla\cdot(\bm\tau\bm v)-\Phi,
 \qquad
 (\nabla p_t)\cdot\bm v=\nabla\cdot(\bm\tau\bm v)-\Phi. \tag{4.6.3--4.6.4}
\end{equation}
对于无黏流，右端为零，所以总压沿流线保持常数。对于黏性流，$-\Phi$ 只会耗散、降低总压；$\nabla\cdot(\bm\tau\bm v)$ 却可能在局部为正并使总压上升。若边界没有外力输入，它对全域的净贡献仍为零。文献 [55] 首先指出了这一现象，并在二维黏性驻点流和绕球三维黏性流中展示了局部总压升高。

\section{固体表面的零散度与黏性力}

稳态可压流连续方程可写成
\begin{equation}
 \nabla\cdot(\rho\bm v)=0,
 \qquad \rho\nabla\cdot\bm v+\nabla\rho\cdot\bm v=0. \tag{4.7.1--4.7.2}
\end{equation}
在固体表面施加无滑移条件 $\bm v=0$，立即得到 $\nabla\cdot\bm v=0$。因此可压性不影响稳态流动作用在固体上的黏性力，表面应力化为
\begin{equation}
 \bm\tau\big|_{\rm body}=\mu\left[\nabla\bm v+(\nabla\bm v)^{\mathsf T}\right]. \tag{4.7.4}
\end{equation}
设 $dS$ 为面元，$\bm n$ 为外法向单位矢量，$q_n=\bm v\cdot\bm n$。把 $\nabla q_n$ 分解到法向和两个切向方向。无滑移且无吸入/喷出意味着切向导数为零；又因为壁面速度散度为零、切向速度沿壁面恒为零，法向导数 $\partial q_n/\partial n$ 也为零。因此 $(\nabla\bm v)^{\mathsf T}\bm n=0$，有
\begin{equation}
 \bm\tau\bm n\,dS
 =\mu\left[\nabla\bm v-(\nabla\bm v)^{\mathsf T}\right]\bm n\,dS
 =\mu(\bm\omega\times\bm n)\,dS. \tag{4.7.13--4.7.14}
\end{equation}
也就是说，可通过对壁面涡量积分计算稳态不可压或可压流的黏性力；CFD 程序若能准确给出壁面涡量，就可据此计算黏性载荷。

\section{一维 Navier--Stokes 方程}

一维守恒式为
\begin{equation}
 \frac{\partial\bm U}{\partial t}+\frac{\partial\bm F}{\partial x}=0,
 \quad
 \bm U=\begin{bmatrix}\rho\\\rho u\\\rho E\end{bmatrix},\quad
 \bm F=\begin{bmatrix}\rho u\\\rho u^2+p-\tau_{xx}\\\rho uH-\tau_{xx}u+q_x\end{bmatrix}. \tag{4.8.1--4.8.2}
\end{equation}
其中
\begin{equation}
 \tau_{xx}=(2\mu+\lambda)u_x,\qquad
 q_x=-\kappa T_x=-\frac{\mu}{\Pr(\gamma-1)}(a^2)_x
 =-\frac{\gamma\mu}{\Pr(\gamma-1)}(p/\rho)_x. \tag{4.8.3--4.8.4}
\end{equation}
采用 Stokes 假设后，$\tau_{xx}=4\mu u_x/3$。耗散函数为 $\Phi=(2\mu+\lambda)u_x^2$，始终非负；不过守恒形式并不需要显式计算该耗散项。

\section{二维 Navier--Stokes 方程}

二维方程写成 $\bm U_t+\bm F_x+\bm G_y=0$。守恒变量、两个方向的含黏通量、应力分量和热流分量如下。
\sourcecrop{149}{42}{95}{568}{710}{0.46}
这里 $\bm U=(\rho,\rho u,\rho v,\rho E)^{\mathsf T}$；$\bm F$、$\bm G$ 分别为 $x$、$y$ 方向总通量。Stokes 假设将法向应力化简为
\begin{equation}
 \tau_{xx}=\frac{2\mu}{3}(2u_x-v_y),\qquad
 \tau_{yy}=\frac{2\mu}{3}(2v_y-u_x). \tag{4.9.10--4.9.11}
\end{equation}
二维耗散函数同样是速度梯度平方项之和并保持非负；在守恒方程的直接离散中无需单独出现。

\section{三维 Navier--Stokes 方程}

三维守恒形式为 $\bm U_t+\bm F_x+\bm G_y+\bm H_z=0$。完整的守恒变量、三个方向通量、六个独立应力分量以及三个热流分量如下。
\sourcecrop{150}{42}{90}{570}{710}{0.46}
\sourcecrop{151}{42}{90}{570}{710}{0.46}
其中 $\bm U=(\rho,\rho u,\rho v,\rho w,\rho E)^{\mathsf T}$，$\bm F,\bm G,\bm H$ 分别是 $x,y,z$ 方向通量；应力张量对称，即 $\tau_{ij}=\tau_{ji}$。采用 Stokes 假设后，各法向应力去掉平均体积膨胀部分；耗散函数仍为非负二次型。

\section{轴对称 Navier--Stokes 方程}

在圆柱坐标 $(r,\theta,z)$ 中考虑无旋流轴对称流，所有量与 $\theta$ 无关且 $u_\theta=0$。此时热流散度、黏性应力散度和黏性功的展开含有几何项 $1/r$。从圆柱坐标张量形式出发，可分离出二维散度部分和只由轴对称几何产生的源项。
\sourcecrop{152}{38}{90}{574}{710}{0.45}
\sourcecrop{153}{38}{90}{574}{710}{0.45}
\sourcecrop{154}{38}{90}{574}{710}{0.45}
将 $(z,r)$ 换记为 $(x,y)$、$(u_z,u_r)$ 换记为 $(u,v)$ 后，轴对称守恒式可写为
\begin{equation}
 (y\bm U)_t+(y\bm F)_x+(y\bm G)_y
 =(y\bm F_v)_x+(y\bm G_v)_y+\bm S. \tag{4.11.13}
\end{equation}
其中 $\bm U,\bm F,\bm G$ 与二维可压方程的守恒量和无黏通量相同，$\bm F_v,\bm G_v$ 是黏性通量，$\bm S$ 的径向动量分量包含压力以及 $v/y$ 引起的几何黏性源项。必须注意：应力张量中的速度散度多出 $v/y$，应力散度的第二分量也有相应附加项。

按照文献 [173]，引入逻辑开关 $\omega$：二维取 $\omega=0$，轴对称取 $\omega=1$，并定义 $\omega_a=(1-\omega)+\omega y$，则两种形式可统一为
\begin{equation}
 (\omega_a\bm U)_t+(\omega_a\bm F)_x+(\omega_a\bm G)_y
 =(\omega_a\bm F_v)_x+(\omega_a\bm G_v)_y+\omega\bm S. \tag{4.11.16}
\end{equation}
因此，一个二维 Navier--Stokes 程序只需很少改动即可处理轴对称流，而轴对称计算本质上已经包含三维几何效应。

\section{三维法向黏性通量及其 Jacobian}

CFD 中常把无黏通量与黏性通量分开处理。三个坐标方向的黏性通量 $\bm F_v,\bm G_v,\bm H_v$ 的质量分量为零，动量分量为负的相应应力列，能量分量为热流减去黏性功。设
\begin{equation}
 \tau_{vx}=\tau_{xx}u+\tau_{xy}v+\tau_{xz}w,\qquad
 \tau_{vy}=\tau_{yx}u+\tau_{yy}v+\tau_{yz}w,\qquad
 \tau_{vz}=\tau_{zx}u+\tau_{zy}v+\tau_{zz}w. \tag{4.12.2--4.12.4}
\end{equation}
与无黏通量相比，这一结构相当简洁。

控制体边界上通常需要计算沿单位法向 $\bm n=(n_x,n_y,n_z)^{\mathsf T}$ 投影的黏性通量
\begin{equation}
 \bm F_n^v=\bm F_vn_x+\bm G_vn_y+\bm H_vn_z
 =\begin{bmatrix}0&-\tau_{nx}&-\tau_{ny}&-\tau_{nz}&-\tau_{nv}+q_n\end{bmatrix}^{\mathsf T}, \tag{4.12.5}
\end{equation}
其中 $\tau_{ni}=\tau_{xi}n_x+\tau_{yi}n_y+\tau_{zi}n_z$，$\tau_{nv}=\tau_{nx}u+\tau_{ny}v+\tau_{nz}w$，$q_n=q_xn_x+q_yn_y+q_zn_z$。横线表示这些量已在某个界面平均状态上求值；有限体积法中，$\bm n$ 就是控制体面的法向。

隐式算法还需要法向黏性通量对状态的 Jacobian。其链式求导、法向应力和温度导数的完整表达如下。
\sourcecrop{155}{38}{88}{574}{710}{0.44}
\sourcecrop{156}{38}{88}{574}{710}{0.44}
\sourcecrop{157}{38}{88}{574}{710}{0.44}
这一结果是通用的，具体矩阵取决于带横线的界面量及其梯度如何离散。例如 $\bar T=(T_L+T_R)/2$ 时，$\partial\bar T/\partial T_L=1/2$；若 $u_n=(u_R-u_L)/\Delta s$，则 $\partial u_n/\partial u_L=-1/\Delta s$。非结构网格中，切向导数常导致非紧致模板，使 Jacobian 过于庞大；构造隐式驱动矩阵时经常忽略这些项。这样得到的 Jacobian 本身不一致，但只要右端残差仍由一致格式构造，收敛解仍保持正确精度，只是收敛速度可能不是最优。

\section{黏性项的展开形式}

在分析或制造解源项计算中，把黏性项完全展开为原始变量导数很方便。下列原式给出三个动量黏性项、能量黏性功与热传导项，并进一步展开 $\mu(T)$ 和 $\lambda(T)$ 的空间导数。
\sourcecrop{157}{38}{88}{574}{710}{0.43}
\sourcecrop{158}{38}{88}{574}{710}{0.43}
\sourcecrop{159}{38}{88}{574}{710}{0.43}
$\lambda$ 的导数可由 Stokes 假设 $\lambda=-2\mu/3$ 直接得到；温度导数可由密度和压力导数换算。这种展开对制造解法产生的解析源项尤其有用。需要注意，黏度和温度的导数形式会随无量纲化选择而改变。

\section{无量纲化}

\subsection{无量纲 Navier--Stokes 方程}

除 Euler 方程已有的 $\rho^*=\rho/\rho_r$、$p^*=p/(\rho_rV_r^2)$、$e^*=e/V_r^2$、$\bm v^*=\bm v/V_r$、$\bm x^*=\bm x/L$、$t^*=t/(L/V_r)$ 外，还需定义
\begin{equation}
 T^*=T/T_r,\qquad \mu^*=\mu/\mu_r. \tag{4.14.2}
\end{equation}
代入后，质量、动量和总能方程的外形与有量纲方程完全相同；差别集中在黏性应力、热流和状态方程中的无量纲系数。完整比例关系如下。
\sourcecrop{159}{38}{90}{574}{706}{0.41}
由于压力与内能的标度相容，热完全气体关系 $p^*=(\gamma-1)\rho^*e^*$ 对任意参考状态均保持原形；但 $p=\rho RT$ 的温度形式以及黏性、热通量的具体系数取决于 $\rho_r,V_r,T_r,\mu_r$ 的选取。适当定义缩放黏度后，控制方程仍几乎与有量纲形式一致。

\subsection{以自由来流量为参考}

取 $\rho_r=\rho_\infty$、$V_r=V_\infty$、$T_r=T_\infty$、$\mu_r=\mu_\infty$。状态方程、应力、热流及缩放黏度为
\begin{align}
 p^*&=\frac{\rho^*T^*}{\gamma M_\infty^2}, \tag{4.14.11}\\
 \bm\tau^*&=\widehat\mu^*\left[-\frac23(\nabla^*\cdot\bm v^*)\bm I+
 \nabla^*\bm v^*+(\nabla^*\bm v^*)^{\mathsf T}\right], \tag{4.14.12}\\
 \bm q^*&=-\frac{\gamma\widehat\mu^*}{(\gamma-1)\Pr}\nabla^*T^*,\qquad
 \widehat\mu^*=\frac{\mu^*}{\mathrm{Re}_\infty},\qquad
 \mathrm{Re}_\infty=\frac{\rho_\infty V_\infty L}{\mu_\infty}. \tag{4.14.13--4.14.14}
\end{align}
自由来流无量纲值为 $\rho_\infty^*=T_\infty^*=\mu_\infty^*=1$、$p_\infty^*=1/(\gamma M_\infty^2)$、$\bm v_\infty^*=\bm n_\infty$。无量纲 Sutherland 定律及其导数见原式。
\sourcecrop{160}{40}{92}{572}{610}{0.31}
设置自由来流入口时，除了 Mach 数和流向，还必须给定 $T_\infty$、$\mathrm{Re}_\infty$ 和 $\Pr$。Reynolds 数必须与参考长度 $L$ 一致：无量纲网格中长度为 $1$ 的物体尺度才对应所输入的 $\mathrm{Re}_\infty$。例如目标圆柱直径 Reynolds 数为 $100$ 时，若网格圆柱直径为 $1$，取 $\mathrm{Re}_\infty=100$；直径为 $2$ 时取 $50$；直径为 $0.1$ 时取 $1000$。

\subsection{以自由来流声速为参考}

也可用 $V_r=a_\infty$。此时 $p^*=\rho^*T^*/\gamma$，自由来流压力不再依赖 $M_\infty$；缩放黏度为 $\widehat\mu^*=M_\infty\mu^*/\mathrm{Re}_\infty$，而 $\bm v_\infty^*=M_\infty\bm n_\infty$。相应的完整无量纲式如下。
\sourcecrop{161}{40}{92}{572}{706}{0.41}
无量纲黏度仍由上一小节的 Sutherland 关系给出。入口同样需要指定 $M_\infty$、$\bm n_\infty$、$T_\infty$、$\mathrm{Re}_\infty$ 和 $\Pr$。由于来流压力不依赖 Mach 数，这套标度在可压缩流计算中使用非常广泛。

\subsection{以驻点量为参考}

第三种选择使用驻点密度、声速和温度作为参考。自由来流量与驻点量之间通过等熵关系相连，缩放黏度含有质量通量因子 $m_\infty=\rho_\infty^*a_\infty^*$。完整定义与来流状态如下。
\sourcecrop{161}{40}{90}{572}{450}{0.25}
\sourcecrop{162}{40}{90}{572}{706}{0.40}
无量纲 Sutherland 定律仍与前述相同；因此入口条件仍须在 $\mathrm{Re}_\infty$ 和 $\Pr$ 之外给定 $M_\infty$、来流方向及 $T_\infty$。

\section{简化 Navier--Stokes 方程}

薄层、抛物化以及锥形 Navier--Stokes 方程都是完整方程的近似形式，因而比求解完整 Navier--Stokes 方程更简单、成本更低，详见文献 [51,155]。

\section{Navier--Stokes 方程的拟线性形式}

为便于分析和格式构造，常把方程写成
\begin{equation}
 \bm V_t+\bm A^v\bm V_x+\bm B^v\bm V_y+\bm C^v\bm V_z
 =\bm D^v\bm V_{xx}+\bm E^v\bm V_{yy}+\bm F^v\bm V_{zz}
 +\bm G^v\bm V_{xy}+\bm H^v\bm V_{yz}+\bm I^v\bm V_{zx}. \tag{4.16.1}
\end{equation}
上标 $v$ 仅说明系数矩阵与变量组 $\bm V$ 相关。得到这一形式需要假定 $\mu$、$\kappa$ 为常数，并忽略导数乘积形成的非线性项，因此它不再是完整系统。不过判定偏微分方程类型并不需要保留所有项；此拟线性系统仍保留 Navier--Stokes 方程的本质特征，标量平流--扩散方程或黏性 Burgers 方程正是它的模型方程。对称化的拟线性系统也已成功用于收敛加速技术。

能量方程中的耗散 $\bm\tau:\nabla\bm v$ 是导数乘积，必须略去；热传导展开后也有两个梯度的乘积项，应略去。利用 $T=p/(R\rho)$，保留线性二阶导数后，压力方程成为
\begin{equation}
 \frac{Dp}{Dt}+\gamma p\nabla\cdot\bm v
 =\frac{\mu\gamma}{\rho\Pr}\nabla^2p-
 \frac{\mu\gamma p}{\rho^2\Pr}\nabla^2\rho. \tag{4.16.8}
\end{equation}
在常黏度和 Stokes 假设下，动量黏性项化为
\begin{equation}
 \nabla\cdot\bm\tau=\frac13\mu\nabla(\nabla\cdot\bm v)+\mu\nabla^2\bm v. \tag{4.16.9}
\end{equation}
故线性化系统由连续方程、上述动量方程和压力方程组成。令 $\nu=\mu/\rho$、$a^2=\gamma p/\rho$ 后，可直接写出一维、二维和三维的矩阵形式。

\subsection{一维拟线性方程}

以原始变量 $\bm V=(\rho,u,p)^{\mathsf T}$ 或熵变量写出的系数矩阵如下页原式所示。对流矩阵包含特征速度 $u,u\pm a$，扩散矩阵只作用于速度和热力学变量。
\sourcecrop{164}{38}{86}{574}{708}{0.45}
\sourcecrop{165}{38}{86}{574}{708}{0.45}
这两组变量只是同一线性化系统的不同表示，均保留声学传播和黏性/热扩散的主导结构。

\subsection{二维拟线性方程}

二维变量取 $\bm V=(\rho,u,v,p)^{\mathsf T}$。$x$、$y$ 对流 Jacobian 以及纯二阶、混合二阶导数矩阵完整列于下式。
\sourcecrop{166}{38}{86}{574}{708}{0.45}
\sourcecrop{167}{38}{86}{574}{708}{0.45}
矩阵中 $\nu$ 为运动黏度；压力扩散还含 $a^2/\Pr$，混合导数来自 $\nabla(\nabla\cdot\bm v)$。

\subsection{三维拟线性方程}

三维变量为 $\bm V=(\rho,u,v,w,p)^{\mathsf T}$。三个一阶对流矩阵、三个纯二阶扩散矩阵和三个混合二阶矩阵如下。
\sourcecrop{168}{38}{86}{574}{708}{0.45}
\sourcecrop{169}{38}{86}{574}{708}{0.45}
与低维情形一样，这些矩阵针对常物性且忽略导数乘积的线性化能量方程，但保留了完整系统的双曲--抛物耦合特征。

\section{不可压 Navier--Stokes 方程}

假定 $\rho$ 为常数（或 Mach 数很小），得到不可压系统。守恒形式为
\begin{align}
 \nabla\cdot(\rho\bm v)&=0, \tag{4.17.1}\\
 \partial_t(\rho\bm v)+\nabla\cdot(\rho\bm v\otimes\bm v)+\nabla p-\nabla\cdot\bm\tau&=0, \tag{4.17.2}\\
 \rho c_v\frac{DT}{Dt}-\nabla\cdot(\kappa\nabla T)-\Phi&=0. \tag{4.17.3}
\end{align}
其中 $\bm\tau=\mu[\nabla\bm v+(\nabla\bm v)^{\mathsf T}]$。若 $\mu$ 为常数，利用 $\nabla\cdot\bm v=0$，有 $\nabla\cdot\bm\tau=\mu\nabla^2\bm v$；若 $\kappa$ 也为常数，动量扩散和热扩散都成为 Laplace 算子。温度方程与质量、动量方程完全解耦：可先求速度场，再在需要时单独求给定速度场上的温度。第 4.21 节的伪可压缩方法正是针对这一系统。

\section{二维不可压 Navier--Stokes 方程}

二维系统可以使用守恒、原始变量或 Laplace 黏性项形式。连续方程为 $u_x+v_y=0$；常物性原始形式为
\begin{align}
 u_t+uu_x+vu_y&=-p_x/\rho+\nu(u_{xx}+u_{yy}), \tag{4.18.11}\\
 v_t+uv_x+vv_y&=-p_y/\rho+\nu(v_{xx}+v_{yy}), \tag{4.18.12}\\
 \rho c_v(T_t+uT_x+vT_y)&=\kappa(T_{xx}+T_{yy})+\Phi. \tag{4.18.13}
\end{align}
守恒形式以及变黏度时未化简的分量式如下。
\sourcecrop{170}{38}{88}{574}{708}{0.42}
\sourcecrop{171}{38}{380}{574}{708}{0.22}

\section{三维不可压 Navier--Stokes 方程}

三维系统同样可写成守恒、原始变量和常物性 Laplace 形式；连续方程为 $u_x+v_y+w_z=0$。完整分量式如下。
\sourcecrop{171}{38}{88}{574}{520}{0.30}
\sourcecrop{172}{38}{88}{574}{708}{0.44}
能量方程仍是给定速度场上的独立温度输运方程；黏性耗散 $\Phi$ 作为热源出现。

\section{轴对称不可压 Navier--Stokes 方程}

轴对称不可压方程比其可压版本简单得多，因为 $\rho$、$\mu$ 和 $\nu=\mu/\rho$ 均为常数。对无旋流情形，$u_\theta=0$，仅保留径向和轴向动量方程；黏性项使用圆柱坐标下的矢量 Laplacian。将 $(r,z)$ 换成 $(y,x)$、$(u_r,u_z)$ 换成 $(v,u)$ 并乘以 $y$ 后，可写成与二维方程相似的守恒形式
\begin{equation}
 (y\bm U)_t+(y\bm F)_x+(y\bm G)_y=(y\bm F_v)_x+(y\bm G_v)_y+\bm S. \tag{4.20.9}
\end{equation}
详细分量和源项如下。
\sourcecrop{172}{38}{88}{574}{708}{0.42}
\sourcecrop{173}{38}{88}{574}{708}{0.42}
同样可用 $\omega=0$ 表示二维、$\omega=1$ 表示轴对称，以 $\omega_a=(1-\omega)+\omega y$ 统一两种方程。轴对称流的流函数由连续方程直接得到：
\begin{equation}
 u=\frac1y\frac{\partial\psi}{\partial y},\qquad
 v=-\frac1y\frac{\partial\psi}{\partial x}, \tag{4.20.14}
\end{equation}
它自动满足连续方程。

\section{伪可压缩 Navier--Stokes 方程}

引入人工密度 $\rho^*=p/a^{*2}$，并把其伪时间导数加入连续方程：
\begin{equation}
 \frac{\partial P}{\partial t^*}+\nabla\cdot(a^{*2}\bm v)=0,
 \qquad P=p/\rho. \tag{4.21.3}
\end{equation}
再与不可压动量方程联立，得到伪可压缩系统
\begin{align}
 P_{t^*}+\nabla\cdot(a^{*2}\bm v)&=0,\\
 \bm v_{t^*}+\nabla\cdot(\bm v\otimes\bm v+P\bm I)&=\nabla\cdot\bm\tau,
 \quad \bm\tau=\nu[\nabla\bm v+(\nabla\bm v)^{\mathsf T}]. \tag{4.21.4--4.21.6}
\end{align}
动量式中的时间导数本来是真实时间导数，但与伪连续方程共同推进时不再代表真实瞬态，故记作 $t^*$。三维 Cartesian 坐标下的守恒通量如下。
\sourcecrop{174}{38}{88}{574}{708}{0.43}
若黏度为常数，黏性项可简化为 $\nu\nabla^2\bm v$，各方向黏性通量仅由相应速度导数组成。

\section{涡量输运方程与流函数}

二维不可压动量方程取旋度，得到
\begin{equation}
 \omega_t+u\omega_x+v\omega_y=\nu(\omega_{xx}+\omega_{yy}),
 \qquad \omega=v_x-u_y. \tag{4.22.1}
\end{equation}
引入 $u=\psi_y$、$v=-\psi_x$ 后，连续方程自动满足，且流函数服从 Poisson 方程
\begin{equation}
 \psi_{xx}+\psi_{yy}=-\omega. \tag{4.22.3}
\end{equation}
因此只需联立求解 $\omega$ 的输运方程和 $\psi$ 的 Poisson 方程即可确定速度场。随后压力可由动量方程取散度所得的 Poisson 方程单独求解：
\begin{equation}
 \nabla^2p=2\rho(u_xv_y-u_yv_x)=2\rho[\psi_{xx}\psi_{yy}-(\psi_{xy})^2]. \tag{4.22.4}
\end{equation}
这就是经典涡量--流函数方法的基本方程；它把连续方程和两条动量方程的三方程联立问题，改为两个方程联立再加一个独立压力方程。

\section{涡量输运方程与矢势}

三维不可压速度可作 Helmholtz 分解
\begin{equation}
 \bm v=\nabla\phi+\nabla\times\bm\Psi. \tag{4.23.1}
\end{equation}
第一项无旋，第二项无散。代入连续方程及 $\bm\omega=\nabla\times\bm v$，得到
\begin{equation}
 \nabla^2\phi=0,\qquad \nabla^2\bm\Psi=-\bm\omega,\qquad
 \frac{D\bm\omega}{Dt}=(\nabla\bm v)\bm\omega+\nu\nabla^2\bm\omega. \tag{4.23.2--4.23.4}
\end{equation}
求解这七个标量方程可确定速度场。随后由
\begin{equation}
 \nabla^2p=-\nabla\cdot[\nabla\cdot(\rho\bm v\otimes\bm v)] \tag{4.23.5}
\end{equation}
求压力；利用连续方程还可把右端写成速度梯度的二次型。原书列出的 Cartesian 展开如下。
\sourcecrop{176}{40}{310}{572}{706}{0.25}
这些方程也是涡方法的基础。

\section{流函数}

流函数既适合生成流场，也适合显示流线。二维只需一个标量流函数即可得到两个速度分量；三维需要两个流函数，但仍比三分量涡量紧凑。

\subsection{三维流函数}

取 $\bm\Psi=\psi\nabla\chi$、$\phi=0$，则
\begin{equation}
 \bm v=\nabla\times(\psi\nabla\chi)=\nabla\psi\times\nabla\chi. \tag{4.24.1}
\end{equation}
速度同时切于 $\psi=\text{常数}$ 和 $\chi=\text{常数}$ 两族曲面，所以两族曲面的交线就是流线，$\psi$、$\chi$ 可视为三维流函数。由于 $\nabla\cdot(\nabla\times\bm\Psi)=0$，连续方程自动满足；许多熟知的二维流函数都可由此特化得到。

\subsection{Cartesian 坐标（二维）}

令 $\psi=\psi(x,y)$、$\chi=z$，则
\begin{equation}
 u=\psi_y,\qquad v=-\psi_x,\qquad
 \psi_{xx}+\psi_{yy}=-\omega_z. \tag{4.24.2--4.24.3}
\end{equation}
无旋流中 $\omega_z=0$，故流函数满足 Laplace 方程。

\subsection{极坐标（二维）}

在圆柱坐标中令 $\psi=\psi(r,\theta)$、$\chi=z$，得到
\begin{equation}
 u_r=\frac1r\psi_\theta,\qquad u_\theta=-\psi_r,\qquad
 \frac1r\frac{\partial}{\partial r}(r\psi_r)+\frac1{r^2}\psi_{\theta\theta}=-\omega_z. \tag{4.24.5--4.24.6}
\end{equation}
无旋流时右端为零，仍是极坐标 Laplace 方程。

\subsection{轴对称圆柱坐标}

无旋流轴对称条件下取 $\psi=\psi(r,z)$、$\chi=\theta$，得
\begin{equation}
 u_r=-\frac1r\psi_z,\qquad u_z=\frac1r\psi_r,\qquad
 \frac1r\left(\psi_{zz}+\psi_{rr}-\frac1r\psi_r\right)=-\omega_\theta. \tag{4.24.8--4.24.10}
\end{equation}
左端并不是圆柱坐标标量 Laplace 算子，因此轴对称无旋流的流函数并不满足 Laplace 方程。

\subsection{轴对称球坐标}

在球坐标 $(r,\theta,\varphi)$ 中取 $\psi=\psi(r,\varphi)$、$\chi=-\theta$，可得
\begin{equation}
 u_r=\frac{1}{r^2\sin\varphi}\psi_\varphi,\qquad
 u_\varphi=-\frac{1}{r\sin\varphi}\psi_r. \tag{4.24.12}
\end{equation}
对应的流函数方程同样不是标量 Laplace 方程，因而轴对称球坐标无旋流的流函数也不满足 Laplace 方程。

\subsection{轴对称（x, y）坐标}

圆柱和球坐标结果可统一写成
\begin{equation}
 u=\frac1y\psi_y,\qquad v=-\frac1y\psi_x,\qquad
 \psi_{xx}+\psi_{yy}-\frac1y\psi_y=-y\omega_\theta. \tag{4.24.14--4.24.15}
\end{equation}
无旋流还可定义速度势 $u=\phi_x$、$v=\phi_y$。$\phi$ 与 $\psi$ 的关系类似 Cauchy--Riemann 系统，但因系数 $1/y$ 不是常数而并不完全相同。$\psi$ 不满足 Laplace 方程，速度势 $\phi$ 却满足。以 $(x,y)$ 平面表示轴对称流，能直观比较二维流与轴对称流的差别。

\subsection{体积流量与流函数}

穿过截面 $S$ 的体积流量为 $Q=\iint_S\bm v\cdot\bm n\,dS$。由 Stokes 定理以及 $\bm\Psi=\psi\nabla\chi$，可把它化为边界线积分。若截面由 $\psi=\psi_1,\psi_2$ 和 $\chi=\chi_1,\chi_2$ 四个流面围成，则
\begin{equation}
 Q=(\psi_2-\psi_1)(\chi_2-\chi_1). \tag{4.24.21}
\end{equation}
因此体积流量就是两流函数差值的乘积。流函数适用于无旋、有旋、无黏和黏性流；可压稳态流中还可定义 $\bm v=\rho^{-1}\nabla\times\bm\Psi$，使 $\nabla\cdot(\rho\bm v)=0$ 自动成立。

\section{Stokes 方程}

忽略不可压 Navier--Stokes 方程中的非线性对流项，可得到线性的 Stokes 方程；它适用于小 Reynolds 数的 Stokes 流或蠕动流：
\begin{equation}
 \nabla\cdot\bm v=0,\qquad
 \rho\bm v_t+\nabla p-\mu\nabla^2\bm v=0. \tag{4.25.1--4.25.2}
\end{equation}
对动量方程取散度并使用连续方程，可得 $\nabla^2p=0$；有对流项时，压力满足的则是 Poisson 方程。又因 $\nabla^2\bm v=-\nabla\times\bm\omega$，方程可写为关于速度、压力和涡量的一阶系统
\begin{equation}
 \nabla\cdot\bm v=0,\qquad \nabla\times\bm v=\bm\omega,\qquad
 \rho\bm v_t+\nabla p+\mu\nabla\times\bm\omega=0. \tag{4.25.8--4.25.10}
\end{equation}
一阶系统通常比二阶导数更易离散。动量方程取旋度还表明
\begin{equation}
 \bm\omega_t=\nu\nabla^2\bm\omega, \tag{4.25.12}
\end{equation}
即涡量服从扩散方程；稳态时进一步退化为 Laplace 方程。

\section{二维 Stokes 方程}

二维 Stokes 方程的速度--压力形式及压力--涡量一阶形式如下。
\sourcecrop{180}{38}{88}{574}{708}{0.42}
稳态时，一阶形式本质上是两组通过 $\omega_z$ 耦合的 Cauchy--Riemann 型方程，可借助复变函数理论构造精确解。不过二维 Stokes 方程无法同时满足远场自由来流条件和物体表面的无滑移条件，因而不能准确描述二维绕体流动；这就是 Stokes 悖论。保留线性化对流项中的主导项可得到 Oseen 近似并解除这一矛盾。

\section{三维 Stokes 方程}

三维情形不存在 Stokes 悖论，可以正确描述绕体流动；涡量也从二维的单一分量变为三分量矢量。速度--压力和压力--涡量的一阶系统如下。
\sourcecrop{181}{38}{88}{574}{708}{0.42}
三维一阶系统不再显现稳态 Cauchy--Riemann 结构，不能直接套用复变函数理论，精确解需采用其他方法。

\section{轴对称 Stokes 方程}

从轴对称不可压 Navier--Stokes 方程中删去对流项即可得到轴对称 Stokes 方程。其守恒形式与压力--涡量一阶形式如下。
\sourcecrop{181}{38}{86}{574}{430}{0.23}
\sourcecrop{182}{38}{300}{574}{708}{0.27}
其中 $\omega$ 是唯一非零的 $\theta$ 向涡量。轴对称系统虽与二维系统很相似，但由于压力源项存在，稳态动量方程不能完全化为 Cauchy--Riemann 系统。连续方程仍可像第 4.20 节那样定义流函数，这对推导精确解有帮助。

\section{边界层方程}

二维不可压边界层方程为
\begin{equation}
 u_x+v_y=0,\qquad uu_x+vu_y=UU_x+\nu u_{yy}, \tag{4.29.1--4.29.2}
\end{equation}
其中 $U=U(x)$ 是边界层外缘速度；平板绕流中它为常数。一种数值解法是：在 $u(x,0)=v(x,0)=0$ 下沿 $y$ 积分连续方程得到 $v$，再以该 $v$ 求解动量方程，反复迭代至收敛。此时动量方程的行为类似黏性 Burgers 方程，因此可用标量格式求解这个非线性系统。该系统还有精确解（见第 7.15.9 节），可直接衡量数值解精度。

\chapter{湍流方程}

\section{平均与滤波}

把任一物理量表示为平均量与脉动量之和，例如
\begin{equation}
 \bm v=\overline{\bm v}+\bm v',\quad
 \rho=\overline\rho+\rho',\quad p=\overline p+p',\quad
 h=\overline h+h',\quad T=\overline T+T',\quad H=\overline H+H'. \tag{5.1.1}
\end{equation}
这称为 Reynolds 分解。必须注意，脉动相对于平均值未必很小；而黏度、热导率和比热的脉动通常较小，因而常被忽略。

湍流计算通常对 Navier--Stokes 方程取平均并求解平均量。平均操作会产生包含脉动、无法仅由平均量计算的附加项。即使为这些项继续推导输运方程，又会产生新的未知相关项，层级会无限延伸。因此必须在某一级截断，并用平均量对未闭合项进行建模。平均或滤波是湍流计算的基础；不同定义具有不同代数性质，也决定平均方程中究竟哪些项需要建模。

\subsection{Reynolds 平均}

Reynolds 平均包括系综平均、时间平均和空间平均。对稳态均匀湍流，三者彼此等价，这就是遍历假设。基于 Reynolds 平均 Navier--Stokes 方程的模拟称为 RANS。

系综平均是 $N$ 次独立实验测量的平均：
\begin{equation}
 \overline f(\bm x,t)=\lim_{N\to\infty}\frac1N\sum_{i=1}^{N}f_i(\bm x,t). \tag{5.1.2}
\end{equation}
实际进行无限次实验并不可行，但这一平均可在理论上用于写出平均控制方程。

时间平均定义为
\begin{equation}
 \overline f(\bm x,t)=\frac1{\mathcal T}\int_t^{t+\mathcal T}f(\bm x,t')\,dt'. \tag{5.1.3}
\end{equation}
平均窗口 $\mathcal T$ 应大于脉动时间尺度而小于平均流变化时间尺度。真实流动中两种尺度可能相近，满足条件的窗口未必存在，但仍可在该假设下形式推导平均方程。稳态时可令 $\mathcal T\to\infty$。

均匀湍流的空间平均定义为
\begin{equation}
 \overline f(t)=\lim_{V\to\infty}\frac1V\iiint_V f(\bm x,t)\,dV. \tag{5.1.5}
\end{equation}
虽然多数实际平均流并非全域均匀，但从足够局部的尺度观察，流动可能近似均匀，空间平均因而仍有意义。

Reynolds 平均满足
\begin{equation}
 \overline{\overline f}=\overline f,\qquad
 \overline{f'}=0,\qquad
 \overline{\overline f\,g'}=0,\qquad
 \overline{\overline f\,\overline g}=\overline f\,\overline g,\qquad
 \overline{f+g}=\overline f+\overline g. \tag{5.1.6--5.1.10}
\end{equation}
这些性质直观且便于记忆。

\subsection{Favre 平均}

Favre 平均以波浪号表示，定义为质量加权平均
\begin{equation}
 \widetilde f=\frac{\overline{\rho f}}{\overline\rho}. \tag{5.1.11}
\end{equation}
它可显著简化可压缩 Navier--Stokes 平均方程，比直接使用 Reynolds 平均得到的形式简洁得多。通常只对速度和热力学变量作 Favre 平均，而密度、压力仍作 Reynolds 平均，例如
\begin{equation}
 \widetilde{\bm v}=\frac{\overline{\rho\bm v}}{\overline\rho},\quad
 \widetilde h=\frac{\overline{\rho h}}{\overline\rho},\quad
 \widetilde T=\frac{\overline{\rho T}}{\overline\rho},\quad
 \widetilde H=\frac{\overline{\rho H}}{\overline\rho}. \tag{5.1.12}
\end{equation}
Favre 脉动记作双撇号 $f''=f-\widetilde f$。它与 Reynolds 脉动的关系及主要性质为
\begin{equation}
 f''=f'-\frac{\overline{\rho'f'}}{\overline\rho},\qquad
 \overline{f''}=-\frac{\overline{\rho'f'}}{\overline\rho}\ne0,\qquad
 \overline{\rho f''}=0,\qquad
 \overline{\rho f}=\overline\rho\,\widetilde f. \tag{5.1.14--5.1.18}
\end{equation}
这些性质虽比 Reynolds 平均复杂，却是可压缩湍流推导的关键。

\subsection{空间滤波}

空间滤波是局部物理空间（或局部波数空间）平均。典型的体积平均盒式滤波器（顶帽滤波器）为
\begin{equation}
 \overline f(\bm x,t)=\frac1{\Delta^3}
 \int_{x-\Delta x/2}^{x+\Delta x/2}
 \int_{y-\Delta y/2}^{y+\Delta y/2}
 \int_{z-\Delta z/2}^{z+\Delta z/2}f(\bm\xi,t)\,d\xi\,d\eta\,d\zeta, \tag{5.1.19}
\end{equation}
其中 $\Delta=(\Delta x\Delta y\Delta z)^{1/3}$。它可理解为计算单元内的积分；$\overline f$ 是网格能够解析的光滑低频分量，称为格尺度（GS）或可解析尺度，其余高频分量 $f'=f-\overline f$ 称为亚格子尺度（SGS）。严格说滤波宽度只需不小于网格尺度，不必与网格尺度相同，因而 GS/SGS 的命名有时会造成混淆。

滤波也可看作把流动分成大涡与小涡。滤波宽度越小，小尺度涡的性质越趋于普适，这与 Kolmogorov 的 K41 理论有关；若滤波尺度足够小，就有希望建立适用于多类流动的通用小尺度模型，但代价是极细网格和高昂计算成本。

一般滤波写成卷积
\begin{equation}
 \overline f(\bm x,t)=\iiint G(\bm x-\bm\xi)f(\bm\xi,t)\,d\bm\xi, \tag{5.1.22}
\end{equation}
$G$ 为滤波函数。盒式滤波器在 $|x_i-\xi_i|<\Delta x_i/2$ 内取 $1/\Delta^3$，其余为零；也可使用 Gaussian 滤波器或 Fourier 截断滤波器。

大涡模拟（LES）求解滤波后的控制方程，并选择足够小的滤波宽度，使亚格子尺度近似均匀。若网格细到能解析所有尺度，就不需要湍流模型，这称为直接数值模拟（DNS）；它概念最简单，却比 LES 更昂贵，实际高 Reynolds 数流动往往需要难以承受的网格。分离涡模拟（DES）在近壁、LES 代价过高的区域使用 RANS，在大涡占优的其余区域使用 LES。

空间滤波一般不具有 Reynolds 平均的幂等性：$\overline{\overline f}\ne\overline f$，并且 $\overline{f'}\ne0$、$\overline{\overline f\,g'}\ne0$。因此滤波方程与 Reynolds 平均方程可能有实质差别。

\section{Reynolds 平均不可压 Navier--Stokes 方程}

对常密度不可压方程作 Reynolds 平均：
\begin{align}
 \nabla\cdot(\rho\overline{\bm v})&=0, \tag{5.2.1}\\
 \partial_t(\rho\overline{\bm v})+
 \nabla\cdot(\rho\overline{\bm v}\otimes\overline{\bm v})+
 \nabla\overline p-\nabla\cdot\widehat{\bm\tau}&=0, \tag{5.2.2}\\
 \widehat{\bm\tau}&=\mu[\nabla\overline{\bm v}+(\nabla\overline{\bm v})^{\mathsf T}]
 -\rho\overline{\bm v'\otimes\bm v'}. \tag{5.2.3}
\end{align}
平均操作新产生的 $\rho\overline{\bm v'\otimes\bm v'}$ 被解释为应力，称为 Reynolds 应力；没有它，平均方程会像层流方程一样闭合。虽然可以继续为 Reynolds 应力推导输运方程，更简便的做法是采用 Boussinesq 近似（梯度输运假设）：
\begin{equation}
 -\rho\overline{\bm v'\otimes\bm v'}
 =\mu_T[\nabla\overline{\bm v}+(\nabla\overline{\bm v})^{\mathsf T}]
 -\frac23\rho k\bm I,\qquad
 k=\frac12\overline{\bm v'\cdot\bm v'}. \tag{5.2.5--5.2.6}
\end{equation}
$\mu_T$ 是涡黏度（湍流黏度），$k$ 是湍动能。第二项保证取迹后两端一致，因为左端迹为 $-2\rho k$。零方程或一方程模型无法求 $k$ 时通常把该项略去；若要保留，则需为 $k$ 再求一条方程，典型二方程模型会这样做。

忽略 $k$ 后，总有效应力相当于把分子黏度替换为 $\mu+\mu_T$，所以 RANS 方程与层流方程形式相同，问题集中为如何估计 $\mu_T$：可代数计算、求一条 $\mu_T$ 输运方程，或联立两条湍流量输运方程。

\section{Favre 平均可压缩 Navier--Stokes 方程}

对可压缩方程采用 Favre 平均，可得到相对简洁的守恒形式。未作闭合近似前，平均总能和总焓含湍动能 $k$，平均应力含 $-\overline{\rho\bm v''\otimes\bm v''}$，热流含 $-\overline{\rho\bm v''h''}$，能量方程还含较小的三阶相关项 $\bm K$。完整定义如下。
\sourcecrop{187}{38}{88}{574}{708}{0.42}
严格说，因 $\overline{\bm v''}\ne0$，分子应力还含脉动速度梯度；这些项通常视为很小而略去。$\bm K$ 以及平均总能、总焓中的 $k$ 也经常略去。随后采用 Boussinesq 近似
\begin{align}
 -\overline{\rho\bm v''\otimes\bm v''}&=
 \mu_T[\nabla\widetilde{\bm v}+(\nabla\widetilde{\bm v})^{\mathsf T}]
 -\frac23\mu_T(\nabla\cdot\widetilde{\bm v})\bm I-\frac23\overline\rho k\bm I,\\
 -\overline{\rho\bm v''h''}&=-\frac{\mu_T}{\Pr_T}\nabla\widetilde h. \tag{5.3.12--5.3.13}
\end{align}
$\Pr_T$ 为湍流 Prandtl 数，海平面空气可近似取 $0.9$。若无从估计 $k$，也将其略去。最终 Favre 平均方程与层流方程的差别仅为
\begin{equation}
 \mu\longrightarrow\mu+\mu_T,\qquad
 \frac{\mu}{\Pr}\longrightarrow\frac{\mu}{\Pr}+\frac{\mu_T}{\Pr_T}. \tag{5.3.22--5.3.23}
\end{equation}
因此与不可压 RANS 一样，核心仍是估计 $\mu_T$。

\section{滤波不可压 Navier--Stokes 方程}

对不可压方程作空间滤波，质量、动量方程保持守恒形式，但出现 SGS 应力
\begin{equation}
 \widehat{\bm\tau}^{\rm SGS}=-\rho(\overline{\bm v\otimes\bm v}-
 \overline{\bm v}\otimes\overline{\bm v}). \tag{5.4.3}
\end{equation}
把 $\bm v=\overline{\bm v}+\bm v'$ 展开后，右端依次分成 Leonard 应力、交叉应力和 SGS Reynolds 应力：
\begin{equation}
 \widehat{\bm\tau}^{\rm SGS}=
 (\overline{\overline{\bm v}\otimes\overline{\bm v}}-
 \overline{\bm v}\otimes\overline{\bm v})+
 \overline{\bm v'\otimes\overline{\bm v}+\overline{\bm v}\otimes\bm v'}+
 \overline{\bm v'\otimes\bm v'}. \tag{5.4.4}
\end{equation}
它比不可压 RANS 中单一 Reynolds 应力稍复杂。

\section{滤波可压缩 Navier--Stokes 方程}

对可压缩方程使用质量加权滤波，可得到与 Favre 平均形式相似的连续、动量和总能方程。新增的 SGS 量包括亚格子动能 $k^{\rm SGS}$、SGS 应力 $\bm\tau^{\rm SGS}$、SGS 热/能量通量 $\bm q^{\rm SGS}$ 以及 $\bm K_{\rm SGS}$。其完整定义为
\sourcecrop{189}{38}{88}{574}{708}{0.41}
Mach 数不大时通常忽略 $k^{\rm SGS}$ 和 $\bm K_{\rm SGS}$，只对 $\bm\tau^{\rm SGS}$ 与 $\bm q^{\rm SGS}$ 建模。尽管方程表面上比 Favre 平均方程复杂，本质仍是在黏性应力和热流中加入闭合项。

\section{湍流模型}

多数方法采用 Boussinesq 型近似，于是问题转化为估计涡黏度 $\mu_T$。最简单的是代数模型，直接由平均流量计算 $\mu_T$，包括 Prandtl 混合长度理论、Cebeci--Smith 模型、Baldwin--Lomax 模型及其众多变体。它们易于实现，在某些流动中能给出良好估计；但涡黏度不同于分子黏度，并非流体的内禀物性，模型成败依赖具体流动，常需针对应用调节常数。Johnson--King 模型引入一条常微分方程，常称半方程模型，但因不直接求 $\mu_T$，适用性仍有限。

比代数模型更进一步，可直接求解 $\mu_T$ 的输运方程，使它随流场共同演化。这类一方程模型包括 Baldwin--Barth、Spalart--Allmaras 模型及其变体。湍流方程可与流动方程耦合求解，也可分离求解。它们对许多流动很成功，但处理分离流仍可能不足。

再为另一个湍流量增加输运方程，就得到二方程模型，例如 Menter 剪切应力输运（SST）模型及其变体。除此以外还有 Reynolds 应力模型、LES 的 Smagorinsky 模型等。湍流建模范围很广，本书建议读者参考文献 [171] 获取更系统的讨论。

\chapter{精确解方法}

精确解广泛用于 CFD 精度验证，但许多资料只给结果而不说明推导过程。掌握多种构造方法后，遇到新的验证需要时便可自行生成合适的精确解。

\section{用变量分离法构造精确解}

变量分离能把某些偏微分方程化成若干常微分方程，是构造精确解的有力工具。

\subsection{Laplace 方程}

考虑
\begin{equation}
 u_{xx}+u_{yy}=0. \tag{6.1.1}
\end{equation}
设 $u=X(x)Y(y)$，代入后得到
\begin{equation}
 \frac1X\frac{d^2X}{dx^2}=-\frac1Y\frac{d^2Y}{dy^2}=\text{常数}. \tag{6.1.4}
\end{equation}
若要求 $x$ 方向周期性，把分离常数取为 $-k^2$，则
\begin{equation}
 X=A\sin kx+B\cos kx,\qquad Y=Ce^{ky}+De^{-ky}, \tag{6.1.5--6.1.6}
\end{equation}
$A,B,C,D$ 由边界条件确定。不同 $k$ 对应解的任意线性组合（Fourier 级数）仍是解，这就是叠加原理；拟合给定区域和边界条件时很有用，但验证单一数值格式通常一个模态已经足够。

例如单位正方形中，若 $u(x,0)=u(0,y)=u(1,y)=0$、$u(x,1)=\sin(\pi x)$，取 $C=-D$、$B=0$、$AC=1/\sinh\pi$，得到
\begin{equation}
 u(x,y)=\frac{\sin(\pi x)\sinh(\pi y)}{\sinh\pi}. \tag{6.1.12}
\end{equation}
交换 $x,y$ 可得到右边界给定 $\sin(\pi y)$ 的第二个解。将两解叠加，得到
\begin{equation}
 u(x,y)=\frac{\sinh(\pi x)\sin(\pi y)+\sinh(\pi y)\sin(\pi x)}{\sinh\pi}, \tag{6.1.16}
\end{equation}
它同时满足上、右边界的正弦条件。能这样直接叠加，根本原因是 Laplace 方程线性。

\subsection{扩散方程}

二维扩散方程
\begin{equation}
 u_t=\nu(u_{xx}+u_{yy}) \tag{6.1.20}
\end{equation}
中令 $u=T(t)X(x)Y(y)$。分别把空间分离常数取为 $-(a\pi)^2$、$-(b\pi)^2$，可得
\begin{equation}
 T=Ce^{-\nu\pi^2(a^2+b^2)t},\quad
 X=A_1\sin(a\pi x)+A_2\cos(a\pi x),\quad
 Y=B_1\sin(b\pi y)+B_2\cos(b\pi y). \tag{6.1.24--6.1.26}
\end{equation}
各常数由边界条件确定，不同 $a,b$ 模态的线性组合仍为精确解。用于验证时可取最简单的 $A_1=B_1=1$、$A_2=B_2=0$：
\begin{equation}
 u=Ce^{-\nu\pi^2(a^2+b^2)t}\sin(a\pi x)\sin(b\pi y). \tag{6.1.28}
\end{equation}
它在正方形齐次 Dirichlet 边界上足以验证扩散格式精度。

\subsection{平流--扩散方程}

考虑常系数二维方程
\begin{equation}
 u_t+au_x+bu_y=\nu(u_{xx}+u_{yy}),\qquad a^2+b^2\ne0. \tag{6.1.29}
\end{equation}
以 $\xi=ax+by$、$\eta=bx-ay$ 把坐标旋转到沿平流方向和法向，方程化为
\begin{equation}
 \frac{u_t}{a^2+b^2}+u_\xi=\nu(u_{\xi\xi}+u_{\eta\eta}). \tag{6.1.30}
\end{equation}
设 $u=T(t)X(\xi)Y(\eta)$ 并分离变量，得到时间指数、$\xi$ 方向两个指数根以及 $\eta$ 方向正弦/余弦。完整通解和特征根如下。
\sourcecrop{193}{38}{88}{574}{708}{0.42}
取 $A=B$ 可直接得到稳态解。通式虽完整但不够直观，常选 $C_1=C_3=0$、$C_2=C_4=1$ 得到较简单的指数--余弦解
\begin{equation}
 u=C_0\cos(B\pi\eta)\exp\left[
 \frac{1-\sqrt{1+4A^2\pi^2\nu^2}}{2\nu}\,\xi\right]. \tag{6.1.39}
\end{equation}

\subsection{平流--扩散--反应方程}

一维方程
\begin{equation}
 u_t+au_x=\nu u_{xx}-ku,\qquad a,\nu,k>0 \tag{6.1.40}
\end{equation}
中令 $u=T(t)X(x)$，并把时间分离常数取为 $-A^2$，则
\begin{equation}
 T=C_0e^{-A^2t},\qquad X=C_1e^{\lambda_1x}+C_2e^{\lambda_2x},\qquad
 \lambda_{1,2}=\frac{a\pm\sqrt{a^2+4\nu(k+A^2)}}{2\nu}. \tag{6.1.44--6.1.46}
\end{equation}
这立即给出精确解。取 $A^2=k$ 时，解按反应速率 $k$ 衰减；取 $A=0$ 则得到稳态指数解。原书给出的两种特化形式如下。
\sourcecrop{194}{38}{88}{574}{520}{0.28}

\subsection{复变量振荡解：扩散方程}

前述解随时间指数衰减，最终趋于零。要构造永不消失的周期解，可在变量分离中使用复数。考虑 $0<y<d$ 上
\begin{equation}
 u_t=\nu u_{yy},\qquad u(0,t)=0,\quad u(d,t)=U\cos\omega t. \tag{6.1.50--6.1.52}
\end{equation}
先把右边界写成 $Ue^{i\omega t}$，求复值问题，再取实部。设 $u=Y(y)T(t)$，边界条件提示分离常数为 $i\omega$，于是
\begin{equation}
 T=Ce^{i\omega t},\qquad Y=Ae^{\lambda y}+Be^{-\lambda y},\qquad
 \lambda=(1+i)\sqrt{\frac{\omega}{2\nu}}. \tag{6.1.55--6.1.57}
\end{equation}
施加两端边界得到复解
\begin{equation}
 u(y,t)=\frac{\sinh(\lambda y)}{\sinh(\lambda d)}Ue^{i\omega t}. \tag{6.1.59}
\end{equation}
取其实部即可得到原问题的实解；显式双曲函数--三角函数形式见原式。
\sourcecrop{195}{40}{90}{572}{580}{0.30}
该解不会随时间消失，适合研究扩散格式的长时间行为。

\subsection{复变量振荡解：平流--扩散方程}

同一程序可用于
\begin{equation}
 u_t+au_y=\nu u_{yy},\qquad u(0,t)=0,\quad u(d,t)=U\cos\omega t. \tag{6.1.62--6.1.64}
\end{equation}
仍取 $u=YT$、分离常数 $i\omega$，有
\begin{equation}
 T=Ce^{i\omega t},\quad Y=Ae^{\lambda_1y}+Be^{\lambda_2y},\quad
 \lambda_{1,2}=\frac{a\pm\sqrt{a^2+4\nu\omega i}}{2\nu}. \tag{6.1.66--6.1.68}
\end{equation}
施加边界条件得到
\begin{equation}
 u(y,t)=\frac{e^{\lambda_1y}-e^{\lambda_2y}}
 {e^{\lambda_1d}-e^{\lambda_2d}}Ue^{i\omega t}. \tag{6.1.69}
\end{equation}
其实部是原实值问题的精确解。即使不把实部手工展开，也可很容易数值计算，用于验证时间精确平流--扩散格式。

\section{用复变量构造精确解}

速度势 $\phi$ 与流函数 $\psi$ 满足 Cauchy--Riemann 系统
\begin{equation}
 \psi_x+\phi_y=0,\qquad \phi_x-\psi_y=0. \tag{6.2.1--6.2.2}
\end{equation}
因此它们分别是解析函数 $F(Z)=\phi+i\psi$ 的实部和虚部，其中 $Z=x+iy$。$F$ 称为复势，其导数
\begin{equation}
 W(Z)=F'(Z)=u-iv \tag{6.2.4}
\end{equation}
称为复速度。任意关于 $Z$ 的解析函数都代表一个稳态、不可压、无黏、无旋二维流动，因此任选解析函数即可生成精确解。基本例子如下。

角度为 $\alpha$ 的楔角绕流可取 $F=CZ^{\pi/\alpha}$；$\alpha=\pi$ 时是速度为常数 $C$ 的均匀流，$\alpha=\pi/2$ 时是冲向水平壁面的驻点流。与 $x$ 轴夹角为 $\alpha$ 的均匀流为 $F=U_\infty Ze^{-i\alpha}$，即 $u=U_\infty\cos\alpha$、$v=U_\infty\sin\alpha$。

位于 $Z_0=x_0+iy_0$、强度为 $\sigma$ 的源/汇为
\begin{equation}
 F=\frac{\sigma}{2\pi}\ln(Z-Z_0),\qquad
 W=\frac{\sigma}{2\pi(Z-Z_0)}. \tag{6.2.11--6.2.12}
\end{equation}
强度为 $m$、方向角为 $\alpha$ 的偶极子取 $F=-me^{i\alpha}/[2\pi(Z-Z_0)]$；位于 $Z_0$、正环量 $\Gamma$ 的顺时针点涡取 $F=i\Gamma\ln(Z-Z_0)/(2\pi)$。各自的实速度分量如下。
\sourcecrop{197}{38}{88}{574}{708}{0.42}

由于控制方程线性，这些基本流可叠加。例如原点单位圆柱的均匀绕流由均匀流与逆向偶极子相加得到
\begin{equation}
 F=U_\infty\left(Z+\frac1Z\right),\qquad
 W=U_\infty\left(1-\frac1{Z^2}\right). \tag{6.2.23--6.2.24}
\end{equation}
取偶极子强度 $m=2\pi U_\infty$ 后，圆柱表面本身成为流线。更多例子见第 7.11 节。

\section{非线性方程中的叠加}

严格叠加原理只适用于线性方程，但对 Euler 等非线性系统，有时仍可把一种特殊扰动“叠加”到均匀流。设均匀速度为 $\bm v_\infty$，扰动以该速度不变形地平移：
\begin{equation}
 \bm v(\bm x,t)=\bm v_\infty+\bm v'(\bm x-\bm x_0-\bm v_\infty t). \tag{6.3.2}
\end{equation}
密度、压力和温度也以同一速度平移。考虑等熵流，有
\begin{equation}
 \frac{p}{p_\infty}=\left(\frac{\rho}{\rho_\infty}\right)^\gamma,\quad
 \frac{\rho}{\rho_\infty}=\left(\frac{T}{T_\infty}\right)^{1/(\gamma-1)},\quad
 \frac{p}{p_\infty}=\left(\frac{T}{T_\infty}\right)^{\gamma/(\gamma-1)}. \tag{6.3.6}
\end{equation}
把平移形式代入 Euler 方程后，均匀平移项抵消；所需条件正是扰动满足稳态等熵 Euler 方程
\begin{equation}
 \nabla\cdot(\rho\bm v')=0,\qquad
 (\nabla\bm v')\bm v'+\frac{\gamma R}{\gamma-1}\nabla T'=0. \tag{6.3.9--6.3.10}
\end{equation}

二维涡扰动是一个重要例子。令极坐标扰动 $u_r'=0$、$u_\theta'=f(r)$，则连续方程和周向动量方程自动满足，径向平衡给出
\begin{equation}
 -T'=\frac{\gamma-1}{\gamma R}\int_r^\infty\frac{f(s)^2}{s}\,ds. \tag{6.3.14}
\end{equation}
选取
\begin{equation}
 f(r)=\frac{K}{2\pi}r e^{\alpha(1-r^2)/2},\qquad
 T'=-\frac{K^2(\gamma-1)}{8\alpha\pi^2\gamma R}e^{\alpha(1-r^2)}, \tag{6.3.15--6.3.16}
\end{equation}
即可完整积分。于是得到初始中心位于 $(x_0,y_0)$、随均匀流平移的等熵涡精确解；速度、温度、密度、压力和移动坐标的完整表达如下。
\sourcecrop{199}{38}{88}{574}{708}{0.42}
这是 Euler 方程的真正精确解：涡不会反作用于均匀背景流，只被平移输运。类似地，也可选用其他稳态扰动构造新的非线性“叠加”解。

\section{用变换构造精确解}

\subsection{Cole--Hopf 变换}

对一维黏性 Burgers 方程
\begin{equation}
 u_t+uu_x=\nu u_{xx}, \tag{6.4.1}
\end{equation}
Cole--Hopf 变换
\begin{equation}
 u=-2\nu\frac{\partial\ln\phi}{\partial x}=-2\nu\frac{\phi_x}{\phi} \tag{6.4.2}
\end{equation}
把任一正的线性扩散方程解 $\phi_t=\nu\phi_{xx}$ 映射为 Burgers 方程解。因此只需把已知扩散精确解代入该变换，就可产生许多黏性 Burgers 精确解。

二维 Burgers 系统中同样有
\begin{equation}
 u=-2\nu\phi_x/\phi,\qquad v=-2\nu\phi_y/\phi,\qquad
 \phi_t=\nu(\phi_{xx}+\phi_{yy}). \tag{6.4.4--6.4.7}
\end{equation}
若 $\phi>0$ 满足二维扩散方程，则 $(u,v)$ 满足二维 Burgers 方程。三维矢量形式进一步统一为
\begin{equation}
 \bm U=-2\nu\nabla\ln\phi,\qquad
 \phi_t=\nu\nabla^2\phi, \tag{6.4.9--6.4.10}
\end{equation}
从而把任一正的三维扩散解映射为矢量 Burgers 解。

还可处理带源项的一维方程
\begin{equation}
 u_t+uu_x=\nu u_{xx}+\epsilon u^2-ku+d. \tag{6.4.11}
\end{equation}
令 $u=p\phi_x/\phi$，要求变换后各个关于 $\phi$ 的独立组合分别为零，可确定 $p=-2\nu$，并把问题化为线性演化方程和常系数二阶空间方程。代入试探解 $\phi=C(t)e^{\lambda x}$ 后，特征方程的根为
\begin{equation}
 \lambda=\frac{-k\pm\sqrt{k^2-4\epsilon d}}{4\epsilon\nu}. \tag{6.4.21}
\end{equation}
判别式分别为正、负和零时，对应两个不同实根、共轭复根及重根。三种情形下 $\phi$ 的通解和由 $u=-2\nu\phi_x/\phi$ 得到的精确解如下。
\sourcecrop{201}{38}{88}{574}{708}{0.42}
\sourcecrop{202}{38}{88}{574}{708}{0.42}
\sourcecrop{203}{38}{250}{574}{708}{0.30}
该方法及更一般方程的精确解见文献 [9]。

\subsection{Hodograph 变换}

二维可压缩连续方程可由流函数自动满足：
\begin{equation}
 u=\frac{\rho_0}{\rho}\psi_y,\qquad
 v=-\frac{\rho_0}{\rho}\psi_x. \tag{6.4.40--6.4.41}
\end{equation}
由无旋条件和动量方程导出的 $\psi(x,y)$ 方程是高度非线性的，甚至比非线性势方程更复杂。Hodograph 变换交换自变量 $(x,y)$ 与因变量 $(V,\theta)$，其中 $V$ 是速度大小、$\theta$ 是流向角，可把它化为关于 $\psi(V,\theta)$ 的线性方程
\begin{equation}
 V^2\psi_{VV}+V\left(1+\frac{V^2}{c^2}\right)\psi_V+
 \left(1-\frac{V^2}{c^2}\right)\psi_{\theta\theta}=0. \tag{6.4.45}
\end{equation}
线性化后既可系统求解，也可直接试代候选函数；任意解的线性组合仍为解。第 7.13.6 节的 Ringleb 流就是著名例子。

\section{线性守恒律系统的精确解}

考虑常系数线性双曲系统
\begin{equation}
 \bm U_t+\bm A\bm U_x=0, \tag{6.5.1}
\end{equation}
$\bm A$ 具有实特征值和完备特征矢量。设 $\bm L$、$\bm R$ 分别为左、右特征矢量矩阵，$\bm\Lambda=\bm L\bm A\bm R$，定义特征变量 $\bm W=\bm L\bm U$，系统即对角化为
\begin{equation}
 \bm W_t+\bm\Lambda\bm W_x=0. \tag{6.5.6}
\end{equation}
每一分量都是标量平流方程，故 $w_k(x,t)=w_k(x-\lambda_kt,0)$。变换回原变量得到
\begin{equation}
 \bm U(x,t)=\sum_k[\bm\ell_k\bm U_0(x-\lambda_kt)]\bm r_k. \tag{6.5.10}
\end{equation}

简单波是变化仅限于某一特征矢量张成的一维子空间的解。若初始解是第 $j$ 个简单波，其余特征分量为常数，则
\begin{equation}
 \bm U(x,t)=\bm U_0(x-\lambda_jt). \tag{6.5.12}
\end{equation}
只有一个传播速度参与，所以波形被完整保持。例如
\begin{equation}
 \bm U_0=\overline{\bm U}+\sigma\sin x\,\bm r_j
 \quad\Longrightarrow\quad
 \bm U=\overline{\bm U}+\sigma\sin(x-\lambda_jt)\bm r_j. \tag{6.5.13--6.5.14}
\end{equation}
这一结果适用于任意该类线性双曲系统，也适用于把 $(x,t)$ 换成两个空间坐标的稳态系统。

\section{非线性守恒律系统的简单波}

考虑
\begin{equation}
 \bm U_t+\bm F(\bm U)_x=0,\qquad
 \bm U_t+\bm A(\bm U)\bm U_x=0. \tag{6.6.1--6.6.2}
\end{equation}
此时特征矩阵依赖解，不能再简单定义 $\bm W=\bm L\bm U$，只能微分地定义 $d\bm W=\bm L\,d\bm U$。对角化后每个特征分量满足非线性平流方程，其解在激波形成前可隐式写成 $w_k(x,t)=w_k(x-\lambda_kt,0)$；但 $\lambda_k$ 依赖解，一般无法给出闭式解，也不能像线性情形那样精确投影任意初值。

若只让第 $j$ 个特征变量变化，取任意光滑函数 $h$ 满足
\begin{equation}
 h_t+\lambda_j(\bm U)h_x=0, \tag{6.6.11}
\end{equation}
其余特征变量保持常数，则沿 $\xi=x-\lambda_jt$ 有
\begin{equation}
 \frac{d\bm U}{d\xi}=\bm r_j(\bm U)\frac{dh}{d\xi}. \tag{6.6.13}
\end{equation}
只要这一常微分方程可积，就能得到第 $j$ 个简单波的精确解。变量 $\bm U$ 可取守恒变量或原始变量，只要特征矢量定义与之相容。

\section{若干简单波精确解}

\subsection{第 j 个特征矢量为常矢量}

若特征矢量与解无关，直接积分得到
\begin{equation}
 \bm U(\xi)=h(\xi)\bm r_j+\text{常量}. \tag{6.7.1}
\end{equation}
Euler 方程的熵波就是例子。对原始变量 $(\rho,u,p)$，对应特征矢量为 $(1,0,0)^{\mathsf T}$、特征值为 $u$，所以密度可以是任意 $h(x-ut)$ 加常数，而速度和压力保持常数。由于传播速度不变，熵波线性退化，任意连续甚至不连续波形都只会平移而不变形。

\subsection{第 j 个特征矢量与状态矢量成线性关系}

若 $\bm r_j=\bm M\bm U$、$\bm M$ 为常矩阵，则
\begin{equation}
 \frac{d\bm U}{d\xi}=\bm M\bm U\,h'(\xi) \tag{6.7.7}
\end{equation}
是线性常微分方程系统，可用标准对角化解析求解。理想磁流体力学的 Alfv\'en 波是重要例子。其密度、流向速度和压力跨波保持常数，故特征值 $v_x\pm B_x/\sqrt{4\pi\rho}$ 也是常数，波线性退化。原书列出完整 MHD Jacobian、Alfv\'en 特征矢量以及四个相关横向量的积分过程。
\sourcecrop{206}{38}{88}{574}{708}{0.42}
\sourcecrop{207}{38}{88}{574}{708}{0.42}
一般解表明 $B_y,B_z$ 在相平面作旋转，而横向速度由 Alfv\'en 关系与磁场相联系。取适当的 $h$ 和常数，可得到正向 Alfv\'en 正弦简单波：
\sourcecrop{208}{38}{360}{574}{708}{0.23}
由于它线性退化，波形应在任意时间保持，可作为 MHD 程序验证的良好算例。

\subsection{其他情形：Euler 声波}

Euler 方程的声学特征矢量（原始变量）为
\begin{equation}
 \bm r_j^\pm=(\pm\rho/a,\,1,\,\pm\rho a)^{\mathsf T},\qquad
 \lambda_j^\pm=u\pm a. \tag{6.7.26--6.7.27}
\end{equation}
正、负声波必须二选一才能构成简单波。速度分量可直接积分为 $u=h+u_\infty$。对绝热关系 $p/\rho^\gamma=p_\infty/\rho_\infty^\gamma$，用 $a/a_\infty=(\rho/\rho_\infty)^{(\gamma-1)/2}$ 积分密度方程，得到
\begin{align}
 \frac{u}{a_\infty}&=M_\infty+\frac{h}{a_\infty},\\
 \frac{\rho}{\rho_\infty}&=\left[1\pm\frac{\gamma-1}{2}
 \left(\frac{u}{a_\infty}-M_\infty\right)\right]^{2/(\gamma-1)},\\
 \frac{p}{p_\infty}&=\left[1\pm\frac{\gamma-1}{2}
 \left(\frac{u}{a_\infty}-M_\infty\right)\right]^{2\gamma/(\gamma-1)}. \tag{6.7.38--6.7.40}
\end{align}
这是隐式解，因为 $\xi=x-(u\pm a)t$ 中的特征速度依赖解；它仅在激波形成之前有效，且 $h$ 必须保证密度、压力为正。

该声学简单波还满足无黏 Burgers 方程。以正向波为例，令 $V=u+a$，则 $V_t+VV_x=0$；又由 Riemann 不变量 $u-2a/(\gamma-1)=\text{常数}$，可把任一 Burgers 解 $V(\xi)$ 转换为 Euler 简单波。也可在平均流上添加任意满足正性约束的 Burgers 扰动。归根结底，若 $d\bm W=\bm L\,d\bm U$ 可积并能显式定义 Riemann 不变量，就可获得闭式简单波；浅水方程和等温 Euler 方程也具有这一性质。

\section{制造解}

制造解法的思想极其直接：任选一个函数代入待验证控制方程，把剩余不平衡量定义为源项，这样便为带该源项的方程“制造”出精确解。

例如一维线性平流方程中选 $u=\sin x$，代入 $au_x=f$ 得 $f=a\cos x$，故带此源项的方程以 $u=\sin x$ 为精确解。Poisson 方程中选 $u=\sin x+\cos y$，则 $u_{xx}+u_{yy}=f$ 要求 $f=-\sin x-\cos y$。

引入源项会增加原方程的复杂性；若源项离散不慎，数值格式甚至可能失去形式精度。也可以不加源项而选择合适系数。例如把 $u=x^2+y^2$ 代入 $au_x+bu_y=0$，取 $a=y$、$b=-x$，就得到
\begin{equation}
 yu_x-xu_y=0,\qquad u=x^2+y^2. \tag{6.8.8--6.8.9}
\end{equation}

制造解对 Navier--Stokes 方程、湍流模型等复杂系统尤其方便。还可把不可压缩流精确解代入可压缩方程并生成相应源项，从而制造外形更像真实流动的解。物理合理性非常重要；例如若制造解导致负黏度，它对实际验证就没有意义。

\section{细网格数值解}

有人把极细网格上的数值解当作“精确解”，据此计算粗网格误差和估计收敛阶。这种做法在严格意义上是错误的，因为细网格数值解仍不是精确解。这样的误差研究只能说明解趋向某个特定细网格解的速度，而不能说明趋向真实精确解的速度。若格式本身不一致，所谓误差和精度阶甚至完全失去意义：它只能告诉你数值解距离一个错误解有多近、或者多快收敛到错误解。必须始终牢记：细网格解不是精确解。


% ---- 合并分片 4 ----
\chapter{精确解}

精确解既能揭示方程的传播、扩散和非线性结构，也是验证 CFD 程序最可靠的基准。
本章汇集线性输运、Burgers 方程、势流、欧拉方程和 Navier--Stokes 方程的典型解析
解。除非另有说明，常数均由初始条件或边界条件确定。

\section{线性对流}
\subsection{一维线性对流}
考虑 $u_t+a u_x=0$，其中 $a$ 为常对流速度。沿特征线 $x-at=$ 常数有
\begin{equation}u(x,t)=u_0(x-at).
\end{equation}
因此任意初始剖面只作刚体平移而不改变形状。周期域内可用正弦波、方波、三角波或
高斯包检验数值格式的相位误差和耗散；间断初值还能检验非振荡性质。若 $a>0$，入口
在左边界；若 $a<0$，入口在右边界。

\subsection{二维直线与圆周对流}
二维常系数方程 $u_t+a u_x+b u_y=0$ 的解为
$u(x,y,t)=u_0(x-at,y-bt)$。沿任意方向平移的光滑隆起可同时检验网格方向偏置。
若速度场为 $a=-\omega y,b=\omega x$，特征线是绕原点旋转的圆，解等于把初值旋转
角 $\omega t$；一个周期后应精确回到初态。三维加入恒定轴向速度后，轨迹成为螺旋线，
可用于检验真正三维的输运与网格适应。

式中 $u_0$ 为初始分布，$(a,b)$ 为平移速度，$\omega$ 为角速度。原书给出的多个紧支撑
脉冲和高斯分布可直接作为代码回归测试。

\section{扩散方程}
\subsection{一维、二维与三维}
一维扩散方程 $u_t=\nu u_{xx}$ 的基本解是宽度随 $\sqrt{\nu t}$ 增长、峰值随时间
降低的高斯核：
\begin{equation}
 u(x,t)=\frac{1}{\sqrt{4\pi\nu t}}
 \exp\!\left[-\frac{(x-x_0)^2}{4\nu t}\right].
\end{equation}
与任意初值卷积即可得到解。正弦模态
$u=e^{-\nu k^2t}\sin(kx)$ 则以速率 $\nu k^2$ 衰减，高波数衰减更快。
二维和三维基本解分别是各坐标一维高斯核的乘积；在 $d$ 维中前因子为
$(4\pi\nu t)^{-d/2}$，指数含到源点的欧氏距离平方。这里 $\nu$ 是扩散系数。

\section{对流--扩散方程}
\subsection{一维}
对 $u_t+a u_x=\nu u_{xx}$，令移动坐标 $\xi=x-at$ 即化为纯扩散方程。因此高斯解
在以 $a$ 平移的同时展宽，傅里叶模态为
\begin{equation}u(x,t)=e^{-\nu k^2t}\sin[k(x-at)].\end{equation}
半无限域和有限区间的解可用误差函数或分离变量构造。稳态边值问题
$a u_x=\nu u_{xx}$ 的解为指数函数；高 P\'eclet 数时边界层集中在下游边界，是检验
迎风格式人工扩散的经典算例。

\subsection{二维}
常速度二维方程的基本解是中心位于 $(x_0+at,y_0+bt)$ 的椭圆/圆形高斯包；各方向
扩散率不同时，指数和归一化因子分别使用 $\nu_x,\nu_y$。原书还给出旋转速度场下
带扩散的精确解、稳态边界层解和用于非结构网格误差分析的光滑解。

\section{对流--反应与对流--扩散--反应}
三维常系数反应输运方程
$u_t+\bm a\cdot\nabla u=ku$ 的解是沿特征平移并乘以 $e^{kt}$：
\begin{equation}u(\bm x,t)=e^{kt}u_0(\bm x-\bm a t).
\end{equation}
加入扩散 $\nu\nabla^2u$ 后，再与移动高斯核卷积。$k<0$ 表示衰减，$k>0$ 表示增长。
这些解可分别隔离检验对流、扩散和刚性源项的时间离散误差。

\section{球面输运方程}
球坐标中纯径向守恒输运包含几何因子：
$u_t+r^{-2}\partial_r(r^2 a u)=0$。若 $a$ 为常数，则 $r^2u$ 沿 $r-at$ 平移，故
$u(r,t)=[(r-at)/r]^2u_0(r-at)$。几何衰减不是物理扩散，而是球面面积随半径增长的
结果；该解适合检验几何源项或曲线坐标守恒性。

\section{Burgers 方程}
\subsection{一维无黏方程}
无黏 Burgers 方程
\begin{equation}u_t+u u_x=0\end{equation}
沿特征线 $dx/dt=u$ 保持 $u$ 不变，隐式解为 $u(x,t)=u_0(x-u t)$。由于特征速度就是
解本身，较大 $u$ 的特征追上较小 $u$ 的特征时形成激波；反向排列则形成稀疏波。
光滑初值首次破裂时间为
\begin{equation}t_s=-1/\min_x u_0'(x)\end{equation}
（若最小导数为负）。激波速度由 Rankine--Hugoniot 条件给出
$s=(u_L+u_R)/2$，熵条件要求 $u_L>u_R$。

\subsection{Riemann 问题}
初值在 $x=0$ 左右分别为常数 $u_L,u_R$。若 $u_L>u_R$，解为以速度
$(u_L+u_R)/2$ 移动的激波；若 $u_L<u_R$，解为自相似稀疏扇：
\begin{equation}
 u(x,t)=\begin{cases}
 u_L,&x/t<u_L,\\ x/t,&u_L\le x/t\le u_R,\\u_R,&x/t>u_R.
 \end{cases}
\end{equation}
若 $u_L=u_R$ 则保持常值。二维 Burgers 型方程仍可用特征构造，书中给出的径向/旋转
光滑解曾用于网格自适应测试。

\section{黏性 Burgers 方程}
\subsection{一维}
方程 $u_t+uu_x=\nu u_{xx}$ 可由 Cole--Hopf 变换
$u=-2\nu(\ln\phi)_x$ 化为热方程 $\phi_t=\nu\phi_{xx}$。这给出平滑行进激波
\begin{equation}
 u=u_R+\frac{u_L-u_R}{1+\exp[(u_L-u_R)(x-st)/(2\nu)]},
 \quad s=\frac{u_L+u_R}{2}.
\end{equation}
厚度与 $\nu/(u_L-u_R)$ 成正比，$\nu\to0$ 时趋于无黏间断。带源项时可反向代入指定
解析函数计算源项，即构造有严格解的受迫问题。二维耦合 Burgers 方程也可通过势函数
或特别选择的速度场得到光滑解。

\section{Laplace 方程与 Poisson 方程}
Laplace 方程 $\nabla^2u=0$ 的解称调和函数。二维可用分离变量得到
$e^{kx}\sin ky$、$\sinh(kx)\sin(ky)$ 等，也可取任意解析复函数的实部或虚部；三维
有 $1/r$、球谐函数以及可分离的三角--双曲函数。Poisson 方程
$\nabla^2u=f$ 可先指定光滑 $u$，再精确微分得到 $f$ 和边界值。原书列出的二维、三维
多项式、三角函数和指数函数解适合分别检验内部离散阶与边界实现。

\section{不可压缩欧拉/Cauchy--Riemann/Laplace 解}
二维不可压缩无旋流可用复势 $W(z)=\phi+i\psi$ 描述，速度满足
$u=\phi_x=\psi_y$、$v=\phi_y=-\psi_x$。因此速度势、流函数和不可压缩欧拉方程三种
表述彼此等价，压强由伯努利关系求得。

\subsection{驻点流与 Rankine 半体}
平面驻点流可取 $u=ax,v=-ay$，流函数 $\psi=axy$，压强随 $x^2+y^2$ 二次变化。
把均匀流与点源叠加得到 Rankine 半体；物体表面是穿过驻点的分离流线，其极坐标
方程由 $\psi=$ 常数确定。该解同时检验曲壁滑移条件与远场边界。

\subsection{方腔、圆柱与椭圆柱绕流}
方腔调和解可用 Fourier 级数满足各壁给定法向速度。半径 $a$ 的圆柱置于速度
$U_\infty$ 的均匀流中，复势为
\begin{equation}W(z)=U_\infty\left(z+\frac{a^2}{z}\right),\end{equation}
表面切向速度 $v_\theta=-2U_\infty\sin\theta$，压强系数
$C_p=1-4\sin^2\theta$。加入环量项
$-i\Gamma\ln z/(2\pi)$ 后产生升力，并由 Kutta--Joukowski 关系
$L'=-\rho U_\infty\Gamma$ 给出单位展向升力。保角映射把圆柱解变换为椭圆柱解。

\subsection{翼型绕流与 Fraenkel 常涡量解}
Joukowski 变换 $\zeta=z+c^2/z$ 把偏心圆映射为带厚度和弯度的翼型。将圆柱复势、
环量与映射导数组合即可得到翼型表面速度；在尖后缘要求速度有限的 Kutta 条件唯一
确定环量。原书进一步给出参数化翼型族、表面压强分布以及多种映射例子。Fraenkel
解则描述圆柱外具有常涡量区域的不可压缩流，展示势流以外仍可存在解析有旋解。

\section{三维不可压缩欧拉方程：轴对称流}
\subsection{驻点流、Rankine 半体与球绕流}
三维轴对称驻点流中，轴向速度与 $z$ 成正比，径向速度与 $r$ 成反比符号的线性函数，
以保证散度为零。把均匀流与三维点源叠加得到轴对称 Rankine 半体。半径 $a$ 的球在
均匀流中的速度势为均匀流势与偶极势之和；球面径向速度为零，切向速度为
$-(3/2)U_\infty\sin\theta$，压强系数为
$1-(9/4)\sin^2\theta$。这些解是三维曲面边界和压力计算的重要验证题。

\section{欧拉方程的精确解}
\subsection{一维简单声波与熵波}
对等熵简单波，一个黎曼不变量保持常数，另一个满足 Burgers 型标量方程。由指定的
速度或声速初值可隐式构造在破裂前的光滑声波。熵波则满足压力和速度恒定、密度任意
剖面以速度 $u_0$ 平移；它是检验接触波分辨率的直接算例。

\subsection{二维非定常等熵涡对流}
在均匀流上叠加局部等熵涡。以涡心距离 $r$ 定义高斯型切向速度扰动，再通过等熵关系
确定温度、密度和压强扰动；整个涡以自由来流速度平移而保持形状。这一光滑周期问题
广泛用于测量高阶格式的收敛阶，因为没有激波或边界误差干扰。

\subsection{线性化势流、喷管流与 Ringleb 流}
线性化势方程在亚声速区为椭圆型、超声速区为双曲型；用 Prandtl--Glauert 坐标缩放
可把亚声速问题化为 Laplace 方程，并由解析势函数计算压强系数。

准一维等熵喷管中，质量流率 $\rho uA$、总焓和总熵沿管恒定。面积--马赫数关系为
\begin{equation}
 \frac{A}{A^*}=\frac1M\left[\frac{2}{\gamma+1}
 \left(1+\frac{\gamma-1}{2}M^2\right)\right]^{(\gamma+1)/[2(\gamma-1)]}.
\end{equation}
喉部 $A=A^*$ 对应 $M=1$。给定背压后可得到全亚声速、喉部临界、超声速膨胀或
含正常激波的分段精确解；激波两侧由 Rankine--Hugoniot 关系连接。

Ringleb 流是通过交换流函数与速度变量构造的二维可压缩无旋精确解。物理坐标由速度
大小和流向隐式给出，边界由等流函数线及等速度线组成。它同时包含明显几何弯曲和
变系数，是验证高阶曲线网格欧拉求解器的经典基准。

\subsection{圆柱超声速绕流的滞止压与正常/斜激波}
对圆柱前方超声速流，激波后的总压由正常激波关系与随后等熵减速关系组合，可精确
给出驻点压力，用于检查壁面压力和激波损失。正常激波满足质量、法向动量和总焓守恒；
切向速度不变。上游马赫数给定后，可显式求下游密度比、压强比、温度比、马赫数和
熵增。斜激波只需把马赫数替换为法向分量，再保留切向速度；激波角、偏转角和马赫数
满足经典 $\theta$--$\beta$--$M$ 关系，并有弱、强两支解。

\section{欧拉方程的 Riemann 问题}
\subsection{基本关系与问题构造}
一维初值在界面左右分别为常状态
$\bm W_L=(\rho_L,u_L,p_L)$、$\bm W_R=(\rho_R,u_R,p_R)$。解由左声波、接触波和右
声波组成；接触波两侧速度与压力相等，密度可跳跃。每一声波可能是激波或稀疏波。
给定希望出现的波型，可从中间状态出发，沿等熵稀疏关系或 Rankine--Hugoniot 激波
关系反推左右状态，从而有控制地构造测试问题。

\subsection{精确 Riemann 求解器}
核心未知量是星区压力 $p^*$。左右波各给出速度变化函数 $f_L(p),f_R(p)$，求解标量
非线性方程
\begin{equation}f_L(p^*)+f_R(p^*)+u_R-u_L=0.
\end{equation}
用 Newton 法迭代得到 $p^*$ 后，
$u^*=(u_L+u_R+f_R-f_L)/2$。再根据 $p^*$ 与各侧压力的大小判定激波或稀疏波，计算
波速并在给定相似坐标 $\xi=x/t$ 采样。稀疏扇内使用黎曼不变量，激波后密度用跳跃
关系；接触面左右分别采用对应星区密度。原书完整列出算法、初值估计和采样分支。

\section{不可压缩 Navier--Stokes 方程}
本节解析解满足 $\nabla\cdot\bm v=0$ 和
$\bm v_t+(\bm v\cdot\nabla)\bm v=-\nabla p/\rho+\nu\nabla^2\bm v$。它们覆盖光滑
时间解、内部流、边界层和低雷诺数外流。

\subsection{光滑二维、三维非定常解}
通过三角函数速度场使连续方程恒等满足，并选择指数时间因子平衡黏性扩散，可得到
Taylor--Green 型二维涡解和三维周期解。压强由非线性项积分得到。这些光滑解特别适合
严格测量空间、时间离散收敛阶。

\subsection{Couette、Poiseuille 及圆管流}
两平板相对运动的稳态 Couette 流速度在壁间线性变化；非定常启动 Couette 流由稳态
线性剖面加衰减 Fourier 正弦级数组成。若同时施加压强梯度，得到线性 Couette 分量与
抛物线 Poiseuille 分量之和。圆管 Hagen--Poiseuille 流为
\begin{equation}u_z(r)=-\frac1{4\mu}\frac{dp}{dz}(R^2-r^2),
\end{equation}
体积流率与 $R^4$ 成正比。

同轴圆柱沿轴向运动时，速度是 $\ln r$ 的线性组合；圆柱旋转时，周向速度为
$Ar+B/r$。常数由内外圆柱壁速确定。这些解检验曲线坐标黏性项和无滑移边界。

\subsection{平板边界层、黏性驻点流与球体 Stokes 流}
零压梯度平板层流边界层采用相似变量
$\eta=y\sqrt{U_\infty/(\nu x)}$，Blasius 函数满足
\begin{equation}f'''+\frac12 f f''=0,\qquad f(0)=f'(0)=0,\quad f'(\infty)=1.
\end{equation}
这是非线性常微分方程的数值“精确”基准，可得到壁面摩擦系数和位移厚度。黏性驻点
流用 Hiemenz 相似变换得到另一边值问题，外流应变率决定边界层尺度。

极低 Reynolds 数球绕流忽略惯性项，满足 Stokes 方程。速度和压强有闭式表达，阻力
$D=6\pi\mu aU_\infty$，即 Stokes 阻力定律。该解检验黏性应力、压力积分和球面
边界。

\section{可压缩 Navier--Stokes 方程}
\subsection{可压缩 Couette 流}
两平板间稳态可压缩 Couette 流在定物性、零压强梯度下速度仍为线性分布，但黏性耗散
使温度成为二次函数。若壁温给定，温度由导热与 $\mu(du/dy)^2$ 平衡；密度再由理想
气体状态方程求得。该解同时检验动量、能量、状态方程和黏性功项。

\subsection{黏性激波结构}
在激波静止坐标中，一维稳态质量流率、动量通量和总能通量均为常数。利用这些积分
关系消去密度和温度，可得到速度的一阶常微分方程；对特殊 Prandtl 数存在闭式或易于
数值积分的精确剖面。黏度把无黏跳跃平滑为有限厚度层，远场极限仍严格满足正常激波
Rankine--Hugoniot 关系。原书给出离散积分算法、边界状态以及 $N=200$ 的参考曲线，
可用于检查热传导、黏性应力和激波厚度。

\subsection{制造解}
制造解方法不要求所选场本身是无源 Navier--Stokes 方程的自然解。先指定足够光滑且
满足所需边界行为的 $\rho,\bm v,p$（或 $T$），将其代入完整可压缩方程，解析计算每个
方程所需源项，再在代码中加入这些源项。若实现正确，数值解应按理论阶收敛到指定场。
该方法能独立覆盖对流、黏性、热传导、几何和边界项，是验证复杂求解器最系统的工具。
制造源项必须由与程序一致的变量定义、物性模型和无量纲化推导，不能遗漏能量方程中
的黏性功或热通量。

\end{document}
